% Generated mechanically from frozen input inputs/p09/ja/src/spaces-perfect.tex.
% Do not edit this generated file; edit the bound locale source instead.

% OK, start here.
%

\setcounter{chapter}{74}
\chapter{代数空間の導来圏}



\phantomsection
\label{spaces-perfect-section-phantom}


\section{はじめに}
\label{spaces-perfect-section-introduction}

\noindent
本章では、代数空間上の加群の導来圏を論じる。この主題を扱う適切な
入門的文献は見当たらないようであり、文献では代数スタック上の加群の
導来圏を扱う論文を参照する形で論じられている。例えば
\cite{olsson_sheaves} を参照されたい。



\section{規約}
\label{spaces-perfect-section-conventions}

\noindent
$\mathcal{A}$ をアーベル圏、$M$ を $\mathcal{A}$ の対象とする。
このとき、$M$ によって $K(\mathcal{A})$ および／または
$D(\mathcal{A})$ の対象で、$M$ を次数 $0$ にもち、それ以外の
次数では零である複体に対応するものも表す。

\medskip\noindent
環 $A$ に対し、$K(A)$ は $A$-加群の複体のホモトピー圏を表し、
$D(A)$ はそれに付随する導来圏を表す。同様に、環付き空間
$(X, \mathcal{O}_X)$ に対し、$K(\mathcal{O}_X)$ は
$\mathcal{O}_X$-加群の複体のホモトピー圏を表し、
$D(\mathcal{O}_X)$ はそれに付随する導来圏を表す。





\section{一般論}
\label{spaces-perfect-section-generalities}

\noindent
本節では、代数空間上の加群の非有界複体のコホモロジーに関する
いくつかの一般的結果を述べる。

\begin{lemma}
\label{spaces-perfect-lemma-restrict-direct-image-open}
$S$ を概型とし、$f : X \to Y$ を $S$ 上の代数空間の射とする。
平展射 $V \to Y$ が与えられたとき、$U = V \times_Y X$ とおき、
射影を $g : U \to V$ と表す。このとき
$(Rf_*E)|_V = Rg_*(E|_U)$ である。ただし $E$ は
$D(\mathcal{O}_X)$ の対象である。
\end{lemma}

\begin{proof}
$E$ を K-単射複体 $\mathcal{I}^\bullet$ で表す。これは
$\mathcal{O}_X$-加群の複体である。すると
$Rf_*(E) = f_*\mathcal{I}^\bullet$ であり、Cohomology on Sites,
Lemma \ref{sites-cohomology-lemma-restrict-K-injective-to-open} により
$Rg_*(E|_U) = g_*(\mathcal{I}^\bullet|_U)$ である。
したがって結論は Properties of Spaces,
Lemma \ref{spaces-properties-lemma-pushforward-etale-base-change-modules}
から従う。
\end{proof}











\begin{definition}
\label{spaces-perfect-definition-supported-on}
$S$ を概型、$X$ を $S$ 上の代数空間とし、
$E$ を $D(\mathcal{O}_X)$ の対象とする。また
$T \subset |X|$ を閉部分集合とする。$E$ が
{\it $T$ 上に台をもつ}とは、コホモロジー層 $H^i(E)$ が
$T$ 上に台をもつことをいう。
\end{definition}









\section{小平展サイト上の準連接加群の導来圏}
\label{spaces-perfect-section-derived-quasi-coherent-etale}

\noindent
$X$ を概型とする。本節では、$D_\QCoh(\mathcal{O}_X)$ を
小平展サイト $X_\etale$（$X$ の小平展サイト）によって
定義できることを示す。
$\mathcal{O}_\etale$ を $X_\etale$ 上の構造層とする。
環付きサイトの射
\begin{equation}
\label{spaces-perfect-equation-epsilon}
\epsilon :
(X_\etale, \mathcal{O}_\etale)
\longrightarrow
(X_{Zar}, \mathcal{O}_X).
\end{equation}
を考える。これは Descent, Lemma
\ref{descent-lemma-compare-sites} において
$\text{id}_{small, \etale, Zar}$ と表されている射である。

\begin{lemma}
\label{spaces-perfect-lemma-epsilon-flat}
(\ref{spaces-perfect-equation-epsilon}) の射 $\epsilon$ は環付きサイトの平坦射である。
特に、函手
$\epsilon^* : \textit{Mod}(\mathcal{O}_X) \to
\textit{Mod}(\mathcal{O}_\etale)$ は完全である。
さらに、$\epsilon^*\mathcal{F} = 0$ ならば
$\mathcal{F} = 0$ である。
\end{lemma}

\begin{proof}
$\epsilon$ の平坦性は Descent, Lemma
\ref{descent-lemma-compare-etale-zariski-flat} である。
別の証明も述べる。示すべきことは、
$\mathcal{O}_\etale$ が平坦な
$\epsilon^{-1}\mathcal{O}_X$-加群であることである。
そのためには、
$\mathcal{O}_{X, x} \to \mathcal{O}_{\etale, \overline{x}}$
が平坦であることを、任意の幾何学的点 $\overline{x}$（$X$ の点）
に対して確かめれば十分である。Modules on Sites,
Lemma \ref{sites-modules-lemma-check-flat-stalks}、Sites, Lemma
\ref{sites-lemma-point-morphism-sites}、および
\'Etale Cohomology, Remarks
\ref{etale-cohomology-remarks-enough-points} を参照されたい。
\'Etale Cohomology, Lemma
\ref{etale-cohomology-lemma-describe-etale-local-ring} により、
$\mathcal{O}_{\etale, \overline{x}}$ は
$\mathcal{O}_{X, x}$ の強ヘンゼル化である。したがって
$\mathcal{O}_{X, x} \to \mathcal{O}_{\etale, \overline{x}}$
は More on Algebra, Lemma
\ref{more-algebra-lemma-dumb-properties-henselization} により忠実平坦である。

\medskip\noindent
$\epsilon^*$ の完全性は、Modules on Sites, Lemma
\ref{sites-modules-lemma-flat-pullback-exact} により
$\epsilon$ の平坦性から従う。

\medskip\noindent
$\mathcal{F}$ を $\mathcal{O}_X$-加群とする。
$\epsilon^*\mathcal{F} = 0$ ならば、上の記法のもとで
$$
0 = \epsilon^*\mathcal{F}_{\overline{x}} =
\mathcal{F}_x \otimes_{\mathcal{O}_{X, x}} \mathcal{O}_{\etale, \overline{x}}
$$
がすべての幾何学的点 $\overline{x}$ に対して成り立つ
（Modules on Sites, Lemma
\ref{sites-modules-lemma-pullback-stalk}）。射
$\mathcal{O}_{X, x} \to \mathcal{O}_{\etale, \overline{x}}$
の忠実平坦性により、$\mathcal{F}_x = 0$ がすべての
$x \in X$ に対して成り立つと結論される。
\end{proof}

\noindent
$X$ を概型とし、記法を (\ref{spaces-perfect-equation-epsilon}) のとおりとする。
Descent, Proposition
\ref{descent-proposition-equivalence-quasi-coherent} および Remark
\ref{descent-remark-change-topologies-ringed-sites} により
$\epsilon^* : \QCoh(\mathcal{O}_X)
\to \QCoh(\mathcal{O}_\etale)$
が圏同値であることを思い出す。さらに、Descent, Lemma
\ref{descent-lemma-equivalence-quasi-coherent-limits} により
$\QCoh(\mathcal{O}_\etale)$ は
$\textit{Mod}(\mathcal{O}_\etale)$ の Serre 部分圏をなす。
そこで、Derived Categories, Section
\ref{derived-section-triangulated-sub} に従い、
$D_\QCoh(\mathcal{O}_\etale)$ を、
準連接なコホモロジー層をもつ複体を対象とする
$D(\mathcal{O}_\etale)$ の三角部分圏と定めることができる。
函手 $\epsilon^*$ は完全である
（Lemma \ref{spaces-perfect-lemma-epsilon-flat}）から
$\epsilon^* :  D(\mathcal{O}_X) \to D(\mathcal{O}_\etale)$
を誘導し、準連接加群の引き戻しは準連接であるから、さらに
$\epsilon^* : D_\QCoh(\mathcal{O}_X) \to
D_\QCoh(\mathcal{O}_\etale)$ を誘導する。

\begin{lemma}
\label{spaces-perfect-lemma-derived-quasi-coherent-small-etale-site}
$X$ を概型とする。上で定義した函手
$\epsilon^* : D_\QCoh(\mathcal{O}_X) \to
D_\QCoh(\mathcal{O}_\etale)$
は圏同値である。
\end{lemma}

\begin{proof}
函手
$R\epsilon_* : D(\mathcal{O}_\etale) \to D(\mathcal{O}_X)$
が準逆を誘導することを示す。ここでは、
$\epsilon_*$ が制限
$X_{Zar} \subset X_\etale$ によって与えられること、および
Descent, Lemma \ref{descent-lemma-compare-sites} における
$\epsilon^* = \text{id}_{small, \etale, Zar}^*$
の記述を自由に用いる。

\medskip\noindent
準連接 $\mathcal{O}_X$-加群 $\mathcal{F}$ に対する随伴射
$\mathcal{F} \to \epsilon_*\epsilon^*\mathcal{F}$
は同型である。実際、$\mathcal{F}^a$
（Descent, Definition
\ref{descent-definition-structure-sheaf}）が層であることは
Descent, Lemma
\ref{descent-lemma-sheaf-condition-holds} で示されている。
逆に、任意の準連接 $\mathcal{O}_\etale$-加群
$\mathcal{H}$ は、Descent, Proposition
\ref{descent-proposition-equivalence-quasi-coherent} により、
$\epsilon^*\mathcal{F}$ の形である。ここで
$\mathcal{O}_X$-加群 $\mathcal{F}$ は準連接である。このとき、今述べたことから
$\mathcal{F} = \epsilon_*\mathcal{H}$ であり、したがって
随伴射
$\epsilon^*\epsilon_*\mathcal{H} \to \mathcal{H}$
はすべての準連接 $\mathcal{O}_\etale$-加群
$\mathcal{H}$ に対して同型である。

\medskip\noindent
$E$ を $D_\QCoh(\mathcal{O}_\etale)$ の対象とし、
$\mathcal{H}^q = H^q(E)$ によってその第 $q$ コホモロジー層を表す。また、
$\mathcal{B}$ を $X_\etale$ のアフィン対象全体の集合とする。
すると
$H^p(U, \mathcal{H}^q) = 0$ はすべての $p > 0$、
すべての $q \in \mathbf{Z}$、および
すべての $U \in \mathcal{B}$ に対して成り立つ。
Descent, Proposition
\ref{descent-proposition-same-cohomology-quasi-coherent} および
Cohomology of Schemes, Lemma
\ref{coherent-lemma-quasi-coherent-affine-cohomology-zero} を参照されたい。
Cohomology on Sites, Lemma
\ref{sites-cohomology-lemma-cohomology-over-U-trivial} により、これは
$$
H^q(U, E) = H^0(U, \mathcal{H}^q)
$$
がすべての $U \in \mathcal{B}$ に対して成り立つことを意味する。
特に、アフィン開集合 $U \subset X$ に対しても成り立つ。
したがって、第 $q$ コホモロジーである $R\epsilon_*E$ の
$U$ 上での値は、層 $\epsilon_*\mathcal{H}^q$ の
$U$ における値である。
層化を施すと
$$
H^q(R\epsilon_*E) = \epsilon_*\mathcal{H}^q
$$
を得る。特に、$R\epsilon_*$ は函手
$D_\QCoh(\mathcal{O}_\etale) \to D_\QCoh(\mathcal{O}_X)$
を誘導する。$\epsilon^*$ は完全であるから
$H^q(\epsilon^*R\epsilon_*E) = \epsilon^*\epsilon_*\mathcal{H}^q =
\mathcal{H}^q$ となる（上の議論による）。
ゆえに随伴射
$\epsilon^*R\epsilon_*E \to E$ は同型である。

\medskip\noindent
逆に、$F \in D_\QCoh(\mathcal{O}_X)$ に対する随伴射
$F \to R\epsilon_*\epsilon^*F$
も同じ理由で同型である。実際、
$R\epsilon_*\epsilon^*F$ のコホモロジー層は
$\epsilon_*H^m(\epsilon^*F) = \epsilon_*\epsilon^*H^m(F) = H^m(F)$
に同型である。
\end{proof}










\section{準連接加群の導来圏}
\label{spaces-perfect-section-derived-quasi-coherent}

\noindent
$S$ を概型とする。Lemma
\ref{spaces-perfect-lemma-derived-quasi-coherent-small-etale-site}
は、圏 $D_\QCoh(\mathcal{O}_S)$ を、
$\mathcal{O}_S$-加群の複体（概型 $S$ 上のもの）によっても、
$\mathcal{O}$-加群の複体（$S$ の小平展サイト上のもの）によっても
定義できることを示している。したがって、次の定義は概型の場合の定義と整合する。

\begin{definition}
\label{spaces-perfect-definition-derived-quasi-coherent}
$S$ を概型、$X$ を $S$ 上の代数空間とする。
{\it 準連接なコホモロジー層をもつ $\mathcal{O}_X$-加群の導来圏}
を $D_\QCoh(\mathcal{O}_X)$ と表す。
\end{definition}

\noindent
この定義が意味をもつことは Properties of Spaces, Lemma
\ref{spaces-properties-lemma-properties-quasi-coherent}
および Derived Categories, Section
\ref{derived-section-triangulated-sub} から従う。
こうして標準的な函手
\begin{equation}
\label{spaces-perfect-equation-compare}
D(\QCoh(\mathcal{O}_X))
\longrightarrow
D_\QCoh(\mathcal{O}_X)
\end{equation}
を得る。Derived Categories, Equation
(\ref{derived-equation-compare}) を参照されたい。

\medskip\noindent
代数空間の平坦射 $f : Y \to X$ は完全函手
$f^* : \textit{Mod}(\mathcal{O}_X) \to \textit{Mod}(\mathcal{O}_Y)$
を誘導する。Morphisms of Spaces, Lemma
\ref{spaces-morphisms-lemma-flat-morphism-sites} および
Modules on Sites, Lemma
\ref{sites-modules-lemma-flat-pullback-exact} を参照されたい。
特に、
$Lf^* : D(\mathcal{O}_X) \to D(\mathcal{O}_Y)$
は任意の代表複体上で計算できる
（Derived Categories, Lemma
\ref{derived-lemma-right-derived-exact-functor}）。
$Lf^* = f^*$ と書くのは $f$ が平坦なときである。
この場合、$H^i(f^*E) = f^*H^i(E)$ が、$E$（
$D(\mathcal{O}_X)$ の対象）に対して成り立つ。
これは $f$ が平展な場合にしばしば用いる。もちろん平展の場合、
引き戻し函手は単に $Y_\etale$ への制限である。
Properties of Spaces, Equation
(\ref{spaces-properties-equation-restrict-modules}) を参照されたい。

\begin{lemma}
\label{spaces-perfect-lemma-check-quasi-coherence-on-covering}
$S$ を概型、$X$ を $S$ 上の代数空間とし、
$E$ を $D(\mathcal{O}_X)$ の対象とする。次の条件は同値である。
\begin{enumerate}
\item $E$ は $D_\QCoh(\mathcal{O}_X)$ に属する。
\item 任意の平展射 $\varphi : U \to X$ で $U$ が
アフィン概型であるものに対し、$\varphi^*E$ は
$D_\QCoh(\mathcal{O}_U)$ の対象である。
\item 任意の平展射 $\varphi : U \to X$ で $U$ が
概型であるものに対し、$\varphi^*E$ は
$D_\QCoh(\mathcal{O}_U)$ の対象である。
\item 全射平展射 $\varphi : U \to X$ で、$U$ が概型であり、
$\varphi^*E$ が
$D_\QCoh(\mathcal{O}_U)$ の対象となるものが存在する。
\item 代数空間の全射平展射 $f : Y \to X$ で、
$Lf^*E$ が $D_\QCoh(\mathcal{O}_Y)$ の対象となるものが存在する。
\end{enumerate}
\end{lemma}

\begin{proof}
これは補題の直前の議論と Properties of Spaces, Lemma
\ref{spaces-properties-lemma-characterize-quasi-coherent}
から直ちに従う。
\end{proof}

\begin{lemma}
\label{spaces-perfect-lemma-quasi-coherence-direct-sums}
$S$ を概型、$X$ を $S$ 上の代数空間とする。
このとき $D_\QCoh(\mathcal{O}_X)$ は直和をもつ。
\end{lemma}

\begin{proof}
Injectives, Lemma \ref{injectives-lemma-derived-products} により、
導来圏 $D(\mathcal{O}_X)$ は直和をもち、
任意の代表について項ごとの直和を取ることで計算される。
直和を取ることは完全函手である
（任意の Grothendieck アーベル圏において）から、
直和のコホモロジー層がコホモロジー層の直和であることは明らかである。
準連接層の直和は準連接であるため、Properties of Spaces, Lemma
\ref{spaces-properties-lemma-properties-quasi-coherent}
から補題が従う。
\end{proof}

\noindent
導来極限について若干の情報が必要となる。次の補題に現れる導来極限は、
通常 $D_\QCoh$ の対象ではないことに注意されたい。

\begin{lemma}
\label{spaces-perfect-lemma-Rlim-quasi-coherent}
$S$ を概型、$X$ を $S$ 上の代数空間とする。
$(K_n)$ を $D_\QCoh(\mathcal{O}_X)$ の逆系とし、その導来極限を
$K = R\lim K_n$ とする（これは $D(\mathcal{O}_X)$ における極限である）。
$H^q(K_{n + 1}) \to H^q(K_n)$ がすべての $q \in \mathbf{Z}$ と
$n \geq 1$ に対して
が全射であると仮定する。このとき
\begin{enumerate}
\item $H^q(K) = \lim H^q(K_n)$、
\item $R\lim H^q(K_n) = \lim H^q(K_n)$、および
\item 任意のアフィン開集合 $U \subset X$ に対し、
$H^p(U, \lim H^q(K_n)) = 0$（$p > 0$）
\end{enumerate}
が成り立つ。
\end{lemma}

\begin{proof}
$\mathcal{B} \subset \Ob(X_\etale)$ をアフィン対象全体の集合とする。
$H^q(K_n)$ は準連接であるから、
$H^p(U, H^q(K_n)) = 0$ が $U \in \mathcal{B}$ に対して成り立つ。
Cohomology of Spaces, Section
\ref{spaces-cohomology-section-higher-direct-image} における議論と
Cohomology of Schemes, Lemma
\ref{coherent-lemma-quasi-coherent-affine-cohomology-zero}
を参照されたい。同様の理由により、射
$H^0(U, H^q(K_{n + 1})) \to H^0(U, H^q(K_n))$ は
$U \in \mathcal{B}$ に対して
は全射である。(1) は、今確認した条件と Cohomology on Sites, Lemma
\ref{sites-cohomology-lemma-derived-limit-suitable-system}
から従う。(2) と (3) は Cohomology on Sites, Lemma
\ref{sites-cohomology-lemma-inverse-limit-is-derived-limit}
から従う。
\end{proof}

\begin{lemma}
\label{spaces-perfect-lemma-quasi-coherence-pullback}
$S$ を概型とし、$f : Y \to X$ を
$S$ 上の代数空間の射とする。函手 $Lf^*$ は
$D_\QCoh(\mathcal{O}_X)$ を $D_\QCoh(\mathcal{O}_Y)$
の中へ送る。
\end{lemma}

\begin{proof}
図式
$$
\xymatrix{
U \ar[d]_a \ar[r]_h & V \ar[d]^b \\
X \ar[r]^f & Y
}
$$
を選ぶ。ここで $U$ と $V$ は概型、縦の射は平展であり、
$a$ は全射である。% P09-ERR-0393
$a^* \circ Lf^* = Lh^* \circ b^*$ であるから、結論は
Lemma \ref{spaces-perfect-lemma-check-quasi-coherence-on-covering}
と概型の場合から従う。後者は Derived Categories of Schemes, Lemma
\ref{perfect-lemma-quasi-coherence-pullback} である。
\end{proof}

\begin{lemma}
\label{spaces-perfect-lemma-quasi-coherence-tensor-product}
$S$ を概型、$X$ を $S$ 上の代数空間とする。
$K, L$（$D_\QCoh(\mathcal{O}_X)$ の対象）に対し、
導来テンソル積 $K \otimes^\mathbf{L} L$ は
$D_\QCoh(\mathcal{O}_X)$ に属する。
\end{lemma}

\begin{proof}
$\varphi : U \to X$ を概型 $U$ からの全射平展射とする。
$$
\varphi^*(K \otimes_{\mathcal{O}_X}^\mathbf{L} L) =
\varphi^*K \otimes_{\mathcal{O}_U}^\mathbf{L} \varphi^*L
$$
であるから、Lemma
\ref{spaces-perfect-lemma-check-quasi-coherence-on-covering} により概型の場合へ帰着される。
その場合は Derived Categories of Schemes, Lemma
\ref{perfect-lemma-quasi-coherence-tensor-product} である。
\end{proof}

\noindent
次の補題は、$D_\QCoh(\mathcal{O}_X)$ の対象に対する
右導来函手の「計算」に役立つ。

\begin{lemma}
\label{spaces-perfect-lemma-nice-K-injective}
$S$ を概型、$X$ を $S$ 上の代数空間とし、
$E$ を $D_\QCoh(\mathcal{O}_X)$ の対象とする。このとき
Derived Categories, Remark
\ref{derived-remark-map-into-derived-limit-truncations}
の射 $E \to R\lim \tau_{\geq -n}E$ は同型である
\footnote{特に、$E$ は K-単射代表をもつ。Derived Categories, Lemma
\ref{derived-lemma-difficulty-K-injectives} を参照されたい。}。
\end{lemma}

\begin{proof}
$\mathcal{H}^i = H^i(E)$ を第 $i$ コホモロジー層
（$E$ のもの）と表す。$\mathcal{B}$ を
$X_\etale$ のアフィン対象全体の集合とする。
$H^p(U, \mathcal{H}^i) = 0$ はすべての $p > 0$、
すべての $i \in \mathbf{Z}$、すべての
$U \in \mathcal{B}$ に対して成り立つ。実際、$U$ はアフィン概型である。
Cohomology of Spaces, Section
\ref{spaces-cohomology-section-higher-direct-image} における議論と
Cohomology of Schemes, Lemma
\ref{coherent-lemma-quasi-coherent-affine-cohomology-zero}
を参照されたい。したがって補題は Cohomology on Sites, Lemma
\ref{sites-cohomology-lemma-is-limit-dimension} において
$d = 0$ とすることから従う。
\end{proof}

\begin{lemma}
\label{spaces-perfect-lemma-application-nice-K-injective}
$S$ を概型、$X$ を $S$ 上の代数空間とする。
$F : \textit{Mod}(\mathcal{O}_X) \to \textit{Ab}$
を函手、$N \geq 0$ を整数とする。次を仮定する。
\begin{enumerate}
\item $F$ は左完全である。
\item $F$ は可算直積と可換である。
\item $R^pF(\mathcal{F}) = 0$ がすべての $p \geq N$ と
準連接な $\mathcal{F}$ に対して成り立つ。
\end{enumerate}
このとき、$E \in D_\QCoh(\mathcal{O}_X)$ に対し
\begin{enumerate}
\item $H^i(RF(\tau_{\leq a}E) \to H^i(RF(E))$ は
$i \leq a$ のとき同型である。% P09-ERR-0394
\item $H^i(RF(E)) \to H^i(RF(\tau_{\geq b - N + 1}E))$ は
$i \geq b$ のとき同型である。
\item $H^i(E) = 0$ が $i \not \in [a, b]$ に対して成り立つとする。
ここである $-\infty \leq a \leq b \leq \infty$ を取る。
このとき $H^i(RF(E)) = 0$ が
$i \not \in [a, b + N - 1]$ に対して成り立つ。
\end{enumerate}
\end{lemma}

\begin{proof}
(1) は Derived Categories, Lemma
\ref{derived-lemma-negative-vanishing} である。

\medskip\noindent
(2) を証明する。$E_n = \tau_{\geq -n}E$ とおく。
Lemma \ref{spaces-perfect-lemma-nice-K-injective} により
$E = R\lim E_n$ である。したがって Injectives, Lemma
\ref{injectives-lemma-RF-commutes-with-Rlim} により、
$RF(E) = R\lim RF(E_n)$ が $D(\textit{Ab})$ において成り立つ。ゆえに各
$i \in \mathbf{Z}$ に対して短完全列
$$
0 \to R^1\lim H^{i - 1}(RF(E_n)) \to H^i(RF(E)) \to \lim H^i(RF(E_n)) \to 0
$$
がある。More on Algebra, Remark
\ref{more-algebra-remark-compare-derived-limit} を参照されたい。
(2) を示すため、左端の項が零であり、右端の項が
$H^i(RF(E_{-b + N - 1})$ に等しいことを示す。% P09-ERR-0395
ここで $b$ は任意であり、$i \geq b$ を満たす。

\medskip\noindent
各 $n$ に対し、区別三角
$$
H^{-n}(E)[n] \to E_n \to E_{n - 1} \to H^{-n}(E)[n + 1]
$$
が $D(\mathcal{O}_X)$ において存在する（Derived Categories, Remark
\ref{derived-remark-truncation-distinguished-triangle}）。
$H^{-n}(E)$ は準連接であるから、
$$
H^i(RF(H^{-n}(E)[n])) = R^{i + n}F(H^{-n}(E)) = 0
$$
は $i + n \geq N$ のとき成り立ち、
$$
H^i(RF(H^{-n}(E)[n + 1])) = R^{i + n + 1}F(H^{-n}(E)) = 0
$$
は $i + n + 1 \geq N$ のとき成り立つ。したがって
$$
H^i(RF(E_n)) \to H^i(RF(E_{n - 1}))
$$
は $n \geq N - i$ のとき同型である。
ゆえに系 $H^i(RF(E_n))$ はすべて ML 条件を満たし、
上の短完全列における $R^1\lim$ の項は零である
（More on Algebra, Section
\ref{more-algebra-section-Rlim} の議論を参照）。
さらに、系 $H^i(RF(E_n))$ は所望のとおり
$n = N - i - 1$ から定値である。

\medskip\noindent
(3) を証明する。$E$ に関する仮定のもとで
$\tau_{\leq a - 1}E = 0$ であり、(1) から
$H^i(RF(E))$ の $i \leq a - 1$ における消滅を得る。
同様に $\tau_{\geq b + 1}E = 0$ であるから、(2) により
$H^i(RF(E))$ の $i \geq b + n$ における消滅を得る。
% P09-ERR-0396
\end{proof}









\section{全順像}
\label{spaces-perfect-section-total-direct-image}

\noindent
次の補題は Cohomology of Spaces, Lemma
\ref{spaces-cohomology-lemma-vanishing-higher-direct-images}
の類似である。

\begin{lemma}
\label{spaces-perfect-lemma-quasi-coherence-direct-image}
$S$ を概型とし、$f : X \to Y$ を
$S$ 上の代数空間の準分離かつ準コンパクトな射とする。
\begin{enumerate}
\item 函手 $Rf_*$ は $D_\QCoh(\mathcal{O}_X)$ を
$D_\QCoh(\mathcal{O}_Y)$ の中へ送る。
\item $Y$ が準コンパクトならば、整数 $N = N(X, Y, f)$ が存在して、
対象 $E$（$D_\QCoh(\mathcal{O}_X)$ の対象）が
$H^m(E) = 0$（$m > 0$）を満たすとき、
$H^m(Rf_*E) = 0$（$m \geq N$）となる。
\item 実際、$Y$ が準コンパクトならば、$N = N(X, Y, f)$ を、
代数空間の任意の射 $Y' \to Y$ に対して函手 $R(f')_*$ が
同じ結論を満たすように選べる。ここで
$f' : X' \to Y'$ は $f$ の基底変換である。
\end{enumerate}
\end{lemma}

\begin{proof}
$E$ を $D_\QCoh(\mathcal{O}_X)$ の対象とする。
(1) を示すには、$Rf_*E$ が準連接なコホモロジー層をもつことを
示せばよい。この問題は $Y$ 上で局所的であるから、
$Y$ は準コンパクトであると仮定してよい。
Cohomology of Spaces, Lemma
\ref{spaces-cohomology-lemma-vanishing-higher-direct-images}
における $N = N(X, Y, f)$ を取る。すると
$R^pf_*\mathcal{F} = 0$ である。ここで $\mathcal{O}_X$-加群
$\mathcal{F}$ は任意の準連接加群であり、$p \geq N$ とする。
さらに Cohomology of Spaces, Lemma
\ref{spaces-cohomology-lemma-higher-direct-image} により、
$R^pf_*\mathcal{F}$ はすべての $p$ に対して準連接である。
これらの主張は基底変換後にも成り立つ。

\medskip\noindent
まず $E$ が下に有界であると仮定する。このような $E$ に対し、
選んだ $N$ によって (1)、(2)、(3) が成り立つことを示す。
この場合、例えばスペクトル系列
$$
R^pf_*H^q(E) \Rightarrow R^{p + q}f_*E
$$
（Derived Categories, Lemma
\ref{derived-lemma-two-ss-complex-functor}）と、
$R^pf_*H^q(E)$ の準連接性、および
$R^pf_*H^q(E)$ の $p \geq N$ における消滅を用いれば、
(1)、(2)、(3) がこの場合に成り立つことが分かる。

\medskip\noindent
次に (2) と (3) を証明する。$H^m(E) = 0$ が $m > 0$ に対して
成り立つとする。$V$ を $Y_\etale$ のアフィン対象とする。
Cohomology of Spaces, Lemma
\ref{spaces-cohomology-lemma-quasi-coherence-higher-direct-images-application}
により、$H^p(V \times_Y X, \mathcal{F}) = 0$ が
$p \geq N$ に対して成り立つ。
したがって函手 $\Gamma(V \times_Y X, -)$ に
Lemma \ref{spaces-perfect-lemma-application-nice-K-injective} を適用でき、
$$
R\Gamma(V, Rf_*E) = R\Gamma(V \times_Y X, E)
$$
の次数 $\geq N$ のコホモロジーは消滅する。
これはすべての $V$（$Y_\etale$ のアフィン対象）に対して成り立つから、
$H^m(Rf_*E) = 0$ が $m \geq N$ に対して成り立つと結論される。

\medskip\noindent
最後に、一般の場合の (1) を証明する。区別三角
$$
\tau_{\leq -n - 1}E \to E \to \tau_{\geq -n}E \to
(\tau_{\leq -n - 1}E)[1]
$$
が $D(\mathcal{O}_X)$ において存在することを思い出す。
Derived Categories, Remark
\ref{derived-remark-truncation-distinguished-triangle} を参照されたい。
(2) により、$Rf_*\tau_{\leq -n - 1}E$ のコホモロジー層は
次数 $\geq -n + N$ で消滅する。
したがって整数 $q$ が与えられたとき、
$R^qf_*E$ は $R^qf_*\tau_{\geq -n}E$ に、ある $n$ に対して等しくなり、
上の結果を適用できる。
\end{proof}

\begin{lemma}
\label{spaces-perfect-lemma-quasi-coherence-pushforward-direct-sums}
$S$ を概型とし、$f : X \to Y$ を $S$ 上の代数空間の
準分離かつ準コンパクトな射とする。このとき
$Rf_* : D_\QCoh(\mathcal{O}_X) \to D_\QCoh(\mathcal{O}_Y)$
は直和と可換である。
\end{lemma}

\begin{proof}
$E_i$ を $D_\QCoh(\mathcal{O}_X)$ の対象の族とし、
$E = \bigoplus E_i$ とおく。射
$$
\bigoplus Rf_*E_i \longrightarrow Rf_*E
$$
が同型であることを示したい。この射が次数 $0$ のコホモロジー層上で
同型を誘導することを示せば、補題が従う。
Lemma \ref{spaces-perfect-lemma-quasi-coherence-direct-image} における整数 $N$ を選ぶ。
すると、同補題により
$R^0f_*E = R^0f_*\tau_{\geq -N}E$ および
$R^0f_*E_i = R^0f_*\tau_{\geq -N}E_i$ である。また
$\tau_{\geq -N}E = \bigoplus \tau_{\geq -N}E_i$ である。
したがって、すべての $E_i$ のコホモロジー層が次数 $< -N$
で消滅すると仮定してよい。次にスペクトル系列
$$
R^pf_*H^q(E) \Rightarrow R^{p + q}f_*E
\quad\text{and}\quad
R^pf_*H^q(E_i) \Rightarrow R^{p + q}f_*E_i
$$
（Derived Categories, Lemma
\ref{derived-lemma-two-ss-complex-functor}）を用い、
準連接層の直和の場合に帰着する。この場合は
Cohomology of Spaces, Lemma
\ref{spaces-cohomology-lemma-colimit-cohomology} によって処理される。
\end{proof}

\begin{remark}
\label{spaces-perfect-remark-match-total-direct-images}
$S$ を概型とする。$f : X \to Y$ を、表現可能な
代数空間 $X$ と $Y$（$S$ 上のもの）の間の射とする。
$f_0 : X_0 \to Y_0$ を $f$ を表現する概型の射とする
（不格好だが一時的な記法である）。このとき図式
$$
\xymatrix{
D_\QCoh(\mathcal{O}_{X_0})
\ar@{=}[rrrrrr]_{\text{Lemma
\ref{spaces-perfect-lemma-derived-quasi-coherent-small-etale-site}}}
& & & & & &
D_\QCoh(\mathcal{O}_X) \\
D_\QCoh(\mathcal{O}_{Y_0})
\ar[u]^{Lf^*_0}
\ar@{=}[rrrrrr]^{\text{Lemma
\ref{spaces-perfect-lemma-derived-quasi-coherent-small-etale-site}}}
& & & & & &
D_\QCoh(\mathcal{O}_Y) \ar[u]_{Lf^*}
}
$$
は可換である（Lemma \ref{spaces-perfect-lemma-quasi-coherence-pullback} および
Derived Categories of Schemes, Lemma
\ref{perfect-lemma-quasi-coherence-pullback}）。
これは、Lemma
\ref{spaces-perfect-lemma-derived-quasi-coherent-small-etale-site} の圏同値
$D_\QCoh(\mathcal{O}_{X_0}) \to D_\QCoh(\mathcal{O}_X)$
および
$D_\QCoh(\mathcal{O}_{Y_0}) \to D_\QCoh(\mathcal{O}_Y)$
が、環付きサイトの（平坦な）射
$\epsilon : X_\etale \to X_{0, Zar}$ および
$\epsilon : Y_\etale \to Y_{0, Zar}$
による引き戻しから生じ、環付きサイトの図式
$$
\xymatrix{
X_{0, Zar} \ar[d]_{f_0} & X_\etale \ar[l]^\epsilon \ar[d]^f \\
Y_{0, Zar} & Y_\etale \ar[l]_\epsilon
}
$$
が可換だからである（詳細は省略する）。$f$ が準コンパクトかつ
準分離であるとする。これは $f_0$ が準コンパクトかつ準分離であることと
同値である。このとき図式
$$
\xymatrix{
D_\QCoh(\mathcal{O}_{X_0})
\ar[d]_{Rf_{0, *}} \ar@{=}[rrrrrr]_{\text{Lemma
\ref{spaces-perfect-lemma-derived-quasi-coherent-small-etale-site}}}
& & & & & &
D_\QCoh(\mathcal{O}_X) \ar[d]^{Rf_*} \\
D_\QCoh(\mathcal{O}_{Y_0})
\ar@{=}[rrrrrr]^{\text{Lemma
\ref{spaces-perfect-lemma-derived-quasi-coherent-small-etale-site}}}
& & & & & &
D_\QCoh(\mathcal{O}_Y)
}
$$
も可換であると主張する（Lemma
\ref{spaces-perfect-lemma-quasi-coherence-direct-image} および
Derived Categories of Schemes, Lemma
\ref{perfect-lemma-quasi-coherence-direct-image}）。
これも上に掲げたサイトの可換図式から従う。実際、Lemma
\ref{spaces-perfect-lemma-derived-quasi-coherent-small-etale-site} の証明が示すように、
函手 $R\epsilon_*$ は圏同値
$D_\QCoh(\mathcal{O}_X) \to D_\QCoh(\mathcal{O}_{X_0})$
および
$D_\QCoh(\mathcal{O}_Y) \to D_\QCoh(\mathcal{O}_{Y_0})$
を与える。
\end{remark}


\begin{lemma}
\label{spaces-perfect-lemma-affine-morphism}
$S$ を概型とし、$f : X \to Y$ を $S$ 上の代数空間の
アフィン射とする。このとき
$Rf_* : D_\QCoh(\mathcal{O}_X) \to D_\QCoh(\mathcal{O}_Y)$
は同型を反映する。
\end{lemma}

\begin{proof}
この主張は、射 $\alpha : E \to F$（
$D_\QCoh(\mathcal{O}_X)$ の射）が、$Rf_*\alpha$ が同型であるならば
同型であることを意味する。これはコホモロジー層上で確認できる。
特に問題は $Y$ 上で平展局所的である。したがって $Y$、ひいては
$X$ がアフィンであると仮定してよい。この場合は概型の場合
（Derived Categories of Schemes, Lemma
\ref{perfect-lemma-affine-morphism}）に帰着する。ここで Lemma
\ref{spaces-perfect-lemma-derived-quasi-coherent-small-etale-site} と Remark
\ref{spaces-perfect-remark-match-total-direct-images} を用いた。
\end{proof}

\begin{lemma}
\label{spaces-perfect-lemma-affine-morphism-pull-push}
$S$ を概型とし、$f : X \to Y$ を $S$ 上の代数空間の
アフィン射とする。$E$（$D_\QCoh(\mathcal{O}_Y)$ の対象）に対して
$Rf_* Lf^* E = E \otimes^\mathbf{L}_{\mathcal{O}_Y} f_*\mathcal{O}_X$
である。
\end{lemma}

\begin{proof}
$f$ はアフィンであるから、射
$f_*\mathcal{O}_X \to Rf_*\mathcal{O}_X$
は同型である（Cohomology of Spaces, Lemma
\ref{spaces-cohomology-lemma-affine-vanishing-higher-direct-images}）。
標準的な射
$E \otimes^\mathbf{L} f_*\mathcal{O}_X =
E \otimes^\mathbf{L} Rf_*\mathcal{O}_X \to Rf_* Lf^* E$
がある。これは射
$$
Lf^*(E \otimes^\mathbf{L} Rf_*\mathcal{O}_X) =
Lf^*E \otimes^\mathbf{L} Lf^*Rf_*\mathcal{O}_X \longrightarrow
Lf^* E \otimes^\mathbf{L} \mathcal{O}_X = Lf^* E
$$
に随伴する。後者は
$1 : Lf^*E \to Lf^*E$ と標準的な射
$Lf^*Rf_*\mathcal{O}_X \to \mathcal{O}_X$ から生じる。
こうして構成した射が同型であることは $Y$ 上で局所的に確認できる。
したがって $Y$、ひいては $X$ がアフィンであると仮定してよい。
この場合は概型の場合
（Derived Categories of Schemes, Lemma
\ref{perfect-lemma-affine-morphism-pull-push}）に帰着する。ここで Lemma
\ref{spaces-perfect-lemma-derived-quasi-coherent-small-etale-site} と Remark
\ref{spaces-perfect-remark-match-total-direct-images} を用いた。
\end{proof}









\section{基底上固有であること}
\label{spaces-perfect-section-proper-over-base}

\noindent
本節は Cohomology of Schemes, Section
\ref{coherent-section-proper-over-base} の類似である。
点集合上の位相に関係する事項では常にそうであるように、
この内容を代数空間へ移す際には注意が必要である。

\begin{lemma}
\label{spaces-perfect-lemma-closed-proper-over-base}
$S$ を概型とする。$f : X \to Y$ を、$S$ 上の代数空間の
局所有限型射とする。$T \subset |X|$ を閉部分集合とする。
次の条件は同値である。
\begin{enumerate}
\item 射 $Z \to Y$ は、$Z$ が $T$ 上の被約誘導代数空間構造であるとき固有である
（Properties of Spaces, Definition
\ref{spaces-properties-definition-reduced-induced-space}）。
\item ある閉部分空間 $Z \subset X$ で $|Z| = T$ を満たすものに対して、
射 $Z \to Y$ は固有である。
\item 任意の閉部分空間 $Z \subset X$ で $|Z| = T$ を満たすものに対して、
射 $Z \to Y$ は固有である。
\end{enumerate}
\end{lemma}

\begin{proof}
(3) $\Rightarrow$ (1) および (1) $\Rightarrow$ (2) は直ちに従う。
したがって (2) が (3) を含意することを示せば十分である。
この事実の証明をまず読者自身で考えることを勧める。
$Z'$ と $Z''$ を $T = |Z'| = |Z''|$ を満たす閉部分空間とし、
$Z' \to Y$ は代数空間の固有射であるとする。
$Z'' \to Y$ も固有であることを示す必要がある。
$Z''' = Z' \cup Z''$ をスキーム論的和とする。Morphisms of Spaces,
Definition
\ref{spaces-morphisms-definition-scheme-theoretic-intersection-union}
を参照されたい。このとき $Z'''$ は $|Z'''| = T$ を満たす
別の閉部分空間である。これは、例えば Morphisms of Spaces, Lemma
\ref{spaces-morphisms-lemma-scheme-theoretic-union}
におけるスキーム論的和の記述から従う。
$Z'' \to Z'''$ は閉埋め込みであるから、
$Z''' \to Y$ が固有であることを示せば十分である
（Morphisms of Spaces, Lemmas
\ref{spaces-morphisms-lemma-closed-immersion-proper} および
\ref{spaces-morphisms-lemma-composition-proper}）。
射 $Z' \to Z'''$ は全単射な閉埋め込みであり、特に全射かつ普遍閉である。
したがって、$Z' \to Y$ が分離的であることから、
Morphisms of Spaces, Lemma
\ref{spaces-morphisms-lemma-image-universally-closed-separated}
により $Z''' \to Y$ は分離的である。
さらに、$Z''' \to Y$ は局所有限型である。これは
$X \to Y$ が局所有限型だからである
（Morphisms of Spaces, Lemmas
\ref{spaces-morphisms-lemma-immersion-locally-finite-type} および
\ref{spaces-morphisms-lemma-composition-finite-type}）。
$Z' \to Y$ は準コンパクトであり、$Z' \to Z'''$ は普遍同相であるから、
$Z''' \to Y$ は準コンパクトである。最後に $Z' \to Y$ は普遍閉であるから、
Morphisms of Spaces, Lemma
\ref{spaces-morphisms-lemma-image-proper-is-proper}
により $Z''' \to Y$ も普遍閉である。これで証明が完了する。
\end{proof}

\begin{definition}
\label{spaces-perfect-definition-proper-over-base}
$S$ を概型とする。$f : X \to Y$ を、$S$ 上の代数空間の
局所有限型射とする。$T \subset |X|$ を閉部分集合とする。
Lemma \ref{spaces-perfect-lemma-closed-proper-over-base} の同値な条件が満たされるとき、
{\it $T$ は $Y$ 上固有である}という。
\end{definition}

\noindent
上の定義に用いた補題は、射 $f : X \to Y$ が局所有限型でなければ
偽である。したがって、$f$ が局所有限型でない場合には
この用語を使用しないよう強く勧める。

\begin{lemma}
\label{spaces-perfect-lemma-closed-closed-proper-over-base}
$S$ を概型とする。$f : X \to Y$ を、$S$ 上の代数空間の
局所有限型射とする。$T' \subset T \subset |X|$ を閉部分集合とする。
$T$ が $Y$ 上固有ならば、$T'$ についても同様である。
\end{lemma}

\begin{proof}
省略する。
\end{proof}

\begin{lemma}
\label{spaces-perfect-lemma-base-change-closed-proper-over-base}
$S$ を概型とする。$S$ 上の代数空間の Cartesian 図式
$$
\xymatrix{
X' \ar[d]_{f'} \ar[r]_{g'} & X \ar[d]^f \\
Y' \ar[r]^g & Y
}
$$
を考える。ここで $f$ は局所有限型である。
$T$ が $|X|$ の閉部分集合で $Y$ 上固有ならば、
$|g'|^{-1}(T)$ は $|X'|$ の閉部分集合で $Y'$ 上固有である。
\end{lemma}

\begin{proof}
Morphisms of Spaces, Lemma
\ref{spaces-morphisms-lemma-base-change-finite-type} により
$f'$ は局所有限型であるから、この主張は意味をもつ。
$Z \subset X$ を $T$ 上の被約誘導閉部分空間構造とする。
$Z' = (g')^{-1}(Z)$ をスキーム論的逆像と表す。このとき
$Z' = X' \times_X Z = (Y' \times_Y X) \times_X Z = Y' \times_Y Z$
は $Y'$ 上固有である。これは $Z$（$Y$ 上のもの）の基底変換だからである
（Morphisms of Spaces, Lemma
\ref{spaces-morphisms-lemma-base-change-proper}）。
一方、$T' = |Z'|$ である。% P09-ERR-0397
したがって補題が従う。
\end{proof}

\begin{lemma}
\label{spaces-perfect-lemma-functoriality-closed-proper-over-base}
$S$ を概型、$B$ を $S$ 上の代数空間とする。
$f : X \to Y$ を、いずれも $B$ 上局所有限型な代数空間の射とする。
\begin{enumerate}
\item $Y$ が $B$ 上分離的で、$T \subset |X|$ が
$B$ 上固有な閉部分集合ならば、$|f|(T)$ は $|Y|$ の閉部分集合で
$B$ 上固有である。
\item $f$ が普遍閉で、$T \subset |X|$ が $B$ 上固有な
閉部分集合ならば、$|f|(T)$ は $Y$ の閉部分集合で
$B$ 上固有である。% P09-ERR-0398
\item $f$ が固有で、$T \subset |Y|$ が $B$ 上固有な
閉部分集合ならば、$|f|^{-1}(T)$ は $|X|$ の閉部分集合で
$B$ 上固有である。
\end{enumerate}
\end{lemma}

\begin{proof}
(1) を証明する。$Y$ は $B$ 上分離的であり、
$T \subset |X|$ は $B$ 上固有な閉部分集合であるとする。
$Z$ を $T$ 上の被約誘導閉部分空間構造とし、
Morphisms of Spaces, Lemma
\ref{spaces-morphisms-lemma-scheme-theoretic-image-is-proper}
を $Z \to Y$ に $B$ 上で適用すればよい。

\medskip\noindent
(2) を証明する。$f$ は普遍閉であり、$T \subset |X|$ は
$B$ 上固有な閉部分集合であるとする。
$Z$ を $T$ 上の被約誘導閉部分空間構造とし、
$Z'$ を $|f|(T)$ 上の被約誘導閉部分空間構造とする。
誘導される射 $Z \to Z'$ を得る。
$Z'' = f^{-1}(Z')$ をスキーム論的逆像と表す。
すると $Z'' \to Z'$ は $f$ の基底変換として普遍閉である
（Morphisms of Spaces, Lemma
\ref{spaces-morphisms-lemma-base-change-proper}）。
したがって $Z \to Z'$ は、閉埋め込み $Z \to Z''$ と
$Z'' \to Z'$ の合成として普遍閉である
（Morphisms of Spaces, Lemmas
\ref{spaces-morphisms-lemma-closed-immersion-proper} および
\ref{spaces-morphisms-lemma-composition-proper}）。
Morphisms of Spaces, Lemma
\ref{spaces-morphisms-lemma-image-universally-closed-separated}
により、$Z' \to B$ は分離的である。
$Z \to B$ は準コンパクトで $Z \to Z'$ は全射であるから、
$Z' \to B$ は準コンパクトである。
$Z' \to B$ は $Z' \to Y$ と $Y \to B$ の合成であるから、
$Z' \to B$ は局所有限型である
（Morphisms of Spaces, Lemmas
\ref{spaces-morphisms-lemma-immersion-locally-finite-type} および
\ref{spaces-morphisms-lemma-composition-finite-type}）。
最後に $Z \to B$ は普遍閉であるから、
Morphisms of Spaces, Lemma
\ref{spaces-morphisms-lemma-image-proper-is-proper}
により $Z' \to B$ も普遍閉である。これで証明が完了する。

\medskip\noindent
(3) を証明する。$f$ は固有であり、$T \subset |Y|$ は
$B$ 上固有な閉部分集合であるとする。
$Z$ を $T$ 上の被約誘導閉部分空間構造とし、
$Z' = f^{-1}(Z)$ をスキーム論的逆像と表す。
すると $Z' \to Z$ は $f$ の基底変換として固有である
（Morphisms of Spaces, Lemma
\ref{spaces-morphisms-lemma-base-change-proper}）。
したがって $Z' \to B$ は $Z' \to Z$ と $Z \to B$ の合成として
固有である（Morphisms of Spaces, Lemma
\ref{spaces-morphisms-lemma-composition-proper}）。
これで証明が完了する。
\end{proof}

\begin{lemma}
\label{spaces-perfect-lemma-union-closed-proper-over-base}
$S$ を概型とする。$f : X \to Y$ を、$S$ 上の代数空間の
局所有限型射とする。$T_i \subset |X|$（$i = 1, \ldots, n$）
を閉部分集合とする。$T_i$（$i = 1, \ldots, n$）が
$Y$ 上固有ならば、$T_1 \cup \ldots \cup T_n$ についても同様である。
\end{lemma}

\begin{proof}
$Z_i$ を $T_i$ 上の被約誘導閉部分スキーム構造とする。
% P09-ERR-0399
射
$$
Z_1 \amalg \ldots \amalg Z_n \longrightarrow X
$$
は Morphisms of Spaces, Lemmas
\ref{spaces-morphisms-lemma-closed-immersion-finite} および
\ref{spaces-morphisms-lemma-finite-union-finite} により有限である。
有限射は普遍閉であり
（Morphisms of Spaces, Lemma
\ref{spaces-morphisms-lemma-finite-proper}）、かつ
$Z_1 \amalg \ldots \amalg Z_n$ は $S$ 上固有であるから、
% P09-ERR-0400
Lemma \ref{spaces-perfect-lemma-functoriality-closed-proper-over-base} の (2) により、
像 $Z_1 \cup \ldots \cup Z_n$ は $S$ 上固有である。
\end{proof}

\noindent
$S$ を概型とする。$f : X \to Y$ を、$S$ 上の代数空間の
局所有限型射とする。$\mathcal{F}$ を有限型準連接
$\mathcal{O}_X$-加群とする。このとき台
$\text{Supp}(\mathcal{F})$（$\mathcal{F}$ の台）は $|X|$ の閉部分集合である。
Morphisms of Spaces, Lemma
\ref{spaces-morphisms-lemma-support-finite-type} を参照されたい。
したがって「$\mathcal{F}$ の台は $Y$ 上固有である」と
いうことが意味をもつ。

\begin{lemma}
\label{spaces-perfect-lemma-module-support-proper-over-base}
$S$ を概型とし、$f : X \to Y$ を、$S$ 上の代数空間の
局所有限型射とする。$\mathcal{F}$ を有限型準連接
$\mathcal{O}_X$-加群とする。次の条件は同値である。
\begin{enumerate}
\item $\mathcal{F}$ の台は $Y$ 上固有である。
\item $\mathcal{F}$ のスキーム論的台
（Morphisms of Spaces, Definition
\ref{spaces-morphisms-definition-scheme-theoretic-support}）
は $Y$ 上固有である。
\item 閉部分空間 $Z \subset X$ と有限型準連接
$\mathcal{O}_Z$-加群 $\mathcal{G}$ が存在して、
(a) $Z \to Y$ は固有であり、(b)
$(Z \to X)_*\mathcal{G} = \mathcal{F}$ である。
\end{enumerate}
\end{lemma}

\begin{proof}
台 $\text{Supp}(\mathcal{F})$（$\mathcal{F}$ の台）は
$|X|$ の閉部分集合である。Morphisms of Spaces, Lemma
\ref{spaces-morphisms-lemma-support-finite-type} を参照されたい。
したがって Definition \ref{spaces-perfect-definition-proper-over-base} を適用できる。
$\mathcal{F}$ のスキーム論的台は、その基礎閉部分集合が
$\text{Supp}(\mathcal{F})$ である閉部分空間だから、
Definition \ref{spaces-perfect-definition-proper-over-base} により (1) と (2) は同値である。
(2) が (3) を含意することは明らかである。逆に (3) が成り立つならば
$\text{Supp}(\mathcal{F}) \subset |Z|$ であり、したがって
$\text{Supp}(\mathcal{F})$ は、例えば Lemma
\ref{spaces-perfect-lemma-closed-closed-proper-over-base} により $Y$ 上固有である。
\end{proof}

\begin{lemma}
\label{spaces-perfect-lemma-base-change-module-support-proper-over-base}
$S$ を概型とする。$S$ 上の代数空間の Cartesian 図式
$$
\xymatrix{
X' \ar[d]_{f'} \ar[r]_{g'} & X \ar[d]^f \\
Y' \ar[r]^g & Y
}
$$
を考える。ここで $f$ は局所有限型である。
$\mathcal{F}$ を有限型準連接 $\mathcal{O}_X$-加群とする。
$\mathcal{F}$ の台が $Y$ 上固有ならば、
$(g')^*\mathcal{F}$ の台は $Y'$ 上固有である。
\end{lemma}

\begin{proof}
この主張が意味をもつことに注意する。実際
$(g')*\mathcal{F}$ は Modules on Sites, Lemma
\ref{sites-modules-lemma-local-pullback} により有限型である。
% P09-ERR-0401
Morphisms of Spaces, Lemma
\ref{spaces-morphisms-lemma-support-finite-type} により
$\text{Supp}((g')^*\mathcal{F}) = |g'|^{-1}(\text{Supp}(\mathcal{F}))$
である。したがって補題は Lemma
\ref{spaces-perfect-lemma-base-change-closed-proper-over-base} から従う。
\end{proof}

\begin{lemma}
\label{spaces-perfect-lemma-cat-module-support-proper-over-base}
$S$ を概型とする。$f : X \to Y$ を、$S$ 上の代数空間の
局所有限型射とする。$\mathcal{F}$、$\mathcal{G}$ を有限型準連接
$\mathcal{O}_X$-加群とする。
\begin{enumerate}
\item $\mathcal{F}$、$\mathcal{G}$ の台が $Y$ 上固有ならば、
$\mathcal{F} \oplus \mathcal{G}$、$\mathcal{G}$ の
$\mathcal{F}$ による任意の拡大、$\Im(u)$ と $\Coker(u)$
（任意の $\mathcal{O}_X$-加群の射 $u : \mathcal{F} \to \mathcal{G}$
に対するもの）、および
$\mathcal{F}$ または $\mathcal{G}$ の任意の準連接商の台も
同じ性質をもつ。
\item $Y$ が局所 Noether 的ならば、連接 $\mathcal{O}_X$-加群で
台が $Y$ 上固有なものの圏は、連接
$\mathcal{O}_X$-加群のアーベル圏の Serre 部分圏である
（Homology, Definition
\ref{homology-definition-serre-subcategory}）。
\end{enumerate}
\end{lemma}

\begin{proof}
(1) を証明する。$T$、$T'$ をそれぞれ
$\mathcal{F}$ と $\mathcal{G}$ の台とする。(1) に現れる層は
すべて $T \cup T'$ に含まれる台をもつ。
したがって、これらの層が有限型かつ準連接であることを確認すれば、
主張そのものは Lemmas
\ref{spaces-perfect-lemma-closed-closed-proper-over-base} および
\ref{spaces-perfect-lemma-union-closed-proper-over-base} から明らかである。
準連接性については Properties of Spaces, Section
\ref{spaces-properties-section-quasi-coherent} を参照されたい。
「有限型」については Properties of Spaces, Section
\ref{spaces-properties-section-properties-modules} を参照されたい。

\medskip\noindent
(2) を証明する。証明は (1) の証明と同じである。
Morphisms of Spaces, Lemma
\ref{spaces-morphisms-lemma-locally-finite-type-locally-noetherian}
および Cohomology of Spaces, Section
\ref{spaces-cohomology-section-coherent} における
連接加群の圏の記述により $X$ は局所 Noether 的であるから、
各主張は意味をもつことに注意する。
\end{proof}

\begin{lemma}
\label{spaces-perfect-lemma-support-proper-over-base-pushforward}
$S$ を概型とし、$f : X \to Y$ を $S$ 上の代数空間の射とする。
$f$ は局所有限型で、$Y$ は局所 Noether 的であると仮定する。
$\mathcal{F}$ を連接 $\mathcal{O}_X$-加群とし、その台は $Y$ 上固有とする。
このとき $R^pf_*\mathcal{F}$ は連接 $\mathcal{O}_Y$-加群である
（すべての $p \geq 0$ に対して）。
\end{lemma}

\begin{proof}
Lemma \ref{spaces-perfect-lemma-module-support-proper-over-base} により、
閉埋め込み $i : Z \to X$ が存在して、
$g = f \circ i : Z \to Y$ は固有であり、
$\mathcal{F} = i_*\mathcal{G}$ となる連接加群 $\mathcal{G}$（$Z$ 上）が存在する。
Cohomology of Spaces, Lemma
\ref{spaces-cohomology-lemma-proper-pushforward-coherent} により
$R^pg_*\mathcal{G}$ は $S$ 上連接である。% P09-ERR-0402
一方、$R^qi_*\mathcal{G} = 0$ である（$q > 0$ に対して）
（Cohomology of Spaces, Lemma
\ref{spaces-cohomology-lemma-finite-pushforward-coherent}）。
Cohomology on Sites, Lemma
\ref{sites-cohomology-lemma-relative-Leray} により
$R^pf_*\mathcal{F} = R^pg_*\mathcal{G}$ を得て、補題が従う。
\end{proof}






\section{連接加群の導来圏}
\label{spaces-perfect-section-derived-coherent}

\noindent
$S$ を概型、$X$ を $S$ 上の局所 Noether 的代数空間とする。
この場合、圏
$\textit{Coh}(\mathcal{O}_X) \subset \textit{Mod}(\mathcal{O}_X)$
（連接 $\mathcal{O}_X$-加群の圏）
は弱 Serre 部分圏である。Homology, Section
\ref{homology-section-serre-subcategories} および
Cohomology of Spaces, Lemma
\ref{spaces-cohomology-lemma-coherent-abelian-Noetherian}
を参照されたい。
$$
D_{\textit{Coh}}(\mathcal{O}_X) \subset D(\mathcal{O}_X)
$$
によって、コホモロジー層が連接である複体の部分圏を表す。
Derived Categories, Section
\ref{derived-section-triangulated-sub} を参照されたい。
こうして標準的な函手
\begin{equation}
\label{spaces-perfect-equation-compare-coherent}
D(\textit{Coh}(\mathcal{O}_X))
\longrightarrow
D_{\textit{Coh}}(\mathcal{O}_X)
\end{equation}
を得る。Derived Categories, Equation
(\ref{derived-equation-compare}) を参照されたい。

\begin{lemma}
\label{spaces-perfect-lemma-direct-image-coherent}
$S$ を概型とし、$f : X \to Y$ を $S$ 上の代数空間の射とする。
$f$ は局所有限型で、$Y$ は Noether 的であると仮定する。
$E$ を $D^b_{\textit{Coh}}(\mathcal{O}_X)$ の対象とし、
$H^i(E)$ の台は $Y$ 上固有であるとする（すべての $i$ に対して）。
このとき $Rf_*E$ は
$D^b_{\textit{Coh}}(\mathcal{O}_Y)$ の対象である。
\end{lemma}

\begin{proof}
スペクトル系列
$$
R^pf_*H^q(E) \Rightarrow R^{p + q}f_*E
$$
を考える。Derived Categories, Lemma
\ref{derived-lemma-two-ss-complex-functor} を参照されたい。
仮定と Lemma \ref{spaces-perfect-lemma-support-proper-over-base-pushforward} により
層 $R^pf_*H^q(E)$ は連接である。したがって
$R^{p + q}f_*E$ は連接であり、すなわち
$E \in D_{\textit{Coh}}(\mathcal{O}_Y)$ である。% P09-ERR-0403
下有界性は明らかである。上有界性は Cohomology of Spaces, Lemma
\ref{spaces-cohomology-lemma-vanishing-higher-direct-images}
または Lemma \ref{spaces-perfect-lemma-quasi-coherence-direct-image} から従う。
\end{proof}

\begin{lemma}
\label{spaces-perfect-lemma-direct-image-coherent-bdd-below}
$S$ を概型とし、$f : X \to Y$ を $S$ 上の代数空間の射とする。
$f$ は局所有限型で、$Y$ は Noether 的であると仮定する。
$E$ を $D^+_{\textit{Coh}}(\mathcal{O}_X)$ の対象とし、
$H^i(E)$ の台は $S$ 上固有であるとする（すべての $i$ に対して）。
% P09-ERR-0404
このとき $Rf_*E$ は
$D^+_{\textit{Coh}}(\mathcal{O}_Y)$ の対象である。
\end{lemma}

\begin{proof}
証明は Lemma \ref{spaces-perfect-lemma-direct-image-coherent} の証明と同じである。
また、完全函手 $Rf_*$ が区別三角
$\tau_{\leq a}E \to E \to \tau_{\geq a + 1}E \to \tau_{\leq a}E[1]$
をどのように移すかを考えることにより、
Lemma \ref{spaces-perfect-lemma-direct-image-coherent} から導くこともできる。
\end{proof}

\begin{lemma}
\label{spaces-perfect-lemma-coherent-internal-hom}
$S$ を概型、$X$ を $S$ 上の局所 Noether 的代数空間とする。
$L$ が $D^+_{\textit{Coh}}(\mathcal{O}_X)$ に属し、
$K$ が $D^-_{\textit{Coh}}(\mathcal{O}_X)$ に属するならば、
$R\SheafHom(K, L)$ は
$D^+_{\textit{Coh}}(\mathcal{O}_X)$ に属する。
\end{lemma}

\begin{proof}
$D(\mathcal{O}_X)$ の対象が
$D_{\textit{Coh}}(\mathcal{O}_X)$ に属するかどうかは
$X$ 上で平展局所的に確認できる。Cohomology of Spaces, Lemma
\ref{spaces-cohomology-lemma-coherent-Noetherian} を参照されたい。
したがって補題は概型の場合、すなわち
Derived Categories of Schemes, Lemma
\ref{perfect-lemma-coherent-internal-hom} から従う。
\end{proof}

\begin{lemma}
\label{spaces-perfect-lemma-ext-finite}
$A$ を Noether 環、$X$ を $A$ 上の固有な代数空間とする。
$L$（$D^+_{\textit{Coh}}(\mathcal{O}_X)$ に属するもの）と
$K$（$D^-_{\textit{Coh}}(\mathcal{O}_X)$ に属するもの）に対して、
$A$-加群 $\Ext_{\mathcal{O}_X}^n(K, L)$ は有限である。
\end{lemma}

\begin{proof}
等式
$$
\Ext_{\mathcal{O}_X}^n(K, L) =
H^n(X, R\SheafHom_{\mathcal{O}_X}(K, L)) =
H^n(\Spec(A), Rf_*R\SheafHom_{\mathcal{O}_X}(K, L))
$$
を思い出す。Cohomology on Sites, Lemma
\ref{sites-cohomology-lemma-section-RHom-over-U} および
Cohomology on Sites, Section
\ref{sites-cohomology-section-leray} を参照されたい。
したがって結論は Lemmas
\ref{spaces-perfect-lemma-coherent-internal-hom} および
\ref{spaces-perfect-lemma-direct-image-coherent-bdd-below} から従う。
\end{proof}




\section{帰納原理}
\label{spaces-perfect-section-induction}

\noindent
本節では、概型に対する Cohomology of Schemes, Lemma
\ref{coherent-lemma-induction-principle} と類似する、
代数空間の帰納原理を論じる。その定式化のために
{\it 初等識別正方形}という概念を導入する。この用語は
\cite{MV} から借用したものである。
ここで定式化する原理は、Raynaud と Gruson の論文
\cite{GruRay} に暗黙に含まれている。
代数スタックに対する関連する原理として、David Rydh による
\cite[Theorem D]{rydh_etale_devissage} がある。

\begin{definition}
\label{spaces-perfect-definition-elementary-distinguished-square}
$S$ を概型とする。代数空間の可換図式
$$
\xymatrix{
U \times_W V \ar[r] \ar[d] & V \ar[d]^f \\
U \ar[r]^j & W
}
$$
が $S$ 上の{\it 初等識別正方形}であるとは、次の条件が成り立つことをいう。
\begin{enumerate}
\item $U$ は $W$ の開部分空間で、$j$ は包含射である。
\item $f$ は平展である。
\item $T = W \setminus U$ と置き（被約誘導部分空間構造を入れる）、
射 $f^{-1}(T) \to T$ は同型である。
\end{enumerate}
これを「$(U \subset W, f : V \to W)$ を初等識別正方形とする」
と表す。
\end{definition}

\noindent
$(U \subset W, f : V \to W)$ が初等識別正方形ならば、
$W = U \cup f(V)$ であることに注意する。したがって
$\{U \to W, V \to W\}$ は $W$ の平展被覆である。
これらの平展被覆は良い性質をもち、ある意味で「十分多く」存在する。

\begin{lemma}
\label{spaces-perfect-lemma-make-more-elementary-distinguished-squares}
$S$ を概型とする。$(U \subset W, f : V \to W)$ を、
$S$ 上の代数空間の初等識別正方形とする。
\begin{enumerate}
\item $V' \subset V$ および
$U \subset U' \subset W$ が開部分空間で、$W' = U' \cup f(V')$
ならば、$(U' \subset W', f|_{V'} : V' \to W')$ は
初等識別正方形である。
\item $p : W' \to W$ が代数空間の射ならば、
$(p^{-1}(U) \subset W', V \times_W W' \to W')$ は
初等識別正方形である。
\item $S' \to S$ が概型の射ならば、
$(S' \times_S U \subset S' \times_S W, S' \times_S V \to S' \times_S W)$
は初等識別正方形である。
\end{enumerate}
\end{lemma}

\begin{proof}
省略する。
\end{proof}

\begin{lemma}
\label{spaces-perfect-lemma-induction-principle}
$S$ を概型とする。$X$ を $S$ 上の準コンパクトかつ準分離的な
代数空間とする。$P$ を $X_{spaces, \etale}$ の準コンパクトかつ
準分離的な対象の性質とする。次を仮定する。
\begin{enumerate}
\item $P$ は $X_{spaces, \etale}$ の任意のアフィン対象について成り立つ。
\item 各初等識別正方形 $(U \subset W, f : V \to W)$ について、
次が成り立つとする。
\begin{enumerate}
\item $W$ は $X_{spaces, \etale}$ の準コンパクトかつ準分離的な対象である。
\item $U$ は準コンパクトである。
\item $V$ はアフィンである。
\item $P$ は $U$、$V$、および $U \times_W V$ について成り立つ。
\end{enumerate}
このとき $P$ は $W$ について成り立つ。
\end{enumerate}
このとき $P$ は $X_{spaces, \etale}$ の任意の準コンパクトかつ
準分離的な対象について成り立ち、特に $X$ について成り立つ。
\end{lemma}

\begin{proof}
まず、$P$ は $X_{spaces, \etale}$ の表現可能な任意の準コンパクトかつ
準分離的な対象について成り立つと主張する。実際、$U \to X$ が平展で、
$U$ が準コンパクトかつ準分離的な概型であるとする。仮定 (1) により、
性質 $P$ は $U$ の任意のアフィン開部分について成り立つ。さらに、
$W, V \subset U$ が準コンパクト開部分で、$V$ がアフィンであり、
$P$ が $W$、$V$、および $W \cap V$ について成り立つならば、
(2) により $P$ は $W \cup V$ について成り立つ
（組 $(W \subset W \cup V, V \to W \cup V)$ は
初等識別正方形だからである）。したがって、概型に対する帰納原理
Cohomology of Schemes, Lemma
\ref{coherent-lemma-induction-principle} により、$P$ は $U$ について成り立つ。

\medskip\noindent
証明を終えるには、$P$ が $X$ について成り立つことを示せば十分である
（単に $X$ を、結論を示したい $X_{spaces, \etale}$ の任意の
準コンパクトかつ準分離的な対象で置き換えればよいからである）。
Decent Spaces, Lemma
\ref{decent-spaces-lemma-filter-quasi-compact-quasi-separated} のフィルトレーション
$$
\emptyset = U_{n + 1} \subset
U_n \subset U_{n - 1} \subset \ldots \subset U_1 = X
$$
および射 $f_p : V_p \to U_p$ を用いる。
$P$ が $U_p$ について成り立つことを、$p$ に関する
降下帰納法によって示す。空代数空間はアフィンだから、(1) により
$P$ は $U_{n + 1}$ について成り立つ。$P$ が $U_{p + 1}$ について
成り立つと仮定する。$(U_{p + 1} \subset U_p, f_p : V_p \to U_p)$ は
初等識別正方形だが、$V_p$ はアフィンとは限らないため (2) を
適用できないことがある。しかし $V_p$ は準コンパクト概型なので、
有限アフィン開被覆
$V_p = V_{p, 1} \cup \ldots \cup V_{p, m}$ を選べる。
$W_{p, 0} = U_{p + 1}$ と置き、
$$
W_{p, i} = U_{p + 1} \cup f_p(V_{p, 1} \cup \ldots \cup V_{p, i})
$$
と定める（$i = 1, \ldots, m$）。これらは $X$ の準コンパクト
開部分空間である。このとき
$$
U_{p + 1} = W_{p, 0} \subset
W_{p, 1} \subset \ldots \subset
W_{p, m} = U_p
$$
であり、組
$$
(W_{p, 0} \subset W_{p, 1}, f_p|_{V_{p, 1}}),
(W_{p, 1} \subset W_{p, 2}, f_p|_{V_{p, 2}}),\ldots,
(W_{p, m - 1} \subset W_{p, m}, f_p|_{V_{p, m}})
$$
は Lemma \ref{spaces-perfect-lemma-make-more-elementary-distinguished-squares} により
初等識別正方形である。$P$ は各 $V_{p, 1}$ について成り立つ
（これらはアフィン概型だからである）。% P09-ERR-0981
また、$W_{p, i} \times_{W_{p, i + 1}} V_{p, i + 1}$ は
$V_{p, i + 1}$ の準コンパクト開部分であるため、証明の第一段落により
$P$ はこれについても成り立つ。したがって (2) をそれぞれに適用し、
帰納的に $P$ は $W_{p, 1}, \ldots, W_{p, m} = U_p$ について成り立つ。
\end{proof}

\begin{lemma}
\label{spaces-perfect-lemma-induction-principle-separated}
$S$ を概型とする。$X$ を $S$ 上の準コンパクトかつ準分離的な
代数空間とする。
$\mathcal{B} \subset \Ob(X_{spaces, \etale})$ とし、
$P$ を $\mathcal{B}$ の元の性質とする。次を仮定する。
\begin{enumerate}
\item 各 $W \in \mathcal{B}$ は準コンパクトかつ準分離的である。
\item $W \in \mathcal{B}$ かつ $U \subset W$ が準コンパクト開部分ならば、
$U \in \mathcal{B}$ である。
\item $V \in \Ob(X_{spaces, \etale})$ がアフィンならば、
(a) $V \in \mathcal{B}$ であり、(b) $P$ は $V$ について成り立つ。
\item 各初等識別正方形 $(U \subset W, f : V \to W)$ について、
次が成り立つとする。
\begin{enumerate}
\item $W \in \mathcal{B}$ である。
\item $U$ は準コンパクトである。
\item $V$ はアフィンである。
\item $P$ は $U$、$V$、および $U \times_W V$ について成り立つ。
\end{enumerate}
このとき $P$ は $W$ について成り立つ。
\end{enumerate}
このとき $P$ は各 $W \in \mathcal{B}$ について成り立つ。
\end{lemma}

\begin{proof}
Lemma \ref{spaces-perfect-lemma-induction-principle} の証明とまったく同じ方法で証明される。
（(4)(d) が意味をもつことにも注意する。実際、
$U \times_W V$ は $V$ の準コンパクト開部分であり、したがって条件
(2) と (3) により $\mathcal{B}$ の元である。）
\end{proof}

\begin{remark}
\label{spaces-perfect-remark-how-to}
Lemma \ref{spaces-perfect-lemma-induction-principle-separated} で集合
$\mathcal{B}$ をどのように選べばよいだろうか。
いくつか例を挙げる。
\begin{enumerate}
\item $X$ が準コンパクトかつ分離的ならば、
$\mathcal{B}$ として $X_{spaces, \etale}$ の準コンパクトかつ
分離的な対象の集合を選べる。このとき $X \in \mathcal{B}$ であり、
$\mathcal{B}$ は (1)、(2)、(3)(a) を満たす。この $\mathcal{B}$ の選択では、
Lemma \ref{spaces-perfect-lemma-induction-principle-separated} は
Lemma \ref{spaces-perfect-lemma-induction-principle} を再現する。
\item $X$ が準コンパクトで、$\mathbf{Z}$ 上アフィンな対角射をもつならば
（Properties of Spaces, Definition
\ref{spaces-properties-definition-separated} の意味で）、
$\mathcal{B}$ として $X_{spaces, \etale}$ の対象のうち、
準コンパクトで $\mathbf{Z}$ 上アフィンな対角射をもつものの集合を選べる。
再び $X \in \mathcal{B}$ であり、$\mathcal{B}$ は
(1)、(2)、(3)(a) を満たす。
\item $X$ が準コンパクトかつ準分離的ならば、(1)、(2)、(3)(a) を満たす
最小の部分集合 $\mathcal{B}$（$X$ を含むもの）は、
$W \in \mathcal{B}$ であることと、$W$ が $X$ の準コンパクト
開部分空間であるか、または $W$ が $X_{spaces, \etale}$ の
アフィン対象の準コンパクト開部分であることが同値になるように定められる。
\end{enumerate}
\end{remark}

\noindent
次は、ある開部分での成立を、より大きな開部分へ拡張する変形である。

\begin{lemma}
\label{spaces-perfect-lemma-induction-principle-enlarge}
$S$ を概型とする。$X$ を $S$ 上の準コンパクトかつ準分離的な
代数空間とする。$W \subset X$ を準コンパクト開部分空間とする。
$P$ を $X$ の準コンパクト開部分空間の性質とする。次を仮定する。
\begin{enumerate}
\item $P$ は $W$ について成り立つ。
\item 次の条件を満たす各初等識別正方形
$(W_1 \subset W_2, f : V \to W_2)$ について、
\begin{enumerate}
\item $W_1$、$W_2$ は $X$ の準コンパクト開部分空間である。
\item $W \subset W_1$ である。
\item $V$ はアフィンである。
\item $P$ は $W_1$ について成り立つ。
\end{enumerate}
このとき $P$ は $W_2$ について成り立つ。
\end{enumerate}
このとき $P$ は $X$ について成り立つ。
\end{lemma}

\begin{proof}
Lemma \ref{spaces-perfect-lemma-induction-principle-separated} から導くこともできるが、
ここでは Lemma \ref{spaces-perfect-lemma-induction-principle} の証明を明示的に
繰り返して直接示す。Decent Spaces, Lemma
\ref{decent-spaces-lemma-filter-quasi-compact-quasi-separated} の
フィルトレーション
$$
\emptyset = U_{n + 1} \subset
U_n \subset U_{n - 1} \subset \ldots \subset U_1 = X
$$
および射 $f_p : V_p \to U_p$ を用いる。
$P$ が $W_p = W \cup U_p$ について成り立つことを、
$p$ に関する降下帰納法によって示す。
$W_1 = X$ なので、これで証明が完了する。
(1) により $P$ は
$W_{n + 1} = W \cap U_{n + 1} = W$ について成り立つことに注意する。
% P09-ERR-0982
$P$ が $W_{p + 1}$ について成り立つと仮定する。
$W_p \setminus W_{p + 1}$（被約誘導部分空間構造を入れる）は
$U_p \setminus U_{p + 1}$ の閉部分空間である。
$(U_{p + 1} \subset U_p, f_p : V_p \to U_p)$ は初等識別正方形だから、
$(W_{p + 1} \subset W_p, f_p : V_p \to W_p)$ についても同様である。
しかし $V_p$ はアフィンとは限らないため、(2) を適用できないことがある。
一方、$V_p$ は準コンパクト概型なので、有限アフィン開被覆
$V_p = V_{p, 1} \cup \ldots \cup V_{p, m}$ を選べる。
$W_{p, 0} = W_{p + 1}$ と置き、
$$
W_{p, i} = W_{p + 1} \cup f_p(V_{p, 1} \cup \ldots \cup V_{p, i})
$$
と定める（$i = 1, \ldots, m$）。これらは $X$ の準コンパクト
開部分空間であり、$W$ を含む。このとき
$$
W_{p + 1} = W_{p, 0} \subset
W_{p, 1} \subset \ldots \subset
W_{p, m} = W_p
$$
であり、組
$$
(W_{p, 0} \subset W_{p, 1}, f_p|_{V_{p, 1}}),
(W_{p, 1} \subset W_{p, 2}, f_p|_{V_{p, 2}}),\ldots,
(W_{p, m - 1} \subset W_{p, m}, f_p|_{V_{p, m}})
$$
は Lemma \ref{spaces-perfect-lemma-make-more-elementary-distinguished-squares} により
初等識別正方形である。ここで (2) をそれぞれに適用し、
帰納的に $P$ は $W_{p, 1}, \ldots, W_{p, m} = W_p$ について成り立つ。
\end{proof}



\section{Mayer--Vietoris}
\label{spaces-perfect-section-mayer-vietoris}

\noindent
本節では、初等識別三角形からさまざまな Mayer--Vietoris 列が
得られることを証明する。% P09-ERR-0983

\medskip\noindent
$S$ を概型とする。$U \to X$ を、$S$ 上の代数空間の平展射とする。
Properties of Spaces, Section
\ref{spaces-properties-section-localize} では、構造層と両立する等式
$U_{spaces, \etale} = X_{spaces, \etale}/U$
が示された。したがって、この状況では射
$j_U : U \to X$ を局所化射とみなすことが多い
（Modules on Sites, Definition
\ref{sites-modules-definition-localize-ringed-site} を参照）。
特に、引き戻し $j_U^*$ を $U$ への制限とみなし、
しばしば $ {}|_U$ と表す。これは Properties of Spaces, Equation
(\ref{spaces-properties-equation-restrict-modules}) と両立する。
とりわけ、
\begin{equation}
\label{spaces-perfect-equation-stalk-restriction}
(\mathcal{F}|_U)_{\overline{u}} = \mathcal{F}_{\overline{x}}
\end{equation}
が成り立つ。ここで $\overline{u}$ は $U$ の幾何点であり、
$\overline{x}$ は $\overline{u}$ の $X$ における像である。
さらに、制限函手は完全な左随伴 $j_{U!}$ をもつ。
Modules on Sites, Lemmas
\ref{sites-modules-lemma-extension-by-zero} および
\ref{sites-modules-lemma-extension-by-zero-exact} を参照されたい。
最後に、$\mathcal{G}$ が $\mathcal{O}_X$-加群ならば、
\begin{equation}
\label{spaces-perfect-equation-stalk-j-shriek}
(j_{U!}\mathcal{G})_{\overline{x}} =
\bigoplus\nolimits_{\overline{u}} \mathcal{G}_{\overline{u}}
\end{equation}
であることを思い出す。これは任意の幾何点
$\overline{x} : \Spec(k) \to X$ に対して成り立ち、直和は射
$\overline{u} : \Spec(k) \to U$ のうち、
$j_U \circ \overline{u} = \overline{x}$ を満たすもの全体にわたる。
Modules on Sites, Lemma
\ref{sites-modules-lemma-stalk-j-shriek} および
Properties of Spaces, Lemma
\ref{spaces-properties-lemma-points-small-etale-site} を参照されたい。

\begin{lemma}
\label{spaces-perfect-lemma-exact-sequence-lower-shriek}
$S$ を概型とする。$(U \subset X, V \to X)$ を、
$S$ 上の代数空間の初等識別正方形とする。
\begin{enumerate}
\item $\mathcal{O}_X$-加群の層 $\mathcal{F}$ に対して、短完全列
$$
0 \to j_{U \times_X V!}\mathcal{F}|_{U \times_X V} \to
j_{U!}\mathcal{F}|_U \oplus j_{V!}\mathcal{F}|_V \to \mathcal{F} \to 0
$$
がある。
\item 対象 $E$（$D(\mathcal{O}_X)$ の対象）に対して、識別三角形
$$
j_{U \times_X V!}E|_{U \times_X V} \to
j_{U!}E|_U \oplus j_{V!}E|_V \to E \to 
j_{U \times_X V!}E|_{U \times_X V}[1]
$$
が $D(\mathcal{O}_X)$ において存在する。
\end{enumerate}
\end{lemma}

\begin{proof}
(1) の列が完全であることは、Properties of Spaces, Theorem
\ref{spaces-properties-theorem-exactness-stalks} により
幾何点での茎について確認すればよい。
$\overline{x}$ を $X$ の幾何点とする。Equations
(\ref{spaces-perfect-equation-stalk-restriction}) および
(\ref{spaces-perfect-equation-stalk-j-shriek}) によって
$\overline{x}$ で茎を取ると、列
$$
0 \to
\bigoplus\nolimits_{(\overline{u}, \overline{v})} \mathcal{F}_{\overline{x}}
\to
\bigoplus\nolimits_{\overline{u}} \mathcal{F}_{\overline{x}}
\oplus
\bigoplus\nolimits_{\overline{v}} \mathcal{F}_{\overline{x}}
\to
\mathcal{F}_{\overline{x}} \to 0
$$
を得る。この列は完全である。実際、各 $\overline{x}$ に対し、
ちょうど一つの $\overline{u}$ が存在して $\overline{x}$ へ写るか、
または $\overline{u}$ は存在せず、ちょうど一つの $\overline{v}$ が
$\overline{x}$ へ写る。

\medskip\noindent
(2) を証明する。Cohomology on Sites, Section
\ref{sites-cohomology-section-properties-K-injective} で見たように、
導来圏上の制限函手と零延長函手は、任意の複体へ函手をそのまま
適用することで計算される。$\mathcal{E}^\bullet$ を
$\mathcal{O}_X$-加群の複体で、$E$ を表すものとする。
補題の識別三角形は、Derived Categories, Section
\ref{derived-section-canonical-delta-functor}、特に Lemma
\ref{derived-lemma-derived-canonical-delta-functor} により、
$\mathcal{O}_X$-加群の複体の短完全列
$$
0 \to j_{U \times_X V!}\mathcal{E}^\bullet|_{U \times_X V} \to
j_{U!}\mathcal{E}^\bullet|_U \oplus j_{V!}\mathcal{E}^\bullet|_V
\to \mathcal{E}^\bullet \to 0
$$
に付随する識別三角形である。この列は (1) により短完全である。
\end{proof}

\begin{lemma}
\label{spaces-perfect-lemma-exact-sequence-j-star}
$S$ を概型とする。$(U \subset X, V \to X)$ を、
$S$ 上の代数空間の初等識別正方形とする。
\begin{enumerate}
\item 任意の $\mathcal{O}_X$-加群の層 $\mathcal{F}$ に対して、短完全列
$$
0 \to \mathcal{F} \to
j_{U, *}\mathcal{F}|_U \oplus j_{V, *}\mathcal{F}|_V \to
j_{U \times_X V, *}\mathcal{F}|_{U \times_X V} \to 0
$$
がある。
\item 任意の対象 $E$（$D(\mathcal{O}_X)$ の対象）に対して、識別三角形
$$
E \to 
Rj_{U, *}E|_U \oplus Rj_{V, *}E|_V \to
Rj_{U \times_X V, *}E|_{U \times_X V} \to
E[1]
$$
が $D(\mathcal{O}_X)$ において存在する。
\end{enumerate}
\end{lemma}

\begin{proof}
$W$ を $X_\etale$ の対象とする。列
$$
0 \to
\mathcal{F}(W) \to
\mathcal{F}(W \times_X U) \oplus \mathcal{F}(W \times_X V) \to
\mathcal{F}(W \times_X U \times_X V)
$$
は完全であり、最後の群の任意の元は $W$ 上局所的に中央の群へ
持ち上げられると主張する。
Lemma \ref{spaces-perfect-lemma-make-more-elementary-distinguished-squares} により、組
$(W \times_X U \subset W, V \times_X W \to W)$ は初等識別正方形である。
したがって $W = X$ と仮定してよく、次の列について同じことを
証明すれば十分である。
$$
0 \to
\mathcal{F}(X) \to
\mathcal{F}(U) \oplus \mathcal{F}(V) \to
\mathcal{F}(U \times_X V)
$$
すでに Lemma \ref{spaces-perfect-lemma-exact-sequence-lower-shriek} で、
$$
0 \to j_{U \times_X V!}\mathcal{O}_{U \times_X V}
\to j_{U!}\mathcal{O}_U \oplus
j_{V!}\mathcal{O}_V \to
\mathcal{O}_X \to 0
$$
が $\mathcal{O}_X$-加群の完全列であることを見た。
ここに右完全函手 $\Hom_{\mathcal{O}_X}(- , \mathcal{F})$ を適用すると、
上の列が得られる。% P09-ERR-0984
これはまた、$s \in \mathcal{F}(U \times_X V)$ を
$\mathcal{F}(U) \oplus \mathcal{F}(V)$ の元へ持ち上げる障害が
$$
\Ext^1_{\mathcal{O}_X}(\mathcal{O}_X, \mathcal{F}) =
H^1(X, \mathcal{F})
$$
に属することも意味する。コホモロジーの局所性
（Cohomology on Sites, Lemma
\ref{sites-cohomology-lemma-kill-cohomology-class-on-covering}）により、
この障害は $X$ 上平展局所的に消える。これで (1) の証明が完了する。

\medskip\noindent
(2) を証明する。K-入射的複体
$\mathcal{I}^\bullet$（$E$ を表すもの）で、各項 $\mathcal{I}^n$ が
$\textit{Mod}(\mathcal{O}_X)$ の入射的対象であるものを選ぶ。
Injectives, Theorem
\ref{injectives-theorem-K-injective-embedding-grothendieck} を参照されたい。
このとき $\mathcal{I}^\bullet|U$ は K-入射的複体である
（Cohomology on Sites, Lemma
\ref{sites-cohomology-lemma-restrict-K-injective-to-open}）。
したがって $Rj_{U, *}E|_U$ は
$j_{U, *}\mathcal{I}^\bullet|_U$ によって表される。
$V$ および $U \times_X V$ についても同様である。
ゆえに補題の識別三角形は、Derived Categories, Section
\ref{derived-section-canonical-delta-functor}、特に Lemma
\ref{derived-lemma-derived-canonical-delta-functor} により、
複体の短完全列
$$
0 \to
\mathcal{I}^\bullet \to
j_{U, *}\mathcal{I}^\bullet|_U \oplus j_{V, *}\mathcal{I}^\bullet|_V \to
j_{U \times_X V, *}\mathcal{I}^\bullet|_{U \times_X V} \to
0.
$$
に付随する識別三角形である。この列は (1) により完全である。
\end{proof}

\begin{lemma}
\label{spaces-perfect-lemma-unbounded-relative-mayer-vietoris}
$S$ を概型とする。$f : X \to Y$ を、$S$ 上の代数空間の射とする。
$(U \subset X, V \to X)$ を初等識別正方形とする。
制限を $a = f|_U : U \to Y$、$b = f|_V : V \to Y$、
$c = f|_{U \times_X V} : U \times_X V \to Y$ と表す。
任意の対象 $E$（$D(\mathcal{O}_X)$ の対象）に対して、識別三角形
$$
Rf_*E \to
Ra_*(E|_U) \oplus Rb_*(E|_V) \to
Rc_*(E|_{U \times_X V}) \to
Rf_*E[1]
$$
が $D(\mathcal{O}_Y)$ において存在する。この三角形は $E$ に関して
函手的である。
\end{lemma}

\begin{proof}
K-入射的複体 $\mathcal{I}^\bullet$（$E$ を表すもの）を選ぶ。
$\mathcal{I}^n$ が $\textit{Mod}(\mathcal{O}_X)$ の
入射的対象であると仮定してよい（すべての $n$ に対して）。
Injectives, Theorem
\ref{injectives-theorem-K-injective-embedding-grothendieck} を参照されたい。
このとき $Rf_*E$ は $f_*\mathcal{I}^\bullet$ で計算される。
Cohomology on Sites, Lemma
\ref{sites-cohomology-lemma-restrict-K-injective-to-open} により、
$U$、$V$、および $U \cap V$ についても同様である。
したがって補題の識別三角形は、Derived Categories, Section
\ref{derived-section-canonical-delta-functor}、特に Lemma
\ref{derived-lemma-derived-canonical-delta-functor} により、
複体の短完全列
$$
0 \to
f_*\mathcal{I}^\bullet \to
a_*\mathcal{I}^\bullet|_U \oplus b_*\mathcal{I}^\bullet|_V \to
c_*\mathcal{I}^\bullet|_{U \times_X V} \to
0.
$$
に付随する識別三角形である。これが複体の短完全列であることを、
次のように示す。入射的対象 $\mathcal{I}$
（$\textit{Mod}(\mathcal{O}_X)$ の対象）を選ぶ。$f_*$ を Lemma
\ref{spaces-perfect-lemma-exact-sequence-j-star} の短完全列
$$
0 \to \mathcal{I} \to
j_{U, *}\mathcal{I}|_U \oplus j_{V, *}\mathcal{I}|_V \to
j_{U \times_X V, *}\mathcal{I}|_{U \times_X V} \to 0
$$
に適用し、$R^1f_*\mathcal{I} = 0$ を用いると、短完全列
$$
0 \to f_*\mathcal{I} \to
f_*j_{U, *}\mathcal{I}|_U \oplus f_*j_{V, *}\mathcal{I}|_V \to
f_*j_{U \times_X V, *}\mathcal{I}|_{U \times_X V} \to 0
$$
を得る。$a_* = f_*j_{U, *}$ であり、$b_*$ と $c_*$ についても
同様であることに注意すれば、証明が完了する。
\end{proof}

\begin{lemma}
\label{spaces-perfect-lemma-mayer-vietoris-hom}
$S$ を概型とする。$(U \subset X, V \to X)$ を、
$S$ 上の代数空間の初等識別正方形とする。
対象 $E$、$F$（$D(\mathcal{O}_X)$ の対象）に対して、
Mayer--Vietoris 列
$$
\xymatrix{
& \ldots \ar[r] &
\Ext^{-1}(E_{U \times_X V}, F_{U \times_X V}) \ar[lld] \\
\Hom(E, F) \ar[r] &
\Hom(E_U, F_U) \oplus
\Hom(E_V, F_V) \ar[r] &
\Hom(E_{U \times_X V}, F_{U \times_X V})
}
$$
がある。下添字は該当する開部分への制限を表し、
$\Hom$ は該当する導来圏で取る。
\end{lemma}

\begin{proof}
Lemma \ref{spaces-perfect-lemma-exact-sequence-lower-shriek} の識別三角形を用い、
$\Hom$ の長完全列を得る
（Derived Categories, Lemma
\ref{derived-lemma-representable-homological}）。
さらに Cohomology on Sites, Lemma
\ref{sites-cohomology-lemma-adjoint-lower-shriek-restrict} による等式
$\Hom(j_{U!}E|_U, F) = \Hom(E|_U, F|_U)$ を用いる。
\end{proof}

\begin{lemma}
\label{spaces-perfect-lemma-unbounded-mayer-vietoris}
$S$ を概型とする。$(U \subset X, V \to X)$ を、
$S$ 上の代数空間の初等識別正方形とする。対象 $E$
（$D(\mathcal{O}_X)$ の対象）に対して、識別三角形
$$
R\Gamma(X, E) \to R\Gamma(U, E) \oplus R\Gamma(V, E) \to
R\Gamma(U \times_X V, E) \to R\Gamma(X, E)[1]
$$
があり、特にコホモロジーの長完全列
$$
\ldots \to
H^n(X, E) \to
H^n(U, E) \oplus H^n(V, E) \to
H^n(U \times_X V, E) \to
H^{n + 1}(X, E) \to \ldots
$$
がある。識別三角形と長完全列の構成は $E$ に関して函手的である。
\end{lemma}

\begin{proof}
K-入射的複体 $\mathcal{I}^\bullet$（$E$ を表すもの）で、
各項 $\mathcal{I}^n$ が $\textit{Mod}(\mathcal{O}_X)$ の
入射的対象であるものを選ぶ。Injectives, Theorem
\ref{injectives-theorem-K-injective-embedding-grothendieck} を参照されたい。
Lemma \ref{spaces-perfect-lemma-exact-sequence-j-star} の証明で、複体の短完全列
$$
0 \to \mathcal{I}^\bullet \to
j_{U, *}\mathcal{I}^\bullet|_U \oplus j_{V, *}\mathcal{I}^\bullet|_V \to
j_{U \times_X V, *}\mathcal{I}^\bullet|_{U \times_X V} \to 0
$$
を得た。$H^1(X, \mathcal{I}^n) = 0$ だから、大域切断を取ると
複体の完全列
$$
0 \to \Gamma(X, \mathcal{I}^\bullet) \to
\Gamma(U, \mathcal{I}^\bullet) \oplus
\Gamma(V, \mathcal{I}^\bullet) \to
\Gamma(U \times_X V, \mathcal{I}^\bullet) \to 0
$$
を得る。これらの複体は
$R\Gamma(X, E)$、$R\Gamma(U, E)$、$R\Gamma(V, E)$、および
$R\Gamma(U \times_X V, E)$ を表すので、Derived Categories, Section
\ref{derived-section-canonical-delta-functor}、特に Lemma
\ref{derived-lemma-derived-canonical-delta-functor} により
識別三角形を得る。
\end{proof}

\begin{lemma}
\label{spaces-perfect-lemma-restrict-lower-shriek}
$S$ を概型とする。$j : U \to X$ を、$S$ 上の代数空間の平展射とする。
平展射 $V \to Y$ が与えられたとし、$W = V \times_X U$ と置く。
$j_W : W \to V$ で射影を表す。このとき
$(j_!E)|_V = j_{W!}(E|_W)$ が、$E$
（$D(\mathcal{O}_U)$ に属するもの）に対して成り立つ。% P09-ERR-0985
\end{lemma}

\begin{proof}
$(j_!\mathcal{F})|_V = j_{W!}(\mathcal{F}|_W)$
は $\mathcal{O}_X$-加群 $\mathcal{F}$ に対して成り立つから、
この主張は従う。これは函手 $j_!$ と $j_{W!}$ の
構成から直ちに従う。Modules on Sites, Lemma
\ref{sites-modules-lemma-extension-by-zero} を参照されたい。
\end{proof}

\begin{lemma}
\label{spaces-perfect-lemma-pushforward-with-support-in-open}
$S$ を概型とする。$(U \subset X, j : V \to X)$ を、
$S$ 上の代数空間の初等識別正方形とする。
$T = |X| \setminus |U|$ と置く。
\begin{enumerate}
\item $E$ が $D(\mathcal{O}_X)$ の対象で、$T$ 上に台をもつならば、
(a) $E \to Rj_*(E|_V)$ と (b) $j_!(E|_V) \to E$ は同型である。
\item $F$ が $D(\mathcal{O}_V)$ の対象で、$j^{-1}T$ 上に台をもつならば、
(a) $F \to (j_!F)|_V$、(b) $(Rj_*F)|_V \to F$、および (c)
$j_!F \to Rj_*F$ は同型である。
\end{enumerate}
\end{lemma}

\begin{proof}
$E$ を $D(\mathcal{O}_X)$ の対象で、そのコホモロジー層が
$T$ 上に台をもつものとする。このとき $E|_U = 0$ および
$E|_{U \times_X V} = 0$ である。実際、$T$ は $U$ と交わらず、
$j^{-1}T$ は $U \times_X V$ と交わらない。
したがって (1)(a) は Lemma \ref{spaces-perfect-lemma-exact-sequence-j-star} から従う。
まったく同様に、(1)(b) は
Lemma \ref{spaces-perfect-lemma-exact-sequence-lower-shriek} から従う。

\medskip\noindent
$F$ を $D(\mathcal{O}_V)$ の対象で、そのコホモロジー層が
$j^{-1}T$ 上に台をもつものとする。
Lemma \ref{spaces-perfect-lemma-restrict-direct-image-open} により
$(Rj_*F)|_U = Rj_{W, *}(F|_W) = 0$ である。これは仮定から
$F|_W = 0$ だからである。同様に、
Lemma \ref{spaces-perfect-lemma-restrict-lower-shriek} により
$(j_!F)|_U = j_{W!}(F|_W) = 0$ である。
したがって $j_!F$ と $Rj_*F$ は $T$ 上に台をもち、
$(j_!F)|_V$ と $(Rj_*F)|_V$ は $j^{-1}(T)$ 上に台をもつ。
導来圏で射 (2)(a)、(b)、(c) が同型であることを確認するには、
Properties of Spaces, Theorem
\ref{spaces-properties-theorem-exactness-stalks} により、
$T$ および $j^{-1}(T)$ の幾何点でコホモロジー層の茎に誘導される射が
同型であることを確認すれば十分である。
Lemmas \ref{spaces-perfect-lemma-restrict-direct-image-open}、
\ref{spaces-perfect-lemma-restrict-lower-shriek}、および
\ref{spaces-perfect-lemma-make-more-elementary-distinguished-squares} によれば、
$X$ を $V$ で、$U$ を $U \times_X V$ で、
$V$ を $V \times_X V$ で、$F$ を
$F|_{V \times_X V}$（第一射影による制限）で置き換えた後に
これを確認してよい。$V \times_X V \to V$ は切断をもつので、
(2) は $j : V \to X$ が切断をもつ場合へ帰着する。

\medskip\noindent
$j$ が切断 $\sigma : X \to V$ をもつと仮定する。
$V' = \sigma(X)$ と置く。これは $V$ の開部分空間である。
$U' = j^{-1}(U)$ と置く。これは $V$ のもう一つの開部分空間である。
このとき $(U' \subset V, V' \to V)$ は初等識別正方形である。
$F|_{U'} = 0$ および $F|_{V' \cap U'} = 0$ であることに注意する。
これは $F$ が $j^{-1}(T)$ 上に台をもつからである。
$j' : V' \to V$ で開埋め込みを表し、
$j_{V'} : V' \to X$ で合成 $V' \to V \to X$ を表す。
後者は $\sigma$ の逆である。$F' = \sigma^*F$ と置く。
Lemmas \ref{spaces-perfect-lemma-exact-sequence-lower-shriek} および
\ref{spaces-perfect-lemma-exact-sequence-j-star} の識別三角形から
$F = j'_!(F|_{V'})$ および $F = Rj'_*(F|_{V'})$ が従う。
したがって $j_!F = j_!j'_!(F|_{V'}) = j_{V'!}F = F'$ である。
実際、$j_{V'} : V' \to X$ は同型であり、$\sigma$ の逆である。
同様に $Rj_*F = Rj_*Rj'_*F = Rj_{V', *}F = F'$ である。
これで (2)(c) が示された。(2)(a) と (2)(b) を示すには、
$F = F'|_V$ を示せば十分である。これは明らかである。実際、
$F$ と $F'|_V$ はともに $U'$ と $U' \cap V'$ への制限が零であり、
$V'$ への制限が同じ対象である。
\end{proof}

\noindent
複体を貼り合わせることができる！

\begin{lemma}
\label{spaces-perfect-lemma-glue}
$S$ を概型とする。$(U \subset X, V \to X)$ を、
$S$ 上の代数空間の初等識別正方形とする。次が与えられたとする。
\begin{enumerate}
\item 対象 $A$（$D(\mathcal{O}_U)$ の対象）、
\item 対象 $B$（$D(\mathcal{O}_V)$ の対象）、
\item 同型 $c : A|_{U \times_X V} \to B|_{U \times_X V}$。
\end{enumerate}
このとき対象 $F$（$D(\mathcal{O}_X)$ の対象）と同型
$f : F|_U \to A$、$g : F|_V \to B$ が存在し、
$c = g|_{U \times_X V} \circ f^{-1}|_{U \times_X V}$ を満たす。
さらに、次が与えられたとする。
\begin{enumerate}
\item 対象 $E$（$D(\mathcal{O}_X)$ の対象）、
\item 射 $a : A \to E|_U$（$D(\mathcal{O}_U)$ の射）、
\item 射 $b : B \to E|_V$（$D(\mathcal{O}_V)$ の射）。
\end{enumerate}
これらが
$$
a|_{U \times_X V}  = b|_{U \times_X V} \circ c.
$$
を満たすとする。このとき射 $F \to E$ が $D(\mathcal{O}_X)$ に存在し、
$U$ への制限は $a \circ f$、$V$ への制限は $b \circ g$ である。
\end{lemma}

\begin{proof}
$j_U$、$j_V$、$j_{U \times_X V}$ で、それぞれの $X$ への射を表す。
識別三角形
$$
F \to Rj_{U, *}A \oplus Rj_{V, *}B \to
Rj_{U \times_X V, *}(B|_{U \times_X V}) \to F[1]
$$
を選ぶ。ここで射
$Rj_{V, *}B \to Rj_{U \times_X V, *}(B|_{U \times_X V})$
は明らかなものである。射
$Rj_{U, *}A \to Rj_{U \times_X V, *}(B|_{U \times_X V})$
は
$Rj_{U, *}A \to Rj_{U \times_X V, *}(A|_{U \times_X V})$
と $Rj_{U \times_X V, *}c$ の合成である。$U$ へ制限すると、
$$
F|_U \to A \oplus (Rj_{V, *}B)|_U \to
(Rj_{U \times_X V, *}(B|_{U \times_X V}))|_U \to F|_U[1]
$$
を得る。$j : U \times_X V \to U$ と表す。
制限と全直像の両立性（Lemma
\ref{spaces-perfect-lemma-restrict-direct-image-open}）により、
$(Rj_{V, *}B)|_U$ と
$(Rj_{U \times_X V, *}(B|_{U \times_X V}))|_U$
はいずれも $Rj_*(B|_{U \times_X V})$ と標準的に同型である。
したがって最後に表示した式の第二の射は切断をもち、
射 $F|_U \to A$ は同型である。

\medskip\noindent
射 $F|_V \to B$ が同型であることを示すため、次の技巧を用いる。
識別三角形
$$
F|_V \to B \to B' \to F[1]|_V
$$
を $D(\mathcal{O}_V)$ において選ぶ。$F|_U \to A$ は同型であり、かつ同型
$c : A|_{U \times_X V} \to B|_{U \times_X V}$ があるので、
$F|_V \to B$ の $U \times_X V$ 上への制限は同型である。
したがって $B'$ は $j_V^{-1}(T)$ 上に台をもつ。ここで
$T = |X| \setminus |U|$ である。一方、識別三角形の射
$$
\xymatrix{
F \ar[r] \ar[d] &
Rj_{U, *}F|_U \oplus Rj_{V, *}F|_V \ar[r] \ar[d] &
Rj_{U \times_X V, *}F|_{U \times_X V} \ar[r] \ar[d] &
F[1] \ar[d] \\
F \ar[r] &
Rj_{U, *}A \oplus Rj_{V, *}B \ar[r] &
Rj_{U \times_X V, *}(B|_{U \times_X V}) \ar[r] &
F[1]
}
$$
がある。この図式の垂直射は
$Rj_{V, *}F|_V \to Rj_{V, *}B$ を除いてすべて同型である。
したがって、この射も同型である
（Derived Categories, Lemma
\ref{derived-lemma-third-isomorphism-triangle}）。
これは $Rj_{V, *}B' = 0$ を含意する。ゆえに Lemma
\ref{spaces-perfect-lemma-pushforward-with-support-in-open} により $B' = 0$ である。

\medskip\noindent
射 $F \to E$ の存在は、$\Hom$ に対する Mayer--Vietoris 列、
すなわち Lemma \ref{spaces-perfect-lemma-mayer-vietoris-hom} から従う。
\end{proof}







\section{コヒーレータ}
\label{spaces-perfect-section-coherator}

\noindent
$S$ を概型とする。$X$ を $S$ 上の代数空間とする。
{\it コヒーレータ}とは函手
$$
Q_X :
\textit{Mod}(\mathcal{O}_X)
\longrightarrow
\QCoh(\mathcal{O}_X)
$$
で、包含函手
$\QCoh(\mathcal{O}_X) \to \textit{Mod}(\mathcal{O}_X)$
の右随伴であるものをいう。これは任意の代数空間 $X$ に対して存在し、
さらに随伴射 $Q_X(\mathcal{F}) \to \mathcal{F}$ は、
任意の準連接加群 $\mathcal{F}$ に対して同型である。
Properties of Spaces, Proposition
\ref{spaces-properties-proposition-coherator} を参照されたい。
$Q_X$ は（右随伴なので）左完全であるから、その右導来延長
$$
RQ_X :
D(\mathcal{O}_X)
\longrightarrow
D(\QCoh(\mathcal{O}_X)).
$$
を考えることができる。$Q_X$ は包含函手
$\QCoh(\mathcal{O}_X) \to \textit{Mod}(\mathcal{O}_X)$
の右随伴なので、Derived Categories, Lemma
\ref{derived-lemma-derived-adjoint-functors} により、$RQ_X$ は標準函手
$D(\QCoh(\mathcal{O}_X)) \to D(\mathcal{O}_X)$ の右随伴である。

\medskip\noindent
本節では函手 $RQ_X$ を調べる。Section
\ref{spaces-perfect-section-better-coherator} では、包含函手
$D_\QCoh(\mathcal{O}_X) \to D(\mathcal{O}_X)$ の右随伴
（存在する場合）を調べる。これは上の函手と密接に関連している。

\begin{lemma}
\label{spaces-perfect-lemma-affine-pushforward}
$S$ を概型とする。$f : X \to Y$ を、$S$ 上の代数空間の
アフィン射とする。このとき $f_*$ は導来函手
$f_* : D(\QCoh(\mathcal{O}_X)) \to D(\QCoh(\mathcal{O}_Y))$
を定める。この函手は、図式
$$
\xymatrix{
D(\QCoh(\mathcal{O}_X)) \ar[d]_{f_*} \ar[r] &
D_\QCoh(\mathcal{O}_X) \ar[d]^{Rf_*} \\
D(\QCoh(\mathcal{O}_Y)) \ar[r] &
D_\QCoh(\mathcal{O}_Y)
}
$$
が可換であるという性質をもつ。
\end{lemma}

\begin{proof}
函手
$f_* : \QCoh(\mathcal{O}_X) \to \QCoh(\mathcal{O}_Y)$
は完全である。Cohomology of Spaces, Lemma
\ref{spaces-cohomology-lemma-affine-vanishing-higher-direct-images}
を参照されたい。したがって $f_*$ は導来函手
$f_* : D(\QCoh(\mathcal{O}_X)) \to D(\QCoh(\mathcal{O}_Y))$
を定める。これは任意の代表複体へ $f_*$ をそのまま適用して得られる。
Derived Categories, Lemma
\ref{derived-lemma-right-derived-exact-functor} を参照されたい。
任意の $\mathcal{O}_X$-加群の複体 $\mathcal{F}^\bullet$ に対して、
標準的な射
$f_*\mathcal{F}^\bullet \to Rf_*\mathcal{F}^\bullet$
がある。証明を終えるには、$\mathcal{F}^\bullet$ が各
$\mathcal{F}^n$ を準連接とする複体ならば、これが準同型であることを
示せばよい。この主張は $Y$ 上平展局所的なので、$Y$ は
アフィンであると仮定してよい。アフィン射は表現可能だから、
Remark \ref{spaces-perfect-remark-match-total-direct-images} の両立性により
概型の場合へ帰着する。概型の場合は Derived Categories of Schemes,
Lemma \ref{perfect-lemma-affine-pushforward} である。
\end{proof}

\begin{lemma}
\label{spaces-perfect-lemma-flat-pushforward-coherator}
$S$ を概型とする。$f : X \to Y$ を、$S$ 上の代数空間の射とする。
$f$ は準コンパクト、準分離的、かつ平坦であると仮定する。
函手
$$
\Phi : D(\QCoh(\mathcal{O}_X)) \to D(\QCoh(\mathcal{O}_Y))
$$
で、左完全函手
$f_* : \QCoh(\mathcal{O}_X) \to \QCoh(\mathcal{O}_Y)$
の右導来函手を表す。このとき
$RQ_Y \circ Rf_* = \Phi \circ RQ_X$ である。
\end{lemma}

\begin{proof}
$RQ_Y \circ Rf_*$ と $\Phi \circ RQ_X$ が同じ函手
$D(\QCoh(\mathcal{O}_Y)) \to D(\mathcal{O}_X)$
の右随伴であることを示して証明する。

\medskip\noindent
$f$ は準コンパクトかつ準分離的なので、$f_*$ は準連接性を保つ。
Morphisms of Spaces, Lemma
\ref{spaces-morphisms-lemma-pushforward} を参照されたい。
$\QCoh(\mathcal{O}_X)$ は Grothendieck アーベル圏であることを思い出す
（Properties of Spaces, Proposition
\ref{spaces-properties-proposition-coherator}）。
したがって、任意の $K$（$D(\QCoh(\mathcal{O}_X))$ に属するもの）は、
K-入射的複体 $\mathcal{I}^\bullet$（$\QCoh(\mathcal{O}_X)$ の複体）
によって表すことができる。
Injectives, Theorem
\ref{injectives-theorem-K-injective-embedding-grothendieck} を参照されたい。
このとき $\Phi(K) = f_*\mathcal{I}^\bullet$ と定義できる。

\medskip\noindent
$f$ は平坦だから、函手 $f^*$ は完全である。したがって $f^*$ は
$f^* : D(\mathcal{O}_Y) \to D(\mathcal{O}_X)$ および
$f^* : D(\QCoh(\mathcal{O}_Y)) \to D(\QCoh(\mathcal{O}_X))$
を定める。函手
$f^* = Lf^* : D(\mathcal{O}_Y) \to D(\mathcal{O}_X)$
は
$Rf_* : D(\mathcal{O}_X) \to D(\mathcal{O}_Y)$
の左随伴である。Cohomology on Sites, Lemma
\ref{sites-cohomology-lemma-adjoint} を参照されたい。同様に、函手
$f^* : D(\QCoh(\mathcal{O}_Y)) \to D(\QCoh(\mathcal{O}_X))$
は、Derived Categories, Lemma
\ref{derived-lemma-derived-adjoint-functors} により、
$\Phi : D(\QCoh(\mathcal{O}_X)) \to D(\QCoh(\mathcal{O}_Y))$
の左随伴である。

\medskip\noindent
$A$ を $D(\QCoh(\mathcal{O}_Y))$ の対象とし、
$E$ を $D(\mathcal{O}_X)$ の対象とする。このとき
\begin{align*}
\Hom_{D(\QCoh(\mathcal{O}_Y))}(A, RQ_Y(Rf_*E))
& =
\Hom_{D(\mathcal{O}_Y)}(A, Rf_*E) \\
& =
\Hom_{D(\mathcal{O}_X)}(f^*A, E) \\
& =
\Hom_{D(\QCoh(\mathcal{O}_X))}(f^*A, RQ_X(E)) \\
& =
\Hom_{D(\QCoh(\mathcal{O}_Y))}(A, \Phi(RQ_X(E)))
\end{align*}
である。これは望む結論を含意する。
\end{proof}

\begin{lemma}
\label{spaces-perfect-lemma-affine-coherator}
$S$ を概型とする。$X$ を $S$ 上のアフィン代数空間とする。
$A = \Gamma(X, \mathcal{O}_X)$ と置く。このとき次が成り立つ。
\begin{enumerate}
\item $Q_X : \textit{Mod}(\mathcal{O}_X) \to \QCoh(\mathcal{O}_X)$ は、
$\mathcal{F}$ を準連接 $\mathcal{O}_X$-加群で、
$A$-加群 $\Gamma(X, \mathcal{F})$ に付随するものへ送る函手である。
\item $RQ_X : D(\mathcal{O}_X) \to D(\QCoh(\mathcal{O}_X))$ は、
$E$ を準連接 $\mathcal{O}_X$-加群の複体で、
$R\Gamma(X, E)$（$D(A)$ の対象）に付随するものへ送る函手である。
\item $D_\QCoh(\mathcal{O}_X)$ に制限した函手 $RQ_X$ は、
(\ref{spaces-perfect-equation-compare}) の準逆を定める。
\end{enumerate}
\end{lemma}

\begin{proof}
$X_0 = \Spec(A)$ を $X$ を表現するアフィン概型とする。
環付きサイトの射
$\epsilon : X_\etale \to X_{0, Zar}$
が存在し、圏同値
$$
\xymatrix{
\QCoh(\mathcal{O}_X) \ar@<1ex>[r]^{{\epsilon_*}} &
\QCoh(\mathcal{O}_{X_0}) \ar@<1ex>[l]^{{\epsilon^*}}
}
$$
を誘導することを思い出す。Lemma
\ref{spaces-perfect-lemma-derived-quasi-coherent-small-etale-site} を参照されたい。
したがって、随伴函手の一意性により
$Q_X = \epsilon^* \circ Q_{X_0} \circ \epsilon_*$ である。
よって (1) は Derived Categories of Schemes, Lemma
\ref{perfect-lemma-affine-coherator} における $Q_{X_0}$ の記述と、等式
$\Gamma(X_0, \epsilon_*\mathcal{F}) = \Gamma(X, \mathcal{F})$
から従う。(2) は (1) と、$A$-加群から準連接
$\mathcal{O}_X$-加群への函手が完全であることから従う。
第三の主張は、概型についての結果
（Derived Categories of Schemes, Lemma
\ref{perfect-lemma-affine-coherator}）と Lemma
\ref{spaces-perfect-lemma-derived-quasi-coherent-small-etale-site} から従う。
\end{proof}

\noindent
次に、函手
$D(\QCoh(\mathcal{O}_X)) \to D_\QCoh(\mathcal{O}_X)$
が圏同値となるための判定条件を証明する。

\begin{lemma}
\label{spaces-perfect-lemma-argument-proves}
$S$ を概型とする。$X$ を $S$ 上の準コンパクトかつ準分離的な
代数空間とする。任意の平展射 $j : V \to W$ で、
$W \subset X$ が準コンパクト開部分、$V$ がアフィンであるものについて、
右導来函手
$$
\Phi : D(\QCoh(\mathcal{O}_U)) \to D(\QCoh(\mathcal{O}_W))
$$
（左完全函手
$j_* : \QCoh(\mathcal{O}_V) \to \QCoh(\mathcal{O}_W)$
の右導来函手）が
が可換図式
$$
\xymatrix{
D(\QCoh(\mathcal{O}_V)) \ar[d]_\Phi \ar[r]_{i_V} &
D_\QCoh(\mathcal{O}_V) \ar[d]^{Rj_*} \\
D(\QCoh(\mathcal{O}_W)) \ar[r]^{i_W} &
D_\QCoh(\mathcal{O}_W)
}
$$
に入ると仮定する。% P09-ERR-0986
このとき函手 (\ref{spaces-perfect-equation-compare})
$$
D(\QCoh(\mathcal{O}_X))
\longrightarrow
D_\QCoh(\mathcal{O}_X)
$$
は圏同値であり、その準逆は $RQ_X$ によって与えられる。
\end{lemma}

\begin{proof}
まず帰納原理を用いて $i_X$ が充満忠実であることを証明する。
より正確には、Lemma \ref{spaces-perfect-lemma-induction-principle-enlarge} を用いる。
$(U \subset W, V \to W)$ を初等識別正方形とし、$V$ はアフィン、
$U, W$ は $X$ の準コンパクト開部分であるとする。
$i_U$ が充満忠実であると仮定し、$i_W$ が充満忠実であることを示す。
$X$ を $W$ で置き換えてよい。すなわち $W = X$ と仮定してよい
（これは記法を簡単にするためだけである。補題の仮定がこの操作で
保たれることに注意する）。

\medskip\noindent
$A, B$ を $D(\QCoh(\mathcal{O}_X))$ の対象とする。写像
$$
\Hom_{D(\QCoh(\mathcal{O}_X))}(A, B)
\longrightarrow
\Hom_{D(\mathcal{O}_X)}(i_X(A), i_X(B))
$$
が全単射であることを示したい。$T = |X| \setminus |U|$ と置く。

\medskip\noindent
まず $i_X(B)$ が $T$ 上に台をもつと仮定する。この場合、射
$$
i_X(B) \to Rj_{V, *}(i_X(B)|_V) = Rj_{V, *}(i_V(B|_V))
$$
は準同型である
（Lemma \ref{spaces-perfect-lemma-pushforward-with-support-in-open}）。
仮定により、同型
$i_X(\Phi(B|_V)) \to Rj_{V, *}(i_V(B|_V))$ が
$D(\mathcal{O}_X)$ に存在する。
さらに、準連接加群の導来圏上で $\Phi$ と $ {-}|_V$ は随伴函手である
（Derived Categories, Lemma
\ref{derived-lemma-derived-adjoint-functors}）。
随伴射 $B \to \Phi(B|_V)$ は、$i_X$ を適用すると同型となるので、
$D(\QCoh(\mathcal{O}_X))$ において同型である。したがって
\begin{align*}
\Mor_{D(\QCoh(\mathcal{O}_X))}(A, B)
& =
\Mor_{D(\QCoh(\mathcal{O}_X))}(A, \Phi(B|_V)) \\
& =
\Mor_{D(\QCoh(\mathcal{O}_V))}(A|_V, B|_V) \\
& =
\Mor_{D(\mathcal{O}_V)}(i_V(A|_V), i_V(B|_V)) \\
& =
\Mor_{D(\mathcal{O}_X)}(i_X(A), Rj_{V, *}(i_V(B|_V))) \\
& =
\Mor_{D(\mathcal{O}_X)}(i_X(A), i_X(B))
\end{align*}
となり、望む結論を得る。ここで $i_V$ が充満忠実であること
（Lemma \ref{spaces-perfect-lemma-affine-coherator}）を用いた。

\medskip\noindent
一般の場合を考える。任意の複体 $\mathcal{B}^\bullet$
（準連接 $\mathcal{O}_X$-加群の複体で $B$ を表すもの）を選ぶ。
次に、任意の準同型
$s : \mathcal{B}^\bullet|_U \to \mathcal{C}^\bullet$
（$U$ 上の準連接加群の複体の準同型）を選ぶ。
$j_U : U \to X$ は準コンパクトかつ準分離的なので、函手
$j_{U, *}$ は準連接加群を準連接加群へ移す
（Morphisms of Spaces, Lemma
\ref{spaces-morphisms-lemma-pushforward}）。
したがって、標準的な射
$\mathcal{B}^\bullet \to j_{U, *}(\mathcal{B}^\bullet|_U) \to
j_{U, *}\mathcal{C}^\bullet$
（$X$ 上の準連接加群の複体の射）がある。
$B'' = j_{U, *}\mathcal{C}^\bullet$ と
$D(\QCoh(\mathcal{O}_X))$ において置き、識別三角形
$$
B \to B'' \to B' \to B[1]
$$
を $D(\QCoh(\mathcal{O}_X))$ において選ぶ。
この三角形の第一の射は $U$ 上への制限が同型なので、
$B'$ は $T$ 上に台をもつ。したがって図式
$$
\xymatrix{
\Hom_{D(\QCoh(\mathcal{O}_X))}(A, B'[-1]) \ar[r] \ar[d] &
\Hom_{D(\mathcal{O}_X)}(i_X(A), i_X(B')[-1]) \ar[d] \\
\Hom_{D(\QCoh(\mathcal{O}_X))}(A, B) \ar[r] \ar[d] &
\Hom_{D(\mathcal{O}_X)}(i_X(A), i_X(B)) \ar[d] \\
\Hom_{D(\QCoh(\mathcal{O}_X))}(A, B'') \ar[r] \ar[d] &
\Hom_{D(\mathcal{O}_X)}(i_X(A), i_X(B'')) \ar[d] \\
\Hom_{D(\QCoh(\mathcal{O}_X))}(A, B') \ar[r] &
\Hom_{D(\mathcal{O}_X)}(i_X(A), i_X(B'))
}
$$
では、列は完全であり、最上段と最下段の水平射は全単射である。
最後に、準連接加群の複体 $\mathcal{A}^\bullet$（$A$ を表すもの）
を選ぶ。

\medskip\noindent
$\alpha : i_X(A) \to i_X(B)$ を $D(\mathcal{O}_X)$ における射とする。
制限 $\alpha|_U$ は $D(\QCoh(\mathcal{O}_U))$ の射から生じる。
これは $i_U$ が充満忠実だからである。したがって、上の
$s : \mathcal{B}^\bullet|_U \to \mathcal{C}^\bullet$ を、
$\alpha|_U$ が実際の複体の射
$\mathcal{A}^\bullet|_U \to \mathcal{C}^\bullet$
によって表されるように選べる。これは複体の射
$\mathcal{A} \to j_{U, *}\mathcal{C}^\bullet$ に対応する。
% P09-ERR-0987
言い換えると、$\alpha$ の
$\Hom_{D(\mathcal{O}_X)}(i_X(A), i_X(B''))$ における像は、
$\Hom_{D(\QCoh(\mathcal{O}_X))}(A, B'')$ の元から生じる。
図式追跡により、$\alpha$ は射 $A \to B$
（$D(\QCoh(\mathcal{O}_X))$ の射）から生じる。

最後に、$a : A \to B$ を $D(\QCoh(\mathcal{O}_X))$ の射で、
$D(\mathcal{O}_X)$ において零となるものとする。
$\mathcal{B}^\bullet$ を適切に選べば、$a$ は複体の射
$a^\bullet : \mathcal{A}^\bullet \to \mathcal{B}^\bullet$
によって表されると仮定してよい。$i_U$ は充満忠実なので、
制限 $a^\bullet|_U$ は $D(\QCoh(\mathcal{O}_U))$ において零である。
% P09-ERR-0988
したがって、$s$ を、合成
$s \circ a^\bullet|_U : \mathcal{A}^\bullet|_U \to \mathcal{C}^\bullet$
が零と鎖ホモトピックになるように選べる。
函手 $j_{U, *}$ を適用すると、
$\mathcal{A}^\bullet \to j_{U, *}\mathcal{C}^\bullet$
は零と鎖ホモトピックである。したがって $a$ は
$\Hom_{D(\QCoh(\mathcal{O}_X))}(A, B'')$ において零へ写る。
ゆえに、$a$ はある
$b \in \Hom_{D(\QCoh(\mathcal{O}_X))}(A, B'[-1])$
の像であると仮定してよい。% P09-ERR-0989
$b$ の $\Hom_{D(\mathcal{O}_X)}(i_X(A), i_X(B')[-1])$ における像は、
ある $\gamma \in \Hom_{D(\mathcal{O}_X)}(A, B''[-1])$ から生じる
（$a$ は右側の群で零へ写るからである）。% P09-ERR-0990
上で水平射が全射であることを見たので、$\gamma$ はある $c$
（$\Hom_{D(\QCoh(\mathcal{O}_X))}(A, B''[-1])$ の元）から生じる。
これは望むとおり $a = 0$ を含意する。

\medskip\noindent
ここまでで、$i_X$ が元の $X$ に対して充満忠実であることが分かった。
$RQ_X$ はその右随伴だから、
$RQ_X \circ i_X = \text{id}$ である
（Categories, Lemma
\ref{categories-lemma-adjoint-fully-faithful}）。
証明を終えるため、任意の $E$（$D_\QCoh(\mathcal{O}_X)$ に属するもの）に対し、
射 $i_X(RQ_X(E)) \to E$ が同型であることを示す。
識別三角形
$$
i_X(RQ_X(E)) \to E \to E' \to i_X(RQ_X(E))[1]
$$
を $D_\QCoh(\mathcal{O}_X)$ において選ぶ。上の議論を用いる形式的な議論により
$i_X(RQ_X(E')) = 0$ である。したがって、$E \in
D_\QCoh(\mathcal{O}_X)$ に対して、条件
$i_X(RQ_X(E)) = 0$ が $E = 0$ を含意することを証明すれば十分である。
$j : V \to X$ を、$V$ をアフィンとする平展射とする。Lemmas
\ref{spaces-perfect-lemma-affine-coherator}、\ref{spaces-perfect-lemma-flat-pushforward-coherator}、
および仮定により、
$$
Rj_*(E|_V) = Rj_*(i_V(RQ_V(E|_V))) = i_X(\Phi(RQ_V(E|_V))) =
i_X(RQ_X(Rj_*(E|_V)))
$$
である。識別三角形
$$
E \to Rj_*(E|_V) \to E' \to E[1]
$$
を選ぶ。$RQ_X$ を適用して識別三角形
$$
0 \to RQ_X(Rj_*(E|_V)) \to RQ_X(E') \to 0[1]
$$
を得る。すなわち中央の射は同型である。
上の一連の等式と合わせると、最初の識別三角形は
識別三角形
$$
E \to i_X(RQ_X(E')) \to E' \to E[1]
$$
となり、中央の射は随伴射である。しかし合成 $E \to E'$ は零なので、
随伴性により $E \to i_X(RQ_X(E'))$ は零である。
この射は、射 $E \to Rj_*(E|_V)$（
$\text{id} : E|_V \to E|_V$ に随伴するもの）と同型だから、
$E|_V$ は零である。これは任意のアフィン $V$（$X$ 上平展）に
対して成り立つので、
望むとおり $E$ は零である。
\end{proof}

\begin{proposition}
\label{spaces-perfect-proposition-quasi-compact-affine-diagonal}
$S$ を概型とする。$X$ を $S$ 上の準コンパクト代数空間で、
$\mathbf{Z}$ 上アフィンな対角射をもつものとする
（Properties of Spaces, Definition
\ref{spaces-properties-definition-separated} の意味で）。
このとき函手 (\ref{spaces-perfect-equation-compare})
$$
D(\QCoh(\mathcal{O}_X))
\longrightarrow
D_\QCoh(\mathcal{O}_X)
$$
は圏同値であり、その準逆は $RQ_X$ によって与えられる。
\end{proposition}

\begin{proof}
$V \to W$ を平展射とし、$V$ はアフィン、$W$ は $X$ の
準コンパクト開部分空間であるとする。射 $V \to W$ はアフィンである。
実際、$W$ は $\mathbf{Z}$ 上アフィンな対角射をもち、$V$ は
アフィンである
（Morphisms of Spaces, Lemma
\ref{spaces-morphisms-lemma-affine-permanence}）。
したがって Lemma \ref{spaces-perfect-lemma-affine-pushforward} により、
Lemma \ref{spaces-perfect-lemma-argument-proves} の仮定が成り立つ。これで結論を得る。
\end{proof}

\begin{lemma}
\label{spaces-perfect-lemma-direct-image-coherator}
$S$ を概型とし、$f : X \to Y$ を $S$ 上の代数空間の射とする。
$X$ と $Y$ は準コンパクトで、$\mathbf{Z}$ 上アフィンな対角射を
もつと仮定する
（Properties of Spaces, Definition
\ref{spaces-properties-definition-separated} の意味で）。
函手
$$
\Phi : D(\QCoh(\mathcal{O}_X)) \to D(\QCoh(\mathcal{O}_Y))
$$
で、
$f_* : \QCoh(\mathcal{O}_X) \to \QCoh(\mathcal{O}_Y)$
の右導来函手を表す。このとき図式
$$
\xymatrix{
D(\QCoh(\mathcal{O}_X)) \ar[d]_\Phi \ar[r] &
D_\QCoh(\mathcal{O}_X) \ar[d]^{Rf_*} \\
D(\QCoh(\mathcal{O}_Y)) \ar[r] &
D_\QCoh(\mathcal{O}_Y)
}
$$
は可換である。
\end{lemma}

\begin{proof}
図式の水平射は Proposition
\ref{spaces-perfect-proposition-quasi-compact-affine-diagonal} により圏同値である。
したがって、これらの圏を同一視できる
（アフィンな対角射をもつ他の準コンパクト代数空間についても同様）。
すると補題の主張は、標準的な射
$\Phi(K) \to Rf_*(K)$ が、すべての $K$
（$D(\QCoh(\mathcal{O}_X))$ に属するもの）に対して
同型であるということである。
$K_1 \to K_2 \to K_3 \to K_1[1]$ が
$D(\QCoh(\mathcal{O}_X))$ の識別三角形であり、主張が三つのうち
二つについて成り立つならば、残る一つについても成り立つことに注意する。

\medskip\noindent
$\mathcal{B} \subset \Ob(X_{spaces, \etale})$ を、準コンパクトで
アフィンな対角射をもつ対象の集合とする。
$U \in \mathcal{B}$ とし、任意の射 $g : U \to Z$ を考える。
ここで $Z$ は $S$ 上の準コンパクト代数空間で、
アフィンな対角射をもつものとする。
$$
\Phi_g : D(\QCoh(\mathcal{O}_U)) \to D(\QCoh(\mathcal{O}_Z))
$$
で $g_*$ の導来延長を表す。次の性質を $P(U) =$ とする。
「任意の $K$（$D(\QCoh(\mathcal{O}_U))$ に属するもの）と、
上のような任意の $g : U \to Z$ に対して、射
$\Phi_g(K) \to Rg_*K$ は同型である。」
Remark \ref{spaces-perfect-remark-how-to} により、Lemma
\ref{spaces-perfect-lemma-induction-principle-separated} の条件 (1)、(2)、(3)(a) は
成り立つので、(3)(b) と (4) を証明すればよい。

\medskip\noindent
条件 (3)(b) を確認する。$U$ を $X$ 上平展なアフィン概型とする。
$g : U \to Z$ を上のような射とする。$Z$ の対角射はアフィンなので、
射 $g : U \to Z$ はアフィンである
（Morphisms of Spaces, Lemma
\ref{spaces-morphisms-lemma-affine-permanence}）。
したがって Lemma \ref{spaces-perfect-lemma-affine-pushforward} により $P(U)$ は成り立つ。

\medskip\noindent
条件 (4) を確認する。
$(U \subset W, V \to W)$ を $X_{spaces, \etale}$ の
初等識別正方形とし、$U, W, V$ は $\mathcal{B}$ に属し、
$V$ はアフィンであるとする。$P$ は $U$、$V$、および
$U \times_W V$ について成り立つと仮定する。
$P$ が $W$ について成り立つことを示す必要がある。
$g : W \to Z$ を、アフィンな対角射をもつ準コンパクト代数空間への
射とする。$K$ を $D(\QCoh(\mathcal{O}_W))$ の対象とする。
識別三角形
$$
K \to Rj_{U, *}K|_U \oplus Rj_{V, *}K|_V \to
Rj_{U \times_W V, *}K|_{U \times_W V} \to K[1]
$$
を $D(\mathcal{O}_W)$ において考える。
上で述べた三つのうち二つという性質により、
$\Phi_g(Rj_{U, *}K|_U) \to Rg_*(Rj_{U, *}K|_U)$ が
同型であることと、$V$ および $U \times_W V$ について同様のことを
示せば十分である。これを次の段落で論じる。

\medskip\noindent
$j : U \to W$ を $X_{spaces, \etale}$ の射とし、
$U, W$ は $\mathcal{B}$ に属し、$P$ は $U$ について成り立つとする。
$g : W \to Z$ を、アフィンな対角射をもつ準コンパクト代数空間への
射とする。証明を終えるには、
$\Phi_g(Rj_*K) \to Rg_*(Rj_*K)$ が任意の $K$
（$D(\QCoh(\mathcal{O}_U))$ に属するもの）に対して
同型であることを示す必要がある。
$\mathcal{I}^\bullet$ を $\QCoh(\mathcal{O}_U)$ の K-入射的複体で
$K$ を表すものとする。$P(U)$ を $j$ に適用すると、
$j_*\mathcal{I}^\bullet$ は $Rj_*K$ を表す。
函手
$j_* : \QCoh(\mathcal{O}_U) \to \QCoh(\mathcal{O}_X)$
は完全な左随伴
$j^* : \QCoh(\mathcal{O}_X) \to \QCoh(\mathcal{O}_U)$
をもつので、$j_*\mathcal{I}^\bullet$ は
$\QCoh(\mathcal{O}_W)$ の K-入射的複体である。% P09-ERR-0991
Derived Categories, Lemma
\ref{derived-lemma-adjoint-preserve-K-injectives} を参照されたい。
したがって $\Phi_g(Rj_*K)$ は
$g_*j_*\mathcal{I}^\bullet = (g \circ j)_*\mathcal{I}^\bullet$
によって表される。$P(U)$ を $g \circ j$ に適用すると、これは
$R_{g \circ j, *}(K) = Rg_*(Rj_*K)$ を表す。これで証明が完了する。
\end{proof}










\section{ネーター空間に対するコヒーレータ}
\label{spaces-perfect-section-coherator-Noetherian}

\noindent
ネーター代数空間の場合を扱うため、入射加群についてもう少し準備が必要である。

\begin{lemma}
\label{spaces-perfect-lemma-affine-injective-colimit-direct-sum-pushforwards-artin}
$S$ をネーター・アフィン概型とする。$\QCoh(\mathcal{O}_S)$ の任意の
入射対象は、次の形の準連接層からなるフィルター余極限
$\colim_i \mathcal{F}_i$ である：
$$
\mathcal{F}_i = (Z_i \to S)_*\mathcal{G}_i
$$
ここで $Z_i$ は Artin 環のスペクトルであり、$\mathcal{G}_i$ は
$Z_i$ 上の連接加群である。
\end{lemma}

\begin{proof}
$S = \Spec(A)$ と置く。$\mathcal{J}$ を $\QCoh(\mathcal{O}_S)$ の
入射対象とする。$\QCoh(\mathcal{O}_S)$ は $A$-加群の圏と同値なので、
$\mathcal{J}$ は $\widetilde{J}$ に等しい。ここで $A$-加群 $J$ は
入射的である。
Dualizing Complexes, Proposition
\ref{dualizing-proposition-structure-injectives-noetherian} により、
$J = \bigoplus E_\alpha$ と書ける。ここで $E_\alpha$ は直既約であり、
したがってある点における剰余体の入射包と同型である。% P09-ERR-0992
ゆえに（Artin 概型の有限非交和は Artin 的なので）、
$J$ は $\kappa(\mathfrak p)$ の入射包であると仮定してよい。
ここで $\mathfrak p$ は $A$ の素イデアルである。このとき
$J = \bigcup J[\mathfrak p^n]$ であり、$J[\mathfrak p^n]$ は
$\kappa(\mathfrak p)$ の、$A_\mathfrak/\mathfrak p^nA_\mathfrak p$ 上の
入射包である。% P09-ERR-0993
Dualizing Complexes, Lemma \ref{dualizing-lemma-union-artinian} を参照されたい。
したがって $\widetilde{J}$ は層 $(Z_n \to X)_*\mathcal{G}_n$ の余極限である。
% P09-ERR-0994
ここで $Z_n = \Spec(A_\mathfrak p/\mathfrak p^nA_\mathfrak p)$ であり、
$\mathfrak G_n$ は有限 $A_\mathfrak/\mathfrak p^nA_\mathfrak p$-加群
$J[\mathfrak p^n]$ に付随する連接層である。% P09-ERR-0995
有限性は Dualizing Complexes, Lemma
\ref{dualizing-lemma-finite} から従う。
\end{proof}

\begin{lemma}
\label{spaces-perfect-lemma-injective-colimit-direct-sum-pushforwards-artin}
$S$ をアフィン概型とする。$X$ を $S$ 上のネーター代数空間とする。
$\QCoh(\mathcal{O}_X)$ の任意の入射対象は、次の形の準連接層からなる
フィルター余極限 $\colim_i \mathcal{F}_i$ の
直和因子である：
$$
\mathcal{F}_i = (Z_i \to X)_*\mathcal{G}_i
$$
ここで $Z_i$ は Artin 環のスペクトルであり、$\mathcal{G}_i$ は
$Z_i$ 上の連接加群である。
\end{lemma}

\begin{proof}
アフィン概型 $U$ と全射平展射 $j : U \to X$ を選ぶ
（Properties of Spaces, Lemma
\ref{spaces-properties-lemma-quasi-compact-affine-cover}）。
このとき $U$ はネーター・アフィン概型である。
入射対象 $\mathcal{J}'$（$\QCoh(\mathcal{O}_U)$ のもの）で、単射
$\mathcal{J}|_U \to \mathcal{J}'$ が存在するものを選ぶ。このとき
$$
\mathcal{J} \to j_*\mathcal{J}'
$$
は $\QCoh(\mathcal{O}_X)$ における単射であり、したがって
$\mathcal{J}$ を $j_*\mathcal{J}'$ の直和因子と同一視する。
ゆえに結果は、Lemma
\ref{spaces-perfect-lemma-affine-injective-colimit-direct-sum-pushforwards-artin} で証明した
$\mathcal{J}'$ に対する対応する結果から従う。
\end{proof}

\begin{lemma}
\label{spaces-perfect-lemma-flat-pullback-injective-quasi-coherent}
$S$ を概型とする。$f : X \to Y$ を $S$ 上の代数空間の平坦、準コンパクト、
かつ準分離的な射とする。$\mathcal{J}$ が $\QCoh(\mathcal{O}_X)$ の
入射対象ならば、$f_*\mathcal{J}$ は $\QCoh(\mathcal{O}_Y)$ の
入射対象である。
\end{lemma}

\begin{proof}
$f$ は準コンパクトかつ準分離的なので、函手 $f_*$ は準連接層を
準連接層へ移す
（Morphisms of Spaces, Lemma \ref{spaces-morphisms-lemma-pushforward}）。
函手 $f^*$ は $f_*$ の左随伴であり、単射を単射へ移す。
したがって結果は Homology, Lemma
\ref{homology-lemma-adjoint-preserve-injectives} から従う。
\end{proof}

\begin{lemma}
\label{spaces-perfect-lemma-injective-pushforward}
$S$ を概型とする。$X$ を $S$ 上のネーター代数空間とする。
$\mathcal{J}$ が $\QCoh(\mathcal{O}_X)$ の入射対象ならば、
\begin{enumerate}
\item $H^p(U, \mathcal{J}|_U) = 0$ である。ここで $p > 0$ であり、
$U$ は $X$ 上平展な任意の準コンパクトかつ準分離的な代数空間である。
\item 任意の射 $f : X \to Y$（$S$ 上の代数空間の射）で、$Y$ が
準分離的なものに対して $R^pf_*\mathcal{J} = 0$ である。ここで
$p > 0$ である。
\end{enumerate}
\end{lemma}

\begin{proof}
(1) の証明。Lemma
\ref{spaces-perfect-lemma-injective-colimit-direct-sum-pushforwards-artin} のように、
$\mathcal{J}$ を $\colim \mathcal{F}_i$ の直和因子として書く。ここで
$\mathcal{F}_i = (Z_i \to X)_*\mathcal{G}_i$ である。
$\colim \mathcal{F}_i$ に対して消滅を証明すれば十分であることは明らかである。
引き戻しは余極限と可換し、$U$ は準コンパクトかつ準分離的なので、
$H^p(U, \mathcal{F}_i|_U) = 0$ を $p > 0$ に対して証明すれば十分である。
Cohomology of Spaces, Lemma
\ref{spaces-cohomology-lemma-colimits} を参照されたい。
射 $Z_i \to X$ はアフィンであることに注意する
（Morphisms of Spaces, Lemma
\ref{spaces-morphisms-lemma-Artinian-affine}）。したがって
$$
\mathcal{F}_i|_U = (Z_i \times_X U \to U)_*\mathcal{G}'_i =
R(Z_i \times_X U \to U)_*\mathcal{G}'_i
$$
である。ここで $\mathcal{G}'_i$ は $\mathcal{G}_i$ の
$Z_i \times_X U$ への引き戻しである。Cohomology of Spaces, Lemma
\ref{spaces-cohomology-lemma-affine-base-change} を参照されたい。
$Z_i \times_X U$ はアフィンなので、$\mathcal{G}'_i$ は
$Z_i \times_X U$ 上で高次コホモロジーをもたない。
Leray スペクトル系列により、$\mathcal{F}_i|_U$ に対しても同じことが成り立つ
（Cohomology on Sites, Lemma
\ref{sites-cohomology-lemma-apply-Leray}）。

\medskip\noindent
(2) の証明。$f : X \to Y$ を $S$ 上の代数空間の射とする。
平展射 $V \to Y$ で、$V$ がアフィンなものをとる。このとき
$V \times_Y X \to X$ は平展射であり、$V \times_Y X$ は
$X$ 上平展な準コンパクトかつ準分離的な代数空間である
（詳細は省略する）。したがって (1) により
$H^p(V \times_Y X, \mathcal{J})$ は零である。
$R^pf_*\mathcal{J}$ は前層
$V \mapsto H^p(V \times_Y X, \mathcal{J})$ に付随する層なので、
結果が従う。
\end{proof}

\begin{lemma}
\label{spaces-perfect-lemma-Noetherian-pushforward}
$S$ を概型とする。$f : X \to Y$ を $S$ 上のネーター代数空間の射とする。
このとき準連接層上の $f_*$ は右導来延長
$\Phi : D(\QCoh(\mathcal{O}_X)) \to D(\QCoh(\mathcal{O}_Y))$
をもち、図式
$$
\xymatrix{
D(\QCoh(\mathcal{O}_X)) \ar[d]_{\Phi} \ar[r] &
D_\QCoh(\mathcal{O}_X) \ar[d]^{Rf_*} \\
D(\QCoh(\mathcal{O}_Y)) \ar[r] &
D_\QCoh(\mathcal{O}_Y)
}
$$
は可換である。
\end{lemma}

\begin{proof}
$X$ と $Y$ はネーターなので、この射は準コンパクトかつ準分離的である
（Morphisms of Spaces, Lemma
\ref{spaces-morphisms-lemma-quasi-compact-quasi-separated-permanence}）。
したがって $f_*$ は準連接性を保つ。Morphisms of Spaces, Lemma
\ref{spaces-morphisms-lemma-pushforward} を参照されたい。
次に、$K$ を $D(\QCoh(\mathcal{O}_X))$ の対象とする。
$\QCoh(\mathcal{O}_X)$ は Grothendieck アーベル圏なので
（Properties of Spaces, Proposition
\ref{spaces-properties-proposition-coherator}）、$K$ を K-入射的複体
$\mathcal{I}^\bullet$ によって表せる。ここで各 $\mathcal{I}^n$ は
$\QCoh(\mathcal{O}_X)$ の入射対象である。Injectives, Theorem
\ref{injectives-theorem-K-injective-embedding-grothendieck} を参照されたい。
したがって、函手 $\Phi$ は
$$
\Phi(K) = f_*\mathcal{I}^\bullet
$$
と置くことで定義される。ここで右辺は
$D(\QCoh(\mathcal{O}_Y))$ の対象とみなす。
補題の証明を終えるには、標準的な射
$$
f_*\mathcal{I}^\bullet \longrightarrow Rf_*\mathcal{I}^\bullet
$$
が $D(\mathcal{O}_Y)$ における同型であることを示せば十分である。
そのためには、この射がコホモロジー層上で同型を誘導することを示せばよい。
任意の $m \in \mathbf{Z}$ を選ぶ。$N = N(X, Y, f)$ を Lemma
\ref{spaces-perfect-lemma-quasi-coherence-direct-image} のものとする。短完全列
$$
0 \to \sigma_{\geq m - N - 1}\mathcal{I}^\bullet \to
\mathcal{I}^\bullet \to \sigma_{\leq m - N - 2}\mathcal{I}^\bullet \to 0
$$
を $X$ 上の準連接層の複体の列として考える。
Lemma \ref{spaces-perfect-lemma-quasi-coherence-direct-image} により、
$Rf_*\sigma_{\leq m - N - 2}\mathcal{I}^\bullet$ のコホモロジー層は
$\geq m - 1$ 次で零である。したがって
$R^mf_*\mathcal{I}^\bullet$ は
$R^mf_*\sigma_{\geq m - N - 1}\mathcal{I}^\bullet$ と同型である。
言い換えると、$\mathcal{I}^\bullet$ は
$\QCoh(\mathcal{O}_X)$ の入射対象からなる下に有界な複体であると
仮定してよい。この場合は Leray の非輪状性補題
（Derived Categories, Lemma
\ref{derived-lemma-leray-acyclicity}）から従い、必要な消滅は Lemma
\ref{spaces-perfect-lemma-injective-pushforward} によって与えられる。
\end{proof}

\begin{proposition}
\label{spaces-perfect-proposition-Noetherian}
$S$ を概型とする。$X$ を $S$ 上のネーター代数空間とする。
このとき函手 (\ref{spaces-perfect-equation-compare})
$$
D(\QCoh(\mathcal{O}_X))
\longrightarrow
D_\QCoh(\mathcal{O}_X)
$$
は圏同値であり、その準逆は $RQ_X$ によって与えられる。
\end{proposition}

\begin{proof}
Lemmas \ref{spaces-perfect-lemma-Noetherian-pushforward} および
\ref{spaces-perfect-lemma-argument-proves} から直ちに従う。
\end{proof}













\section{擬連接複体と完全複体}
\label{spaces-perfect-section-spell-out}

\noindent
本節では、代数空間の平展サイトに対して Cohomology on Sites, Sections
\ref{sites-cohomology-section-strictly-perfect}、
\ref{sites-cohomology-section-pseudo-coherent}、
\ref{sites-cohomology-section-tor}、および
\ref{sites-cohomology-section-perfect} で定義された一般概念を調べる。
特に、これらを概型の場合に起こることと対応させる。

\medskip\noindent
まず、概型上の擬連接複体という概念と、それに付随する小平展サイト上の
擬連接複体という概念を比較する。

\begin{lemma}
\label{spaces-perfect-lemma-descend-finite-type}
$X$ を概型とする。$\mathcal{F}$ を $\mathcal{O}_X$-加群とする。
次の条件は同値である。
\begin{enumerate}
\item $\mathcal{F}$ は $\mathcal{O}_X$-加群として有限型である。
\item $\epsilon^*\mathcal{F}$ は $\mathcal{O}_\etale$-加群として有限型である。
ここで考えるのは $X$ の小平展サイトである。
\end{enumerate}
ここで $\epsilon$ は (\ref{spaces-perfect-equation-epsilon}) のものである。
\end{lemma}

\begin{proof}
(1) $\Rightarrow$ (2) は一般的事実である。Modules on Sites, Lemma
\ref{sites-modules-lemma-local-pullback} を参照されたい。
(2) を仮定する。仮定により平展被覆 $\{f_i : X_i \to X\}$ で、
$\epsilon^*\mathcal{F}|_{(X_i)_\etale}$ が有限個の切断で生成されるものが
存在する。$x \in X$ とする。$\mathcal{F}$ が $x$ の近傍で有限個の
切断によって生成されることを示す。$x$ は $X_i \to X$ の像に属するとし、
$X' = X_i$ と書く。生成切断を
$s_1, \ldots, s_n \in
\Gamma(X', \epsilon^*\mathcal{F}|_{X'_\etale})$
とする。等式
$\epsilon^*\mathcal{F} =
\epsilon^{-1}\mathcal{F} \otimes_{\epsilon^{-1}\mathcal{O}_X}
\mathcal{O}_\etale$
により、平展射 $X'' \to X'$ で、$x$ が $X' \to X$ の像に属し、
% P09-ERR-0996
$s_i|_{X''} = \sum s_{ij} \otimes a_{ij}$ となるものを見つけられる。
ここで $s_{ij} \in \epsilon^{-1}\mathcal{F}(X'')$ および
$a_{ij} \in \mathcal{O}_\etale(X'')$ は切断である。
$U \subset X$ を $X'' \to X$ の像とする。$f'' : X'' \to X$ は平展なので、
これは開部分概型である
（Morphisms, Lemma \ref{morphisms-lemma-etale-open}）。
$X''$ をさらに縮小すれば、$s_{ij}$ はある元
$t_{ij} \in \mathcal{F}(U)$ から来ると仮定してよい。
これは逆像函手 $\epsilon^{-1}$ の構成から従う。
このとき $t_{ij}$ が $\mathcal{F}|_U$ を生成すると主張する。
これで補題の証明が終わる。実際、対応する写像
$\mathcal{O}_U^{\oplus N} \to \mathcal{F}|_U$ は、その $f''$ による
$X''$ への引き戻しが全射であるという性質をもつ。
$f'' : X'' \to U$ は概型の全射平坦射なので、
$\mathcal{O}_U^{\oplus N} \to \mathcal{F}|_U$ は全射である。
実際、茎を調べ、$\mathcal{O}_{U, f''(z)} \to \mathcal{O}_{X'', z}$ が
すべての $z \in X''$ に対して忠実平坦であることを用いればよい。
\end{proof}

\noindent
上の状況では、サイトの射 $\epsilon$ は平坦であり、したがって加群の複体上の
引き戻しを定める。

\begin{lemma}
\label{spaces-perfect-lemma-descend-pseudo-coherent}
$X$ を概型とする。$E$ を $D(\mathcal{O}_X)$ の対象とする。
次の条件は同値である。
\begin{enumerate}
\item $E$ は $m$-擬連接である。
\item $\epsilon^*E$ は $m$-擬連接である。ここで考えるのは
$X$ の小平展サイトである。
\end{enumerate}
ここで $\epsilon$ は (\ref{spaces-perfect-equation-epsilon}) のものである。
\end{lemma}

\begin{proof}
(1) $\Rightarrow$ (2) は一般的事実である。Cohomology on Sites, Lemma
\ref{sites-cohomology-lemma-pseudo-coherent-pullback} を参照されたい。
$\epsilon^*E$ が $m$-擬連接であると仮定する。以下、$\epsilon^*$ は
完全函手であり、したがって
$$
\epsilon^*H^i(E) = H^i(\epsilon^*E).
$$
であることを断りなく用いる。$E$ が $m$-擬連接であることを示すには
$X$ 上局所的に議論してよいので、$X$ は準コンパクトである
（例えばアフィンである）と仮定してよい。$X$ は準コンパクトなので、
任意の平展被覆 $\{U_i \to X\}$ は有限細分をもつ。
したがって $\epsilon^*E$ は $D^{-}(\mathcal{O}_\etale)$ の対象である。
Cohomology on Sites, Definition
\ref{sites-cohomology-definition-pseudo-coherent} に続くコメントを参照されたい。
Lemma \ref{spaces-perfect-lemma-epsilon-flat} により、$E$ は
$D^-(\mathcal{O}_X)$ の対象である。

\medskip\noindent
$n \in \mathbf{Z}$ を $H^n(E)$ が非零となる最大の整数とする。このとき
$n$ は $H^n(\epsilon^*E)$ が非零となる最大の整数でもある。
$n - m$ に関する帰納法で補題を証明する。$n < m$ ならば補題は明らかである。
$n \geq m$ ならば、$H^n(\epsilon^*E)$ は有限
$\mathcal{O}_\etale$-加群である。Cohomology on Sites, Lemma
\ref{sites-cohomology-lemma-finite-cohomology} を参照されたい。
したがって Lemma \ref{spaces-perfect-lemma-descend-finite-type} により、$H^n(E)$ は
有限 $\mathcal{O}_X$-加群である。$X$ を開被覆の各要素で置き換えれば、
全射 $\mathcal{O}_X^{\oplus t} \to H^n(E)$ が存在すると仮定してよい。
$X$ 上局所的に、これを複体の射
$\alpha : \mathcal{O}_X^{\oplus t}[-n] \to E$ に持ち上げられる
（詳細は省略する）。識別三角形
$$
\mathcal{O}_X^{\oplus t}[-n] \to E \to C \to \mathcal{O}_X^{\oplus t}[-n + 1]
$$
を選ぶ。このとき $C$ のコホモロジーは $\geq n$ 次で消滅する。
一方、複体 $\epsilon^*C$ は $m$-擬連接である。Cohomology on Sites, Lemma
\ref{sites-cohomology-lemma-cone-pseudo-coherent} を参照されたい。
したがって帰納法により $C$ は $m$-擬連接である。
Cohomology on Sites, Lemma
\ref{sites-cohomology-lemma-cone-pseudo-coherent} をもう一度適用して結論を得る。
\end{proof}

\begin{lemma}
\label{spaces-perfect-lemma-descend-tor-amplitude}
$X$ を概型とする。$E$ を $D(\mathcal{O}_X)$ の対象とする。このとき
\begin{enumerate}
\item $E$ が $[a, b]$ に Tor 振幅をもつことと、$\epsilon^*E$ が
$[a, b]$ に Tor 振幅をもつこととは同値である。
\item $E$ が有限 Tor 次元をもつことと、$\epsilon^*E$ が有限
Tor 次元をもつこととは同値である。
\end{enumerate}
ここで $\epsilon$ は (\ref{spaces-perfect-equation-epsilon}) のものである。
\end{lemma}

\begin{proof}
容易な向きは Cohomology on Sites, Lemma
\ref{sites-cohomology-lemma-tor-amplitude-pullback} から従う。
逆向きについて、$\epsilon^*E$ が $[a, b]$ に Tor 振幅をもつと仮定する。
$\mathcal{F}$ を $\mathcal{O}_X$-加群とする。$\epsilon$ は環付きサイトの
平坦射なので（Lemma \ref{spaces-perfect-lemma-epsilon-flat}）、
$$
\epsilon^*(E \otimes^\mathbf{L}_{\mathcal{O}_X} \mathcal{F})
=
\epsilon^*E
\otimes^\mathbf{L}_{\mathcal{O}_\etale}
\epsilon^*\mathcal{F}
$$
である。したがって、右辺のコホモロジー層が消滅するという仮定から、
Lemma \ref{spaces-perfect-lemma-epsilon-flat} を介して
$E \otimes^\mathbf{L}_{\mathcal{O}_X} \mathcal{F}$ のコホモロジー層に
望まれる消滅が従う。
\end{proof}

\begin{lemma}
\label{spaces-perfect-lemma-tor-dimension-rel}
$f : X \to Y$ を概型の射とする。$E$ を $D(\mathcal{O}_X)$ の対象とする。
このとき
\begin{enumerate}
\item $E$ が $D(f^{-1}\mathcal{O}_Y)$ の対象として $[a, b]$ に
Tor 振幅をもつことと、$\epsilon^*E$ が $[a, b]$ に Tor 振幅をもつこととは
同値である。後者では $D(f_{small}^{-1}\mathcal{O}_{Y_\etale})$ の対象として
考える。
\item $E$ が $D(f^{-1}\mathcal{O}_Y)$ の対象として局所有限 Tor 次元を
もつことと、$\epsilon^*E$ が
$D(f_{small}^{-1}\mathcal{O}_{Y_\etale})$ の対象として局所有限 Tor 次元を
もつこととは同値である。
\end{enumerate}
ここで $\epsilon$ は (\ref{spaces-perfect-equation-epsilon}) のものである。
\end{lemma}

\begin{proof}
(1) の容易な向きは Cohomology on Sites, Lemma
\ref{sites-cohomology-lemma-tor-amplitude-pullback} から従う。
$x \in X$ を点とし、$\overline{x}$ を $x$ 上の幾何学的点とする。
$y = f(x)$ と置き、$\overline{y}$ を $Y$ の幾何学的点で、
$\overline{x}$ から得られるものと書く。このとき
$(f^{-1}\mathcal{O}_Y)_x = \mathcal{O}_{Y, y}$
（Sheaves, Lemma \ref{sheaves-lemma-stalk-pullback}）であり、
$$
(f_{small}^{-1}\mathcal{O}_{Y_\etale})_{\overline{x}} =
\mathcal{O}_{Y_\etale, \overline{y}} =
\mathcal{O}_{Y, y}^{sh}
$$
は強 Hensel 化である
（\'Etale Cohomology, Lemmas
\ref{etale-cohomology-lemma-stalk-pullback} および
\ref{etale-cohomology-lemma-describe-etale-local-ring}）。
$\mathcal{O}_{X_\etale}$ の $X$ における茎は
$\mathcal{O}_{X, x}^{sh}$ なので、% P09-ERR-0997
引き戻しの推移性により
$$
(\epsilon^*E)_{\overline{x}} =
E_x \otimes_{\mathcal{O}_{X, x}}^\mathbf{L} \mathcal{O}_{X, x}^{sh}
$$
を得る。$\epsilon^*E$ が $[a, b]$ に Tor 振幅をもつ
$f_{small}^{-1}\mathcal{O}_{Y_\etale}$-加群の複体ならば、
$(\epsilon^*E)_{\overline{x}}$ は $[a, b]$ に Tor 振幅をもつ
$\mathcal{O}_{Y, y}^{sh}$-加群の複体である
（茎をとることは引き戻しであり、上で引用した補題を用いる）。
More on Flatness, Lemma
\ref{flat-lemma-tor-amplitude-up-down-henselization} により、
$(\epsilon^*E)_{\overline{x}}$ の $\mathcal{O}_{Y, y}$-加群の複体としての
Tor 振幅は $[a, b]$ に含まれる。射
$\mathcal{O}_{X, x} \to \mathcal{O}_{X, x}^{sh}$ は忠実平坦であり
（More on Algebra, Lemma
\ref{more-algebra-lemma-dumb-properties-henselization}）、かつ
$$
(\epsilon^*E)_{\overline{x}} =
E_x \otimes_{\mathcal{O}_{X, x}}^\mathbf{L} \mathcal{O}_{X, x}^{sh}
$$
なので、More on Algebra, Lemma
\ref{more-algebra-lemma-no-change-tor-amplitude} を適用できる。
これにより、$E_x$ の $\mathcal{O}_{Y, y}$-加群の複体としての
Tor 振幅は $[a, b]$ に含まれる。Cohomology, Lemma
\ref{cohomology-lemma-tor-amplitude-stalk} により、$E$ は
$D(f^{-1}\mathcal{O}_Y)$ の対象として $[a, b]$ に Tor 振幅をもつ。
これで (1) の逆向きを得る。(2) は (1) から形式的に従う。
\end{proof}

\begin{lemma}
\label{spaces-perfect-lemma-descend-perfect}
$X$ を概型とする。$E$ を $D(\mathcal{O}_X)$ の対象とする。
$E$ が $D(\mathcal{O}_X)$ の完全対象であることと、
$\epsilon^*E$ が $D(\mathcal{O}_\etale)$ の完全対象であることとは同値である。
ここで $\epsilon$ は (\ref{spaces-perfect-equation-epsilon}) のものである。
\end{lemma}

\begin{proof}
容易な含意は Cohomology on Sites, Lemma
\ref{sites-cohomology-lemma-perfect-pullback} に含まれる一般的結果から従う。
逆向きには、Cohomology on Sites, Lemma
\ref{sites-cohomology-lemma-perfect} の同値と、擬連接性および有限 Tor 次元の
複体に対する対応する結果、すなわち Lemmas
\ref{spaces-perfect-lemma-descend-pseudo-coherent} および
\ref{spaces-perfect-lemma-descend-tor-amplitude} を用いればよい。
いくつかの詳細は省略する。
\end{proof}


\begin{lemma}
\label{spaces-perfect-lemma-pseudo-coherent}
$S$ を概型とする。$X$ を $S$ 上の代数空間とする。
$E$ が $m$-擬連接対象（$D(\mathcal{O}_X)$ のもの）ならば、
$H^i(E)$ は準連接 $\mathcal{O}_X$-加群である。ここで $i > m$ である。
$E$ が擬連接ならば、$E$ は $D_\QCoh(\mathcal{O}_X)$ の対象である。
\end{lemma}

\begin{proof}
局所的に $H^i(E)$ は $H^i(\mathcal{E}^\bullet)$ と同型であり、
$\mathcal{E}^\bullet$ は厳密完全である。層 $\mathcal{E}^i$ は有限自由加群の
直和因子なので、準連接である。これで補題が従う。
\end{proof}

\begin{lemma}
\label{spaces-perfect-lemma-identify-pseudo-coherent-noetherian}
$S$ を概型とする。$X$ を $S$ 上のネーター代数空間とする。
$E$ を $D_\QCoh(\mathcal{O}_X)$ の対象とする。
$m \in \mathbf{Z}$ に対して、次の条件は同値である。
\begin{enumerate}
\item $H^i(E)$ は $i \geq m$ に対して連接であり、$i \gg 0$ に対して零である。
\item $E$ は $m$-擬連接である。
\end{enumerate}
特に、$E$ が擬連接であることと、$E$ が
$D^-_{\textit{Coh}}(\mathcal{O}_X)$ の対象であることとは同値である。
\end{lemma}

\begin{proof}
$X$ は準コンパクトなので、アフィン概型 $U$ と全射平展射
$U \to X$ を見つけられる
（Properties of Spaces, Lemma
\ref{spaces-properties-lemma-quasi-compact-affine-cover}）。
$U$ はネーターであることに注意する。$E$ が $m$-擬連接であることと
$E|_U$ が $m$-擬連接であることとは同値である
（定義、または Cohomology on Sites, Lemma
\ref{sites-cohomology-lemma-pseudo-coherent-independent-representative} から従う）。
同様に、$H^i(E)$ が連接であることと、
$H^i(E)|_U = H^i(E|_U)$ が連接であることとは同値である
（Cohomology of Spaces, Lemma
\ref{spaces-cohomology-lemma-coherent-Noetherian}）。
したがって $X$ は表現可能であると仮定してよい。

\medskip\noindent
$X$ が概型 $X_0$ によって表現されるならば
（Lemma \ref{spaces-perfect-lemma-derived-quasi-coherent-small-etale-site}）、
$E = \epsilon^*E_0$ と書ける。ここで $E_0$ は
$D_\QCoh(\mathcal{O}_{X_0})$ の対象であり、
$\epsilon : X_\etale \to (X_0)_{Zar}$ は
(\ref{spaces-perfect-equation-epsilon}) のものである。この場合、Lemma
\ref{spaces-perfect-lemma-descend-pseudo-coherent} により、$E$ が $m$-擬連接であることと
$E_0$ がそうであることとは同値である。同様に、Lemma
\ref{spaces-perfect-lemma-descend-finite-type} により、$H^i(E_0)$ が有限型
（すなわち連接）であることと $H^i(E)$ がそうであることとは同値である。
最後に、Lemma \ref{spaces-perfect-lemma-epsilon-flat} により、
$H^i(E_0) = 0$ であることと $H^i(E) = 0$ であることとは同値である。
したがって概型の場合に帰着し、これは Derived Categories of Schemes, Lemma
\ref{perfect-lemma-identify-pseudo-coherent-noetherian} である。
\end{proof}

\begin{lemma}
\label{spaces-perfect-lemma-tor-qc-qs}
$S$ を概型とする。$X$ を $S$ 上の準分離的代数空間とする。
$E$ を $D_\QCoh(\mathcal{O}_X)$ の対象とする。$a \leq b$ とする。
次の条件は同値である。
\begin{enumerate}
\item $E$ は $[a, b]$ に Tor 振幅をもつ。
\item すべての $\mathcal{F}$（$\QCoh(\mathcal{O}_X)$ に属するもの）に対して、
$H^i(E \otimes_{\mathcal{O}_X}^\mathbf{L} \mathcal{F}) = 0$ である。
ここで $i \not \in [a, b]$ である。
\end{enumerate}
\end{lemma}

\begin{proof}
(1) が (2) を含意することは明らかである。(2) を仮定する。
$j : U \to X$ を $U$ がアフィンである平展射とする。
$X$ は準分離的なので、$j$ は準コンパクトかつ分離的である。
したがって $j_*$ は準連接加群を準連接加群へ移す
（Morphisms of Spaces, Lemma \ref{spaces-morphisms-lemma-pushforward}）。
任意の準連接加群 $\mathcal{G}$（$U$ 上のもの）をとる。仮定により
$H^i(E \otimes_{\mathcal{O}_X}^\mathbf{L} j_*\mathcal{G})$ は
$i \not \in [a, b]$ に対して消滅する。平坦射 $j$ で引き戻すと、
$H^i(E|_U \otimes_{\mathcal{O}_U}^\mathbf{L} j^*j_*\mathcal{G})$ は
$i \not \in [a, b]$ に対して消滅する。Cohomology of Spaces, Lemma
\ref{spaces-cohomology-lemma-etale-pull-push-split} により、
$\mathcal{G}$ は $j^*j_*\mathcal{G}$ の直和因子である。
したがって条件 (2) から、
$H^i(E|_U \otimes_{\mathcal{O}_U}^\mathbf{L} \mathcal{G})$ が
$i \not \in [a, b]$ のとき、すべての準連接 $\mathcal{O}_U$-加群
$\mathcal{G}$ に対して消滅することが従う。
$E|_U$ が $[a, b]$ に Tor 振幅をもつことを証明すれば十分なので、
$X$ が表現可能な場合に帰着する。

\medskip\noindent
$X$ が概型 $X_0$ によって表現されるならば
（Lemma \ref{spaces-perfect-lemma-derived-quasi-coherent-small-etale-site}）、
$E = \epsilon^*E_0$ と書ける。ここで $E_0$ は
$D_\QCoh(\mathcal{O}_{X_0})$ の対象であり、
$\epsilon : X_\etale \to (X_0)_{Zar}$ は
(\ref{spaces-perfect-equation-epsilon}) のものである。任意の準連接加群
$\mathcal{F}_0$（$X_0$ 上のもの）に対して、加群 $\epsilon^*\mathcal{F}_0$ は
$X$ 上準連接であり、
$$
H^i(E \otimes_{\mathcal{O}_X}^\mathbf{L} \epsilon^*\mathcal{F}_0)
=
\epsilon^*H^i(E_0 \otimes_{\mathcal{O}_{X_0}}^\mathbf{L} \mathcal{F}_0)
$$
である。これは $\epsilon$ が平坦だからである
（Lemma \ref{spaces-perfect-lemma-epsilon-flat}）。さらに、これらの層が
$i \not \in [a, b]$ に対して消滅することから、同じ補題により
$H^i(E_0 \otimes_{\mathcal{O}_{X_0}}^\mathbf{L} \mathcal{F}_0)$ に対する
同じ消滅が従う。したがって問題は概型の場合に帰着し、これは
Derived Categories of Schemes, Lemma \ref{perfect-lemma-tor-qc-qs} で扱われる。
\end{proof}

\begin{lemma}
\label{spaces-perfect-lemma-descend-RHom}
$X$ を概型とする。$E, F$ を $D(\mathcal{O}_X)$ の対象とする。
次のいずれかを仮定する。
\begin{enumerate}
\item $E$ は擬連接であり、$F$ は $D^+(\mathcal{O}_X)$ に属する。
\item $E$ は完全であり、$F$ は任意である。
\end{enumerate}
このとき標準的な同型
$$
\epsilon^*R\SheafHom(E, F) \longrightarrow R\SheafHom(\epsilon^*E, \epsilon^*F)
$$
が存在する。ここで $\epsilon$ は (\ref{spaces-perfect-equation-epsilon}) のものである。
\end{lemma}

\begin{proof}
$\epsilon$ は平坦であり（Lemma \ref{spaces-perfect-lemma-epsilon-flat}）、したがって
$\epsilon^* = L\epsilon^*$ であることを思い出す。
Cohomology on Sites, Remark
\ref{sites-cohomology-remark-prepare-fancy-base-change} により、左辺から右辺への
標準的な射が存在する。これが同型であることを示すには局所的に議論してよい。
すなわち $X$ はアフィン概型であると仮定してよい。

\medskip\noindent
(1) の場合、$E$ を上に有界な複体 $\mathcal{E}^\bullet$
（有限自由 $\mathcal{O}_X$-加群からなるもの）で表せる。
Derived Categories of Schemes, Lemma
\ref{perfect-lemma-lift-pseudo-coherent} を参照されたい。
また、$F$ を下に有界な複体 $\mathcal{F}^\bullet$
（$\mathcal{O}_X$-加群からなるもの）で表せる。Cohomology, Lemma
\ref{cohomology-lemma-Rhom-complex-of-direct-summands-finite-free} を適用すると、
$R\SheafHom(E, F)$ は、各項が
$$
\bigoplus\nolimits_{n = - p + q}
\SheafHom_{\mathcal{O}_X}(\mathcal{E}^p, \mathcal{F}^q)
$$
である複体によって表される。Cohomology on Sites, Lemma
\ref{sites-cohomology-lemma-Rhom-complex-of-direct-summands-finite-free} を
適用すると、$R\SheafHom(\epsilon^*E, \epsilon^*F)$ は、各項が
$$
\bigoplus\nolimits_{n = - p + q}
\SheafHom_{\mathcal{O}_\etale}
(\epsilon^*\mathcal{E}^p, \epsilon^*\mathcal{F}^q)
$$
である複体によって表される。したがって補題の主張は、標準的な射
$$
\epsilon^*\SheafHom_{\mathcal{O}_X}(\mathcal{E}, \mathcal{F})
\longrightarrow
\SheafHom_{\mathcal{O}_\etale}
(\epsilon^*\mathcal{E}, \epsilon^*\mathcal{F})
$$
が、任意の $\mathcal{O}_X$-加群 $\mathcal{F}$ と有限自由
$\mathcal{O}_X$-加群 $\mathcal{E}$ に対して同型であるという真の事実に帰着する。

\medskip\noindent
(2) の場合、$E$ を厳密完全複体 $\mathcal{E}^\bullet$
（$\mathcal{O}_X$-加群からなるもの）で表せる。
Derived Categories of Schemes, Lemmas
\ref{perfect-lemma-affine-compare-bounded} および
\ref{perfect-lemma-perfect-affine} と、加群の完全複体が有限生成射影加群の
有限複体によって表されるという事実を用いる。
したがって、上とまったく同じ証明を行える。ただし Cohomology, Lemma
\ref{cohomology-lemma-Rhom-complex-of-direct-summands-finite-free} への参照を
Cohomology, Lemma \ref{cohomology-lemma-Rhom-strictly-perfect} への参照で
置き換える。
\end{proof}

\begin{lemma}
\label{spaces-perfect-lemma-quasi-coherence-internal-hom}
$S$ を概型とする。$X$ を $S$ 上の代数空間とする。
$L, K$ を $D(\mathcal{O}_X)$ の対象とする。次のいずれかを仮定する。
\begin{enumerate}
\item $L$ は $D^+_\QCoh(\mathcal{O}_X)$ に属し、$K$ は擬連接である。
\item $L$ は $D_\QCoh(\mathcal{O}_X)$ に属し、$K$ は完全である。
\end{enumerate}
このとき $R\SheafHom(K, L)$ は $D_\QCoh(\mathcal{O}_X)$ に属する。
\end{lemma}

\begin{proof}
これは、概型に対する類似の結果
（Derived Categories of Schemes, Lemma
\ref{perfect-lemma-quasi-coherence-internal-hom}）から、Lemma
\ref{spaces-perfect-lemma-check-quasi-coherence-on-covering} の判定条件、Lemmas
\ref{spaces-perfect-lemma-descend-pseudo-coherent} および
\ref{spaces-perfect-lemma-descend-perfect} の判定条件、ならびに Lemma
\ref{spaces-perfect-lemma-descend-RHom} の結果を通じて従う。
\end{proof}

\begin{lemma}
\label{spaces-perfect-lemma-internal-hom-evaluate-tensor-isomorphism}
$S$ を概型とする。$X$ を $S$ 上の代数空間とする。
$K, L, M$ を $D_\QCoh(\mathcal{O}_X)$ の対象とする。
Cohomology on Sites, Lemma
\ref{sites-cohomology-lemma-internal-hom-diagonal-better} の射
$$
K \otimes_{\mathcal{O}_X}^\mathbf{L} R\SheafHom(M, L)
\longrightarrow
R\SheafHom(M, K \otimes_{\mathcal{O}_X}^\mathbf{L} L)
$$
は、次の場合に同型である。
\begin{enumerate}
\item $M$ が完全である。
\item $K$ が完全である。
\item $M$ が擬連接で、$L \in D^+(\mathcal{O}_X)$ であり、$K$ が有限
Tor 次元をもつ。
\end{enumerate}
\end{lemma}

\begin{proof}
射が同型であるか否かは平展局所的に確認できるので、$X$ はアフィン概型であると
仮定してよい。このとき、Lemma
\ref{spaces-perfect-lemma-derived-quasi-coherent-small-etale-site} により、$K, L, M$ を
$\epsilon^*K_0, \epsilon^*L_0, \epsilon^*M_0$ と書ける。
ここで $K_0, L_0, M_0$ は $D_\QCoh(\mathcal{O}_X)$ に属する。
すると Lemma \ref{spaces-perfect-lemma-descend-RHom} により、場合 (1) と (3) は概型の場合、
すなわち Derived Categories of Schemes, Lemma
\ref{perfect-lemma-internal-hom-evaluate-tensor-isomorphism} に帰着する。
$K$ が完全であるが他の仮定を置かない場合、矢印のいずれの側も
$D_\QCoh(\mathcal{O}_X)$ に属するとは限らない。しかし結果はなお真である。
なぜなら、$K$ は（さらに局所化した後）有限自由加群の有限複体によって
表され、この場合は明らかだからである。
\end{proof}









\section{完全複体による近似}
\label{spaces-perfect-section-approximation}

\noindent
本節では、Derived Categories of Schemes, Section
\ref{perfect-section-approximation} で始めた議論を続ける。

\begin{definition}
\label{spaces-perfect-definition-approximation-holds}
$S$ を概型とする。$X$ を $S$ 上の代数空間とする。次を満たす三つ組
$(T, E, m)$ を考える。
\begin{enumerate}
\item $T \subset |X|$ は閉部分集合である。
\item $E$ は $D_\QCoh(\mathcal{O}_X)$ の対象である。
\item $m \in \mathbf{Z}$ である。
\end{enumerate}
{it 三つ組 $(T, E, m)$ に対して近似が成り立つ}とは、完全対象 $P$
（$D(\mathcal{O}_X)$ のもの）で $T$ 上に台をもつものと、射
$\alpha : P \to E$ で、同型 $H^i(P) \to H^i(E)$ を $i > m$ に対して、
また全射 $H^m(P) \to H^m(E)$ を誘導するものが存在することをいう。
\end{definition}

\noindent
すべての三つ組に対して近似が成り立つわけではない。その理由については、
Derived Categories of Schemes, Definition
\ref{perfect-definition-approximation-holds} に続く注意を読まれたい。

\begin{definition}
\label{spaces-perfect-definition-approximation}
$S$ を概型とする。$X$ を $S$ 上の代数空間とする。
$X$ 上で{it 完全複体による近似が成り立つ}とは、任意の閉部分集合
$T \subset |X|$ で、射 $X \setminus T \to X$ が準コンパクトであるものに
対して、ある整数 $r$ が存在し、Definition
\ref{spaces-perfect-definition-approximation-holds} の任意の三つ組 $(T, E, m)$ で
次を満たすものに対して近似が成り立つことをいう。
\begin{enumerate}
\item $E$ は $(m - r)$-擬連接である。
\item $H^i(E)$ は $T$ 上に台をもつ。ここで $i \geq m - r$ である。
\end{enumerate}
\end{definition}

\begin{lemma}
\label{spaces-perfect-lemma-pushforward-perfect}
$S$ を概型とする。$(U \subset X, j : V \to X)$ を $S$ 上の代数空間の
初等識別正方形とする。$E$ を $D(\mathcal{O}_V)$ の完全対象で、
$j^{-1}(T)$ 上に台をもつものとする。ここで $T = |X| \setminus |U|$ である。
このとき $Rj_*E$ は $D(\mathcal{O}_X)$ の完全対象である。
\end{lemma}

\begin{proof}
完全であることは $X_\etale$ 上局所的である。したがって、$Rj_*E$ の
$U$ および $V$ への制限が完全であることを確認すれば十分である。
Lemma \ref{spaces-perfect-lemma-pushforward-with-support-in-open} により
$Rj_*E|_V = E$ であり、これは完全である。また、
$Rj_*E|_U = 0$ である。これは $E|_{V \setminus j^{-1}(T)} = 0$ だからである
（Lemma \ref{spaces-perfect-lemma-restrict-direct-image-open} を用いる）。
\end{proof}

\begin{lemma}
\label{spaces-perfect-lemma-open}
$S$ を概型とする。$(U \subset X, j : V \to X)$ を $S$ 上の代数空間の
初等識別正方形とする。$T$ を $|X| \setminus |U|$ の閉部分集合とし、
$(T, E, m)$ を Definition \ref{spaces-perfect-definition-approximation-holds} の三つ組とする。
次を仮定する。
\begin{enumerate}
\item $(j^{-1}T, E|_V, m)$ に対して近似が成り立つ。
\item 層 $H^i(E)$ は $i \geq m$ に対して $T$ 上に台をもつ。
\end{enumerate}
このとき $(T, E, m)$ に対して近似が成り立つ。
\end{lemma}

\begin{proof}
$P \to E|_V$ を、三つ組 $(j^{-1}T, E|_V, m)$ の $V$ 上の近似とする。
Lemma \ref{spaces-perfect-lemma-pushforward-perfect} により、$Rj_*P$ は
$D(\mathcal{O}_X)$ の完全対象である。一方、Lemma
\ref{spaces-perfect-lemma-pushforward-with-support-in-open} により $Rj_*P = j_!P$ である。
例えば (\ref{spaces-perfect-equation-stalk-j-shriek}) により、$j_!P$ は $T$ 上に台をもつ。
したがって近似
$Rj_*P = j_!P \to j_!(E|_V) \to E$ を得る。
\end{proof}

\begin{lemma}
\label{spaces-perfect-lemma-approximation-affine}
$S$ を概型とする。$X$ を $S$ 上の代数空間で、アフィン概型によって
表現されるものとする。このとき Definition
\ref{spaces-perfect-definition-approximation-holds} の任意の三つ組 $(T, E, m)$ で、
次を満たす整数 $r \geq 0$ が存在するものに対して近似が成り立つ。
\begin{enumerate}
\item $E$ は $m$-擬連接である。
\item $H^i(E)$ は $T$ 上に台をもつ。ここで $i \geq m - r + 1$ である。
\item $X \setminus T$ は $r$ 個のアフィン開部分の和である。
\end{enumerate}
特に、アフィン概型に対して完全複体による近似が成り立つ。
\end{lemma}

\begin{proof}
$X_0$ を $X$ を表現するアフィン概型とする。
$T_0 \subset X_0$ を $T$ に対応する閉部分集合とする。% P09-ERR-0998
$\epsilon : X_\etale \to X_{0, Zar}$ を (\ref{spaces-perfect-equation-epsilon}) の射とする。
Lemma \ref{spaces-perfect-lemma-derived-quasi-coherent-small-etale-site} により、
$E = \epsilon^*E_0$ と書ける。ここで $E_0$ は
$D_\QCoh(\mathcal{O}_{X_0})$ の対象である。すると Lemma
\ref{spaces-perfect-lemma-descend-pseudo-coherent} により $E_0$ は $m$-擬連接である。
コホモロジー層の茎を比較すると（Lemma
\ref{spaces-perfect-lemma-epsilon-flat} の証明を参照）、$H^i(E_0)$ は
$T_0$ 上に台をもつことが分かる。ここで $i \geq m - r + 1$ である。
Derived Categories of Schemes, Lemma
\ref{perfect-lemma-approximation-affine} により、近似
$P_0 \to E_0$（$(T_0, E_0, m)$ のもの）が存在する。
Lemma \ref{spaces-perfect-lemma-descend-perfect} により、
$P = \epsilon^*P_0$ は $D(\mathcal{O}_X)$ の完全対象である。
引き戻すと、望む近似
$P = \epsilon^*P_0 \to \epsilon^*E_0 = E$ を得る。
\end{proof}

\begin{lemma}
\label{spaces-perfect-lemma-induction-step}
$S$ を概型とする。$(U \subset X, j : V \to X)$ を $S$ 上の代数空間の
初等識別正方形とする。$U$ は準コンパクト、$V$ はアフィン、
$U \times_X V$ は準コンパクトであると仮定する。
$U$ 上で完全複体による近似が成り立つならば、$X$ 上でも
完全複体による近似が成り立つ。
\end{lemma}

\begin{proof}
$T \subset |X|$ を閉部分集合で、$X \setminus T \to X$ が
準コンパクトであるものとする。$r_U$ を Definition
\ref{spaces-perfect-definition-approximation} の整数で、組 $(U, T \cap |U|)$ に
適合するものとする。$T' = T \setminus |U|$ と置き、$T'$ に誘導された
被約部分空間構造を入れる。$|T'|$ は $|X| \setminus |U|$ に含まれるので、
$j^{-1}(T') \to T'$ は同型である。さらに
$V \setminus j^{-1}(T')$ は準コンパクトである。実際、これは
$U \times_X V$ と $X \setminus T$ の $X$ 上のファイバー積であり、
$U \times_X V$ は準コンパクト、$X \setminus T \to X$ は
準コンパクトであると仮定した。$r'$ を
$V \setminus j^{-1}(T')$ を被覆するのに必要なアフィンの個数とする。
$r = \max(r_U, r')$ が組 $(X, T)$ に対して機能すると主張する。

\medskip\noindent
これを見るため、三つ組 $(T, E, m)$ で、$E$ が
$(m - r)$-擬連接であり、$H^i(E)$ が $T$ 上に台をもち、
$i \geq m - r$ であるものを選ぶ。$t$ を $H^t(E)|_U$ が非零となる
最大の整数とする（$U$ は準コンパクトで、$E|_U$ は
$(m - r)$-擬連接なので、このような整数は存在する）。
$E$ を近似できることを $t$ に関する帰納法により証明する。

\medskip\noindent
基底の場合：$t \leq m - r'$ とする。これは $H^i(E)$ が
$T'$ 上に台をもち、$i \geq m - r'$ であることを意味する。
したがって Lemma \ref{spaces-perfect-lemma-approximation-affine} により、近似
$P \to E|_V$（$(T', E|_V, m)$ の $V$ 上のもの）が存在する。
Lemma \ref{spaces-perfect-lemma-open} を適用すると、$(T', E, m)$ を近似できる。
このような近似は $(T, E, m)$ の近似でもある。

\medskip\noindent
帰納段階。近似 $P \to E|_U$（$(T \cap |U|, E|_U, m)$ のもの）を選ぶ。
これは特に全射 $H^t(P) \to H^t(E|_U)$ を与える。
証明の残りでは、表現可能な代数空間 $V$ および $U \times_X V$ に対して
Lemma \ref{spaces-perfect-lemma-derived-quasi-coherent-small-etale-site} の同値
（および Remark \ref{spaces-perfect-remark-match-total-direct-images} の両立性）を用いる。
また、Zariski サイトと平展サイトにおける $(m - r)$-擬連接性、
それぞれ完全性が一致することも用いる。Lemmas
\ref{spaces-perfect-lemma-descend-pseudo-coherent} および
\ref{spaces-perfect-lemma-descend-perfect} を参照されたい。
したがって、開埋め込み $U \times_X V \subset V$ に対して
Derived Categories of Schemes, Section \ref{perfect-section-lift} の結果を
用いることができる。このようにして Derived Categories of Schemes, Lemma
\ref{perfect-lemma-lift-perfect-complex-plus-shift-support} により、
完全対象 $Q$（$D(\mathcal{O}_V)$ のもの）で $j^{-1}(T)$ 上に台をもつものと、同型
$Q|_{U \times_X V} \to (P \oplus P[1])|_{U \times_X V}$ が存在する。
Derived Categories of Schemes, Lemma \ref{perfect-lemma-lift-map} により、
$Q$ を $Q \otimes^\mathbf{L} I$ で置き換え、次の写像が% P09-ERR-0999
$$
Q|_{U \times_X V} \longrightarrow
(P \oplus P[1])|_{U \times_X V} \longrightarrow
P|_{U \times_X V} \longrightarrow E|_{U \times_X V}
$$
$Q \to E|_V$ に持ち上がると仮定してよい。Lemma \ref{spaces-perfect-lemma-glue} により、
射 $a : R \to E$（$D(\mathcal{O}_X)$ のもの）で、% P09-ERR-1000
$a|_U$ が $P \oplus P[1] \to E|_U$ と同型であり、$a|_V$ が
$Q \to E|_V$ と同型であるものを得る。したがって $R$ は完全で
$T$ 上に台をもち、写像 $H^t(R) \to H^t(E)$ は $U$ への制限上全射である。
識別三角形
$$
R \to E \to E' \to R[1]
$$
を選ぶ。このとき $E'$ は $(m - r)$-擬連接であり
（Cohomology on Sites, Lemma
\ref{sites-cohomology-lemma-cone-pseudo-coherent}）、
$H^i(E')|_U = 0$ は $i \geq t$ に対して成り立ち、
$H^i(E')$ は $T$ 上に台をもつ。ここで $i \geq m - r$ である。
帰納法により、近似 $R' \to E'$（$(T, E', m)$ のもの）を得る。
合成 $R' \to E' \to R[1]$ を識別三角形
$R \to R'' \to R' \to R[1]$ に組み込み、射 $R' \to E'$ および
$R[1] \to R[1]$ を識別三角形の射
$$
\xymatrix{
R \ar[r] \ar[d] & R'' \ar[d] \ar[r] & R' \ar[d] \ar[r] & R[1] \ar[d] \\
R \ar[r] &  E \ar[r] & E' \ar[r] & R[1]
}
$$
へ TR3 を用いて拡張する。このとき $R''$ は完全複体であり
（Cohomology on Sites, Lemma
\ref{sites-cohomology-lemma-two-out-of-three-perfect}）、$T$ 上に台をもつ。
容易な図式追跡により、$R'' \to E$ が望む近似であることが分かる。
\end{proof}

\begin{theorem}
\label{spaces-perfect-theorem-approximation}
$S$ を概型とする。$X$ を $S$ 上の準コンパクトかつ準分離的な代数空間とする。
このとき $X$ 上で完全複体による近似が成り立つ。
\end{theorem}

\begin{proof}
これは Lemma \ref{spaces-perfect-lemma-induction-principle} の帰納原理と Lemmas
\ref{spaces-perfect-lemma-induction-step} および \ref{spaces-perfect-lemma-approximation-affine} から従う。
\end{proof}






\section{導来圏の生成}
\label{spaces-perfect-section-generating}

\noindent
本節は Derived Categories of Schemes, Section
\ref{perfect-section-generating} の類似である。ただし、まず
Derived Categories of Schemes, Lemma
\ref{perfect-lemma-lift-map-from-perfect-complex-with-support} の類似である
次の補題を証明する。

\begin{lemma}
\label{spaces-perfect-lemma-lift-map-from-perfect-complex-with-support}
$S$ を概型とする。$X$ を $S$ 上の準コンパクトかつ準分離的な代数空間とする。
$W \subset X$ を準コンパクト開部分とする。$T \subset |X|$ を閉部分集合で、
$X \setminus T \to X$ が準コンパクト射であるものとする。
$E$ を $D_\QCoh(\mathcal{O}_X)$ の対象とする。
$\alpha : P \to E|_W$ を射とし、$P$ は $D(\mathcal{O}_W)$ の完全対象で
$T \cap W$ 上に台をもつものとする。このとき射 $\beta : R \to E$ で、
$R$ が $D(\mathcal{O}_X)$ の完全対象で $T$ 上に台をもち、$P$ が
$R|_W$ の直和因子（$D(\mathcal{O}_W)$ におけるもの）となり、$\alpha$ と
$\beta|_W$ に両立するものが存在する。
\end{lemma}

\begin{proof}
これを証明するため Lemma \ref{spaces-perfect-lemma-induction-principle-enlarge} の
帰納原理を用いる。したがって直ちに、初等識別正方形
$(W \subset X, f : V \to X)$ があり、$V$ がアフィンで、
$P \to E|_W$ が補題の主張のとおりである場合に帰着する。
証明の残りでは、表現可能な代数空間 $V$ および $W \times_X V$ に対して
Lemma \ref{spaces-perfect-lemma-derived-quasi-coherent-small-etale-site} の同値
（および Remark \ref{spaces-perfect-remark-match-total-direct-images} の両立性）を用いる。
また、Zariski サイトと平展サイトにおける完全性が一致することも用いる。
Lemma \ref{spaces-perfect-lemma-descend-perfect} を参照されたい。

\medskip\noindent
Derived Categories of Schemes, Lemma
\ref{perfect-lemma-lift-perfect-complex-plus-shift-support} により、完全対象
$Q$（$D(\mathcal{O}_V)$ に属するもの）で $f^{-1}T$ 上に台をもつものと、同型
$Q|_{W \times_X V} \to (P \oplus P[1])|_{W \times_X V}$ を選べる。
Derived Categories of Schemes, Lemma \ref{perfect-lemma-lift-map} により、
$Q$ を $Q \otimes^\mathbf{L} I$ で置き換え（なお $f^{-1}T$ 上に台をもつ）、
写像
$$
Q|_{W \times_X V} \to (P \oplus P[1])|_{W \times V}
\longrightarrow P|_{W \times_X V}
\longrightarrow
E|_{W \times_X V}
$$
が $Q \to E|_V$ に持ち上がると仮定してよい。% P09-ERR-1001
Lemma \ref{spaces-perfect-lemma-glue} により、射 $a : R \to E$
（$D(\mathcal{O}_X)$ のもの）で、$a|_W$ が
$P \oplus P[1] \to E|_W$ と同型であり、
$a|_V$ が $Q \to E|_V$ と同型であるものを得る。% P09-ERR-1002
したがって $R$ は望むとおり完全で $T$ 上に台をもつ。
\end{proof}

\begin{remark}
\label{spaces-perfect-remark-addendum}
Lemma \ref{spaces-perfect-lemma-lift-map-from-perfect-complex-with-support} の証明は、
$$
R|_W = P \oplus P^{\oplus n_1}[1] \oplus \ldots \oplus P^{\oplus n_m}[m]
$$
となることを示す。ここで $m \geq 0$ および $n_j \geq 0$ である。
したがって $R|_W$ の最高次数のコホモロジー層は $P$ のそれに等しい。
写像
$P^{\oplus n_1}[1] \oplus \ldots \oplus P^{\oplus n_m}[m] \to R|_W$
に対してこの構成を繰り返し、錐をとり、帰納法を用いることにより、
任意に与えた次数より上で $R|_W$ と $P$ のコホモロジー層を等しくできる。
\end{remark}

\begin{lemma}
\label{spaces-perfect-lemma-direct-summand-of-a-restriction}
$S$ を概型とする。$X$ を $S$ 上の準コンパクトかつ準分離的な代数空間とする。
$W$ を $X$ の準コンパクト開部分空間とする。$P$ を
$D(\mathcal{O}_W)$ の完全対象とする。このとき $P$ は
$D(\mathcal{O}_X)$ のある完全対象の制限の直和因子である。
\end{lemma}

\begin{proof}
Lemma \ref{spaces-perfect-lemma-lift-map-from-perfect-complex-with-support} の特別な場合である。
\end{proof}

\begin{theorem}
\label{spaces-perfect-theorem-bondal-van-den-Bergh}
$S$ を概型とする。$X$ を $S$ 上の準コンパクトかつ準分離的な代数空間とする。
圏 $D_\QCoh(\mathcal{O}_X)$ は一つの完全対象によって生成される。
より正確には、完全対象 $P$（$D(\mathcal{O}_X)$ のもの）で、任意の
$E \in D_\QCoh(\mathcal{O}_X)$ に対して次が同値となるものが存在する。
\begin{enumerate}
\item $E = 0$ である。
\item $\Hom_{D(\mathcal{O}_X)}(P[n], E) = 0$ である。ここで
$n \in \mathbf{Z}$ は任意である。
\end{enumerate}
\end{theorem}

\begin{proof}
Lemma \ref{spaces-perfect-lemma-induction-principle} の帰納原理を用いてこれを証明する。

\medskip\noindent
$X$ がアフィンならば $\mathcal{O}_X$ は完全生成対象である。
これは Lemma \ref{spaces-perfect-lemma-derived-quasi-coherent-small-etale-site} と
Derived Categories of Schemes, Lemma
\ref{perfect-lemma-affine-compare-bounded} から従う。

\medskip\noindent
$(U \subset X, f : V \to X)$ が初等識別正方形で、$U$ は準コンパクト、
定理は $U$ に対して成り立ち、$V$ はアフィン概型であると仮定する。
$P$ を $D(\mathcal{O}_U)$ の完全対象で、
$D_\QCoh(\mathcal{O}_U)$ の生成対象であるものとする。
Lemma \ref{spaces-perfect-lemma-direct-summand-of-a-restriction} を用いて、完全対象
$Q$（$D(\mathcal{O}_X)$ のもの）で、その $U$ への制限が直和であり、直和因子の一つが
$P$ であるものを選べる。$V = \Spec(A)$ と書く。$Z \subset V$ を、
$X \setminus U$ の逆像で、これへ同型に写る被約閉部分概型とする
（Definition \ref{spaces-perfect-definition-elementary-distinguished-square} を参照）。
これは $V$ のレトロコンパクト閉部分集合である。
$f_1, \ldots, f_r \in A$ を $Z = V(f_1, \ldots, f_r)$ となるように選ぶ。
$K \in D(\mathcal{O}_V)$ を、$f_1, \ldots, f_r$ の Koszul 複体
（$A$ 上のもの）に
対応する完全対象とする。$K$ は $Z$ 上に台をもつので、順像
$K' = Rf_*K$ は $D(\mathcal{O}_X)$ の完全対象で、その $V$ への制限は
$K$ であることに注意する（Lemmas \ref{spaces-perfect-lemma-pushforward-perfect} および
\ref{spaces-perfect-lemma-pushforward-with-support-in-open} を参照）。
$Q \oplus K'$ は $D_\QCoh(\mathcal{O}_X)$ の生成対象であると主張する。

\medskip\noindent
$E$ を $D_\QCoh(\mathcal{O}_X)$ の対象で、$Q \oplus K'$ の任意の
シフトから $E$ への非自明な射が存在しないものとする。Lemma
\ref{spaces-perfect-lemma-pushforward-with-support-in-open} により $K' = f_! K$ であり、
したがって
$$
\Hom_{D(\mathcal{O}_X)}(K'[n], E) = \Hom_{D(\mathcal{O}_V)}(K[n], E|_V)
$$
である。Derived Categories of Schemes, Lemma
\ref{perfect-lemma-orthogonal-koszul-complex} により
（Lemma \ref{spaces-perfect-lemma-derived-quasi-coherent-small-etale-site} も用いる）、
これらの群の消滅から $E|_V$ は
$R(U \times_X V \to V)_*E|_{U \times_X V}$ と同型であることが従う。
これは $E = R(U \to X)_*E|_U$ を含意する（小さな詳細は省略する）。
この場合
$$
\Hom_{D(\mathcal{O}_X)}(Q[n], E) = \Hom_{D(\mathcal{O}_U)}(Q|_U[n], E|_U)
$$
であり、これは $\Hom_{D(\mathcal{O}_U)}(P[n], E|_U)$ を直和因子として含む。
したがって $P$ の選び方により、これらの群の消滅から $E|_U$ が零であることが
従う。よって $E$ は零である。
\end{proof}

\begin{remark}
\label{spaces-perfect-remark-pullback-generator}
$S$ を概型とする。$f : X \to Y$ を $S$ 上の準コンパクトかつ準分離的な
代数空間の射とする。$E \in D_\QCoh(\mathcal{O}_Y)$ を生成対象とする
（Theorem \ref{spaces-perfect-theorem-bondal-van-den-Bergh} を参照）。
このとき次の条件は同値である。
\begin{enumerate}
\item $K \in D_\QCoh(\mathcal{O}_X)$ に対して、$Rf_*K = 0$ であることと
$K = 0$ であることとは同値である。
\item $Rf_* : D_\QCoh(\mathcal{O}_X) \to D_\QCoh(\mathcal{O}_Y)$ は
同型を反映する。
\item $Lf^*E$ は $D_\QCoh(\mathcal{O}_X)$ の生成対象である。
\end{enumerate}
(1) と (2) の同値は、
$Rf_* : D_\QCoh(\mathcal{O}_X) \to D_\QCoh(\mathcal{O}_Y)$ が
三角圏の完全函手であることの形式的帰結である。同様に、(1) と (3) の同値は
$Lf^*$ が $Rf_*$ の左随伴であることから形式的に従う。
これらの条件は、$f$ がアフィンである場合
（Lemma \ref{spaces-perfect-lemma-affine-morphism}）、$f$ が開埋め込みである場合、または
$f$ がそのような射の合成である場合に成り立つ。
\end{remark}

\noindent
次の結果は Theorem \ref{spaces-perfect-theorem-bondal-van-den-Bergh} の強化であり、
% P09-ERR-1003
まったく同じ方法で証明される。$T \subset |X|$ を閉部分集合とし、
ここで $X$ は代数空間とする。$D_T(\mathcal{O}_X)$ により、
コホモロジー層が $T$ 上に台をもつ複体からなる真に充満な飽和三角部分圏を
表すことにする。

\begin{lemma}
\label{spaces-perfect-lemma-generator-with-support}
$S$ を概型とする。$X$ を $S$ 上の準コンパクトかつ準分離的な代数空間とする。
$T \subset |X|$ を閉部分集合で、$|X| \setminus T$ が準コンパクトであるものとする。
上の記法のもとで、圏 $D_{\QCoh, T}(\mathcal{O}_X)$ は一つの完全対象によって
生成される。
\end{lemma}

\begin{proof}
Lemma \ref{spaces-perfect-lemma-induction-principle} の帰納原理を用いてこれを証明する。
この性質は、サイト $X_{spaces, \etale}$ の表現可能な準コンパクトかつ
準分離的対象に対して、Derived Categories of Schemes, Lemma
\ref{perfect-lemma-generator-with-support} により成り立つ。

\medskip\noindent
$(U \subset X, f : V \to X)$ が初等識別正方形で、補題は $U$ に対して
成り立ち、$V$ はアフィンであると仮定する。証明を終えるには結果が $X$ に
対して成り立つことを示す必要がある。$P$ を $D(\mathcal{O}_U)$ の完全対象で
$T \cap U$ 上に台をもち、$D_{\QCoh, T \cap U}(\mathcal{O}_U)$ の
生成対象であるものとする。Lemma
\ref{spaces-perfect-lemma-lift-map-from-perfect-complex-with-support} を用いると、
完全対象 $Q$（$D(\mathcal{O}_X)$ のもの）で $T$ 上に台をもち、その $U$ への制限が
直和で、直和因子の一つが $P$ であるものを選べる。
$V = \Spec(B)$ と書く。$Z = X \setminus U$ と置く。このとき
$f^{-1}Z$ は $V$ の閉部分集合で、$V \setminus f^{-1}Z$ は準コンパクトである。
$X$ は準分離的なので、
$f^{-1}Z \cap f^{-1}T = f^{-1}(Z \cap T)$ は $V$ の閉部分集合で、
$W = V \setminus f^{-1}(Z \cap T)$ は準コンパクトである。
したがって $g_1, \ldots, g_s \in B$ を
$f^{-1}(Z \cap T) = V(g_1, \ldots, g_r)$ となるように選べる。
% P09-ERR-1004
$K \in D(\mathcal{O}_V)$ を、$g_1, \ldots, g_s$ の Koszul 複体
（$B$ 上のもの）に
対応する完全対象とする。$K$ は閉部分集合
$f^{-1}(Z \cap T) \subset V$ 上に台をもつので、順像
$K' = R(V \to X)_*K$ は $D(\mathcal{O}_X)$ の完全対象で、その $V$ への
制限は $K$ であることに注意する（Lemmas
\ref{spaces-perfect-lemma-pushforward-perfect} および
\ref{spaces-perfect-lemma-pushforward-with-support-in-open} を参照）。
$Q \oplus K'$ は $D_{\QCoh, T}(\mathcal{O}_X)$ の生成対象であると主張する。

\medskip\noindent
$E$ を $D_{\QCoh, T}(\mathcal{O}_X)$ の対象で、$Q \oplus K'$ の任意の
シフトから $E$ への非自明な射が存在しないものとする。Lemma
\ref{spaces-perfect-lemma-pushforward-with-support-in-open} により
$K' = R(V \to X)_! K$ であり、したがって
$$
\Hom_{D(\mathcal{O}_X)}(K'[n], E) = \Hom_{D(\mathcal{O}_V)}(K[n], E|_V)
$$
である。Derived Categories of Schemes, Lemma
\ref{perfect-lemma-orthogonal-koszul-complex} により、
$E|_V = Rj_*E|_W$ である。ここで $j : W \to V$ は包含である。図式は
$$
\xymatrix{
W \ar[r]_j & V & Z \cap T \ar[l] \ar[d] \\
V \setminus f^{-1}Z \ar[u]^{j'} \ar[ru]_{j''} & & Z \ar[lu]
}
$$
である。% P09-ERR-1005
$E$ は $T$ 上に台をもつので、$E|_W$ は
$f^{-1}T \cap W = f^{-1}T \cap (V \setminus f^{-1}Z)$ 上に台をもち、
これは $W$ で閉である。したがって
$$
E|_V = Rj_*(E|_W) = Rj_*(Rj'_*(E|_{U \cap V})) = Rj''_*(E|_{U \cap V})
$$
を得る。% P09-ERR-1006
ここで第二の等式は Cohomology, Lemma
\ref{cohomology-lemma-pushforward-restriction} の (1) である。
これを適用できるのは、$V$ が概型で、$E$ が準連接コホモロジー層をもち、
したがって準コンパクト開埋め込み $j'$ に沿う順像が台となる概型上の順像と
一致するからである。Remark \ref{spaces-perfect-remark-match-total-direct-images} を
参照されたい。これは $E = R(U \to X)_*E|_U$ を含意する
（小さな詳細は省略する）。この場合
$$
\Hom_{D(\mathcal{O}_X)}(Q[n], E) = \Hom_{D(\mathcal{O}_U)}(Q|_U[n], E|_U)
$$
であり、これは $\Hom_{D(\mathcal{O}_U)}(P[n], E|_U)$ を直和因子として含む。
したがって $P$ の選び方により、これらの群の消滅から $E|_U$ が零であることが
従う。よって $E$ は零である。
\end{proof}










\section{コンパクト対象と完全対象}
\label{spaces-perfect-section-compact}

\noindent
本節は Derived Categories of Schemes, Section
\ref{perfect-section-compact} の類似である。

\begin{proposition}
\label{spaces-perfect-proposition-compact-is-perfect}
$S$ を概型とする。$X$ を $S$ 上の準コンパクトかつ準分離的な代数空間とする。
$D_\QCoh(\mathcal{O}_X)$ の対象がコンパクトであることと完全であることとは
同値である。
\end{proposition}

\begin{proof}
$K$ を $D(\mathcal{O}_X)$ の完全対象とし、その双対を $K^\vee$ とする
（Cohomology on Sites, Lemma
\ref{sites-cohomology-lemma-dual-perfect-complex}）。このとき
$$
\Hom_{D(\mathcal{O}_X)}(K, M) =
H^0(X, K^\vee \otimes_{\mathcal{O}_X}^\mathbf{L} M)
$$
が $M$ に函手的に成り立つ。$K^\vee \otimes_{\mathcal{O}_X}^\mathbf{L} -$ は直和と可換し、
Lemma \ref{spaces-perfect-lemma-quasi-coherence-pushforward-direct-sums} により
$H^0(X, -)$ は $D_\QCoh(\mathcal{O}_X)$ 上で直和と可換する。
したがって $K$ は $D_\QCoh(\mathcal{O}_X)$ においてコンパクトである。

\medskip\noindent
逆に、$K$ を $D_\QCoh(\mathcal{O}_X)$ のコンパクト対象とする。
$K$ が完全であることを示すには、$K|_U$ が完全であることを、任意の
アフィン概型 $U$ が $X$ 上平展である場合に示せば十分である。Cohomology on Sites,
Lemma \ref{sites-cohomology-lemma-perfect-independent-representative} を
参照されたい。$j : U \to X$ は準コンパクトかつ分離的な射であることに
注意する。したがって Lemma
\ref{spaces-perfect-lemma-quasi-coherence-pushforward-direct-sums} により
$$
Rj_* : D_\QCoh(\mathcal{O}_U) \to D_\QCoh(\mathcal{O}_X)
$$
は直和と可換する。ゆえに $U$ への制限と $Rj_*$ の随伴性から、$K|_U$ は
$D_\QCoh(\mathcal{O}_U)$ の完全対象であることが従う。% P09-ERR-1007
したがって $X$ がアフィン、特に準コンパクトかつ準分離的な概型である場合に
帰着する。Lemmas \ref{spaces-perfect-lemma-derived-quasi-coherent-small-etale-site} および
\ref{spaces-perfect-lemma-descend-perfect} を介して概型の場合、すなわち
Derived Categories of Schemes, Proposition
\ref{perfect-proposition-compact-is-perfect} に帰着する。
\end{proof}

\begin{remark}
\label{spaces-perfect-remark-classical-generator}
$S$ を概型とする。$X$ を $S$ 上の準コンパクトかつ準分離的な代数空間とする。
$G$ を $D(\mathcal{O}_X)$ の完全対象で、$D_\QCoh(\mathcal{O}_X)$ の
生成対象であるものとする。Theorem \ref{spaces-perfect-theorem-bondal-van-den-Bergh} により、
このようなものが少なくとも一つ存在する。Lemma
\ref{spaces-perfect-lemma-quasi-coherence-direct-sums}、Proposition
\ref{spaces-perfect-proposition-compact-is-perfect}、および Derived Categories, Proposition
\ref{derived-proposition-generator-versus-classical-generator} を合わせると、
$G$ は $D_{perf}(\mathcal{O}_X)$ の古典的生成対象であることが分かる。
\end{remark}

\noindent
次の結果は Proposition \ref{spaces-perfect-proposition-compact-is-perfect} の強化である。
$T \subset |X|$ を閉部分集合とし、ここで $X$ は代数空間とする。
前と同様、$D_T(\mathcal{O}_X)$ は、コホモロジー層が $T$ 上に台をもつ
複体からなる真に充満な飽和三角部分圏を表す。直和をとることは
コホモロジー層をとることと可換するので、$D_T(\mathcal{O}_X)$ は直和をもち、
それらは $D(\mathcal{O}_X)$ における直和と等しい。

\begin{lemma}
\label{spaces-perfect-lemma-compact-is-perfect-with-support}
$S$ を概型とする。$X$ を $S$ 上の準コンパクトかつ準分離的な代数空間とする。
$T \subset |X|$ を閉部分集合で、$|X| \setminus T$ が準コンパクトであるものとする。
$D_{\QCoh, T}(\mathcal{O}_X)$ の対象がコンパクトであることと、
$D(\mathcal{O}_X)$ の対象として完全であることとは同値である。
\end{lemma}

\begin{proof}
補題の直前の注意により、$D_{\QCoh, T}(\mathcal{O}_X)$ は直和をもつ三角圏である。
Proposition \ref{spaces-perfect-proposition-compact-is-perfect} により、完全対象は
$D(\mathcal{O}_X)$ のコンパクト対象を定め、したがって直和をとることで
保たれる任意の部分圏でもコンパクト対象を定める。% P09-ERR-1008
逆向きには、生成対象 $E \in D_{\QCoh, T}(\mathcal{O}_X)$ で、完全
$\mathcal{O}_X$-加群複体でもあるものが存在することを用いる。% P09-ERR-1009
Lemma \ref{spaces-perfect-lemma-generator-with-support} を参照されたい。
したがって上により $E$ はコンパクトである。Derived Categories, Proposition
\ref{derived-proposition-generator-versus-classical-generator} により、$E$ は
$D_{\QCoh, T}(\mathcal{O}_X)$ のコンパクト対象からなる充満部分圏の
古典的生成対象である。したがって任意のコンパクト対象は、$E$ から、
(a) シフト、(b) 有限直和、(c) 錐、(d) 直和因子をとる操作からなる有限列に
よって構成できる。これらの各操作は完全であるという性質を保つので、結果が従う。
\end{proof}

\begin{remark}
\label{spaces-perfect-remark-classical-generator-with-support}
$S$ を概型とする。$X$ を $S$ 上の準コンパクトかつ準分離的な代数空間とする。
$T \subset |X|$ を閉部分集合で、$|X| \setminus T$ が準コンパクトであるものとする。
$G$ を $D_{\QCoh, T}(\mathcal{O}_X)$ の完全対象で、
$D_{\QCoh, T}(\mathcal{O}_X)$ の生成対象であるものとする。Lemma
\ref{spaces-perfect-lemma-generator-with-support} により、このようなものが少なくとも一つ存在する。
$D_{\QCoh, T}(\mathcal{O}_X)$ が直和をもつという事実と、Lemma
\ref{spaces-perfect-lemma-compact-is-perfect-with-support} および Derived Categories, Proposition
\ref{derived-proposition-generator-versus-classical-generator} を合わせると、
$G$ は $D_{perf, T}(\mathcal{O}_X)$ の古典的生成対象であることが分かる。
\end{remark}
\noindent
次の補題は、直前の補題の証明に用いられた考え方の一応用である。

\begin{lemma}
\label{spaces-perfect-lemma-map-from-pseudo-coherent-to-complex-with-support}
$S$ を概型とする。$X$ を $S$ 上の準コンパクトかつ準分離的な代数空間とする。
$T \subset |X|$ を閉部分集合とし、その補集合 $U \subset X$ は準コンパクトで
あるとする。$\alpha : P \to E$ を $D_\QCoh(\mathcal{O}_X)$ の射で、次の
いずれかを満たすものとする。
\begin{enumerate}
\item $P$ は完全であり、$E$ は $T$ 上に台をもつ。
\item $P$ は擬連接であり、$E$ は $T$ 上に台をもち、かつ $E$ は下に有界である。
\end{enumerate}
このとき、$\mathcal{O}_X$-加群の完全複体 $I$ と射
$I \to \mathcal{O}_X[0]$ であって、$I \otimes^\mathbf{L} P \to E$ が零であり、
$I|_U \to \mathcal{O}_U[0]$ が同型となるものが存在する。
\end{lemma}

\begin{proof}
$\mathcal{D} = D_{\QCoh, T}(\mathcal{O}_X)$ と置く。いずれの場合にも、複体
$K = R\SheafHom(P, E)$ は $\mathcal{D}$ の対象である。準連接性については
Lemma \ref{spaces-perfect-lemma-quasi-coherence-internal-hom} を参照されたい。
$K$ が $T$ 上に台をもつことは明らかである。実際、$R\SheafHom$ の形成は
開部分への制限と可換する。射 $\alpha$ は
$H^0(K) = \Hom_{D(\mathcal{O}_X)}(\mathcal{O}_X[0], K)$ の元を定める。
したがって射 $\alpha : \mathcal{O}_X[0] \to K$ に対して結果を証明すれば
十分である。

\medskip\noindent
$E \in \mathcal{D}$ を完全生成対象とする。Lemma
\ref{spaces-perfect-lemma-generator-with-support} を参照されたい。Derived Categories, Lemma
\ref{derived-lemma-write-as-colimit} のように
$$
K = \text{hocolim} K_n
$$
と書く。ここでは生成対象 $E$ を用いる。函手
$\mathcal{D} \to D(\mathcal{O}_X)$ は直和と可換するので、
$K = \text{hocolim} K_n$ は $D(\mathcal{O}_X)$ においても成り立つ。
$\mathcal{O}_X$ は $D(\mathcal{O}_X)$ のコンパクト対象であるから、ある $n$ と射
$\alpha_n : \mathcal{O}_X \to K_n$ が存在し、これが $\alpha$ を誘導する。
Derived Categories, Lemma
\ref{derived-lemma-commutes-with-countable-sums} を参照されたい。
射 $\mathcal{O}_X[0] \to K_n$（周囲の圏 $D(\mathcal{O}_X)$ におけるもの）に
Derived Categories, Lemma
\ref{derived-lemma-factor-through} を適用すると、$\alpha_n$ は
$\mathcal{O}_X[0] \to Q \to K_n$ と分解し、ここで $Q$ は
$\langle E \rangle$ の対象である。したがって $Q$ は $T$ 上に台をもつ
完全複体である。

\medskip\noindent
識別三角
$$
I \to \mathcal{O}_X[0] \to Q \to I[1]
$$
を選ぶ。構成により $I$ は完全であり、射 $I \to \mathcal{O}_X[0]$ は
$U$ 上で同型に制限され、合成 $I \to K$ は零である。実際、$\alpha$ は
$Q$ を経由する。これで補題が証明された。
\end{proof}










\section{加群圏としての導来圏}
\label{spaces-perfect-section-derived-is-dga}

\noindent
本節は Derived Categories of Schemes, Section
\ref{perfect-section-derived-is-dga} の類似である。

\begin{lemma}
\label{spaces-perfect-lemma-tensor-with-QCoh-complex}
$S$ を概型とする。$X$ を $S$ 上の代数空間とする。$K^\bullet$ を、
コホモロジー層が準連接である $\mathcal{O}_X$-加群の複体とする。
自己準同型微分次数付き代数を
$(E, d) = \Hom_{\text{Comp}^{dg}(\mathcal{O}_X)}(K^\bullet, K^\bullet)$
と置く。このとき Differential Graded Algebra, Lemma
\ref{dga-lemma-tensor-with-complex-derived} の函手
$$
- \otimes_E^\mathbf{L} K^\bullet :
D(E, \text{d}) \longrightarrow D(\mathcal{O}_X)
$$
の像は $D_\QCoh(\mathcal{O}_X)$ に含まれる。
\end{lemma}

\begin{proof}
$P$ を微分次数付き $E$-加群で性質 $P$ をもつものとする。$F_\bullet$ を、
Differential Graded Algebra, Section \ref{dga-section-P-resolutions} のような
$P$ 上のフィルトレーションとする。このとき
$$
P \otimes_E K^\bullet = \text{hocolim}\ F_iP \otimes_E K^\bullet
$$
である。各 $F_iP$ は有限フィルトレーションをもち、その次数付き成分は
$E[k]$ の直和である。結果は容易に従う。
\end{proof}

\noindent
次の補題は強めることができる（コホモロジー層が固定された範囲内でのみ
非零であるすべての $L$ にわたって、消滅を一様に選べる）。

\begin{lemma}
\label{spaces-perfect-lemma-ext-from-perfect-into-bounded-QCoh}
$S$ を概型とする。$X$ を $S$ 上の準コンパクトかつ準分離的な代数空間とする。
$K$ を $D(\mathcal{O}_X)$ の完全対象とする。このとき
\begin{enumerate}
\item ある整数 $a \leq b$ が存在し、
$\Hom_{D(\mathcal{O}_X)}(K, L) = 0$ が、$L \in D_\QCoh(\mathcal{O}_X)$ であって
$H^i(L) = 0$ を $i \in [a, b]$ に対して満たすものについて成り立つ。
\item $L$ が有界ならば、$\Ext^n_{D(\mathcal{O}_X)}(K, L)$ は有限個を除く
すべての $n$ に対して零である。
\end{enumerate}
\end{lemma}

\begin{proof}
(2) は (1) と
$\Ext^n_{D(\mathcal{O}_X)}(K, L) =
\Hom_{D(\mathcal{O}_X)}(K, L[n])$
から従う。(1) を証明する。$K$ は完全であるから
$$
\Ext^i_{D(\mathcal{O}_X)}(K, L) =
H^i(X, K^\vee \otimes_{\mathcal{O}_X}^\mathbf{L} L)
$$
である。ここで $K^\vee$ は $K$ の「双対」完全複体である。Cohomology on Sites,
Lemma \ref{sites-cohomology-lemma-dual-perfect-complex} を参照されたい。
$P = K^\vee \otimes_{\mathcal{O}_X}^\mathbf{L} L$ は
$D_\QCoh(X)$ に属する。これは Lemmas
\ref{spaces-perfect-lemma-quasi-coherence-tensor-product} および
\ref{spaces-perfect-lemma-pseudo-coherent} による（後者は完全複体が準連接コホモロジー層を
もつことを見るために用いる）。$K^\vee$ の Tor 振幅が $[a, b]$ に含まれるとする。
このときスペクトル系列
$$
E_1^{p, q} = H^p(K^\vee \otimes_{\mathcal{O}_X}^\mathbf{L} H^q(L))
\Rightarrow
H^{p + q}(K^\vee \otimes_{\mathcal{O}_X}^\mathbf{L} L)
$$
から、$H^j(K^\vee \otimes_{\mathcal{O}_X}^\mathbf{L} L)$ は、
$H^q(L) = 0$ が $q \in [j - b, j - a]$ に対して成り立つならば
零であることが分かる。
Cohomology of Spaces, Lemma
\ref{spaces-cohomology-lemma-vanishing-quasi-separated} の整数 $N$、すなわち
$\max(d_p + p)$ をとる。このとき
$H^0(X, K^\vee \otimes_{\mathcal{O}_X}^\mathbf{L} L)$ は、コホモロジー層
$$
H^{-N}(K^\vee \otimes_{\mathcal{O}_X}^\mathbf{L} L),
\ H^{-N + 1}(K^\vee \otimes_{\mathcal{O}_X}^\mathbf{L} L),
\ \ldots,
\ H^0(K^\vee \otimes_{\mathcal{O}_X}^\mathbf{L} L)
$$
が零ならば消滅する。実際、引用した補題と Lemma
\ref{spaces-perfect-lemma-application-nice-K-injective} により
$$
H^0(X, K^\vee \otimes_{\mathcal{O}_X}^\mathbf{L} L) =
H^0(X, \tau_{\geq -N}(K^\vee \otimes_{\mathcal{O}_X}^\mathbf{L} L))
$$
であり、コホモロジー層の消滅により、これは
$H^0(X, \tau_{\geq 1}(K^\vee \otimes_{\mathcal{O}_X}^\mathbf{L} L))$
に等しく、さらに Derived Categories, Lemma
\ref{derived-lemma-negative-vanishing} により零である。
したがって $\Hom_{D(\mathcal{O}_X)}(K, L)$ は、$H^i(L) = 0$ が
$i \in [-b - N, -a]$ に対して成り立つならば零である。
\end{proof}

\noindent
次は Derived Categories of Schemes, Theorem
\ref{perfect-theorem-DQCoh-is-Ddga} の類似である。

\begin{theorem}
\label{spaces-perfect-theorem-DQCoh-is-Ddga}
$S$ を概型とする。$X$ を $S$ 上の準コンパクトかつ準分離的な代数空間とする。
微分次数付き代数 $(E, \text{d})$ で、有限個を除いてコホモロジー群
$H^i(E)$ が零であるものが存在し、$D_\QCoh(\mathcal{O}_X)$ は
$D(E, \text{d})$ と同値になる。
\end{theorem}

\begin{proof}
$K^\bullet$ を $\mathcal{O}$-加群の K-入射的複体で、完全であり、
$D_\QCoh(\mathcal{O}_X)$ を生成するものとする。そのようなものは Theorem
\ref{spaces-perfect-theorem-bondal-van-den-Bergh} と K-入射的分解の存在により存在する。
次のものについて定理が成り立つことを示す。
$$
(E, \text{d}) = \Hom_{\text{Comp}^{dg}(\mathcal{O}_X)}(K^\bullet, K^\bullet)
$$
ここで $\text{Comp}^{dg}(\mathcal{O}_X)$ は $\mathcal{O}$-加群複体の
微分次数付き圏である。Differential Graded Algebra, Section
\ref{dga-section-variant-base-change} を参照されたい。
$K^\bullet$ は K-入射的であるから、すべての $n \in \mathbf{Z}$ に対して
\begin{equation}
\label{spaces-perfect-equation-E-is-OK}
H^n(E) = \Ext^n_{D(\mathcal{O}_X)}(K^\bullet, K^\bullet)
\end{equation}
である。Lemma \ref{spaces-perfect-lemma-ext-from-perfect-into-bounded-QCoh} により、
これらの Ext のうち非零なのは有限個だけである。Differential Graded Algebra,
Lemma \ref{dga-lemma-tensor-with-complex-derived} の函手
$$
- \otimes_E^\mathbf{L} K^\bullet :
D(E, \text{d}) \longrightarrow D(\mathcal{O}_X)
$$
を考える。$K^\bullet$ は完全なので、$D(\mathcal{O}_X)$ のコンパクト対象を
定める。Proposition \ref{spaces-perfect-proposition-compact-is-perfect} を参照されたい。
% P09-ERR-1010
これと (\ref{spaces-perfect-equation-E-is-OK}) を合わせると、上の函手は忠実充満である。
これは Differential Graded Algebra, Lemmas
\ref{dga-lemma-fully-faithful-in-compact-case} から従う。また
Differential Graded Algebra, Lemmas
\ref{dga-lemma-tensor-with-complex-hom-adjoint} により、右随伴
$$
R\Hom(K^\bullet, - ) : D(\mathcal{O}_X) \longrightarrow D(E, \text{d})
$$
をもち、随伴函手に関する一般論により、これは左準逆函手である。
一方、Lemma \ref{spaces-perfect-lemma-tensor-with-QCoh-complex} から
$$
- \otimes_E^\mathbf{L} K^\bullet :
D(E, \text{d}) \longrightarrow D_\QCoh(\mathcal{O}_X)
$$
を得る。$K^\bullet$ を $D_\QCoh(\mathcal{O}_X)$ の生成対象として選んだので、
$D_\QCoh(\mathcal{O}_X)$ に制限した随伴函手の核は零である。形式的な議論から
所望の同値が得られる。Derived Categories, Lemma
\ref{derived-lemma-fully-faithful-adjoint-kernel-zero} を参照されたい。
\end{proof}

\begin{remark}[台付きの場合]
\label{spaces-perfect-remark-DQCoh-is-Ddga-with-support}
$S$ を概型とする。$X$ を準コンパクトかつ準分離的な代数空間とする。
$T \subset |X|$ を閉部分集合で、$|X| \setminus T$ が準コンパクトであるものとする。
Theorem \ref{spaces-perfect-theorem-DQCoh-is-Ddga} の類似が
$D_{\QCoh, T}(\mathcal{O}_X)$ に対して成り立つ。これは定理の証明とまったく
同じ議論から従う。その際、Lemmas \ref{spaces-perfect-lemma-generator-with-support}、
\ref{spaces-perfect-lemma-compact-is-perfect-with-support}、および台付きの
Lemma \ref{spaces-perfect-lemma-tensor-with-QCoh-complex} の変種を用いる。
これが必要になれば、ここに結果を正確に述べ、詳細な証明を与える。
\end{remark}

\begin{remark}[dga の一意性]
\label{spaces-perfect-remark-independence-choice}
$X$ を環 $R$ 上の準コンパクトかつ準分離的な代数空間とする。
Theorem \ref{spaces-perfect-theorem-DQCoh-is-Ddga} の証明の構成により、微分次数付き代数
$(A, \text{d})$ で $R$ 上のものが存在し、三角圏として
$D_\QCoh(X)$ は $R$-線型に $D(A, \text{d})$ と同値である。
$(A, \text{d})$ はどの程度一意なのか、と問うことができる。
答えは、$(A, \text{d})$ が導来同値を除いて定まると言うだけよりも
（ほんの）少し良い。実際、$(B, \text{d})$ を第二のそのような対と仮定する。
このとき
$$
(A, \text{d}) = \Hom_{\text{Comp}^{dg}(\mathcal{O}_X)}(K^\bullet, K^\bullet)
$$
および
$$
(B, \text{d}) = \Hom_{\text{Comp}^{dg}(\mathcal{O}_X)}(L^\bullet, L^\bullet)
$$
である。ここで $K^\bullet$ と $L^\bullet$ は、$\mathcal{O}_X$-加群の
K-入射的複体で、$D_\QCoh(\mathcal{O}_X)$ の完全生成対象に対応する。次のように置く。
$$
\Omega = \Hom_{\text{Comp}^{dg}(\mathcal{O}_X)}(K^\bullet, L^\bullet)
\quad
\Omega' = \Hom_{\text{Comp}^{dg}(\mathcal{O}_X)}(L^\bullet, K^\bullet)
$$
このとき $\Omega$ は微分次数付き $B^{opp} \otimes_R A$-加群であり、
$\Omega'$ は微分次数付き $A^{opp} \otimes_R B$-加群である。さらに、同値
$$
D(A, \text{d}) \to D_\QCoh(\mathcal{O}_X) \to
D(B, \text{d})
$$
は函手 $- \otimes_A^\mathbf{L} \Omega'$ によって与えられ、準逆についても
同様である。したがって Differential Graded Algebra, Remark
\ref{dga-remark-hochschild-cohomology} の状況にある。この注意が必要になれば、
ここに正確な主張と詳細な証明を与える。
\end{remark}











\section{擬連接複体の特徴付け、I}
\label{spaces-perfect-section-pseudo-coherent-hocolim}

\noindent
この内容は More on Morphisms of Spaces, Section
\ref{spaces-more-morphisms-section-characterize-pseudo-coherent} で続ける。
擬連接対象は、完全対象による近似の導来ホモトピー極限として特徴付けることができる。
% P09-ERR-1012

\begin{lemma}
\label{spaces-perfect-lemma-pseudo-coherent-hocolim}
$S$ を概型とする。$X$ を $S$ 上の準コンパクトかつ準分離的な代数空間とする。
$K \in D(\mathcal{O}_X)$ とする。次は同値である。
\begin{enumerate}
\item $K$ は擬連接である。
\item $K = \text{hocolim} K_n$ であり、ここで $K_n$ は完全で、
$\tau_{\geq -n}K_n \to \tau_{\geq -n}K$ はすべての $n$ に対して同型である。
\end{enumerate}
\end{lemma}

\begin{proof}
(2) $\Rightarrow$ (1) は任意の環付きサイト上で成り立つ。実際、(2) を仮定する。
導来圏の完全対象は擬連接であることを思い出そう。Cohomology on Sites, Lemma
\ref{sites-cohomology-lemma-perfect} を参照されたい。すると定義から、
$\tau_{\geq -n}K_n$ は $(-n + 1)$-擬連接であり、したがって
$\tau_{\geq -n}K$ は $(-n + 1)$-擬連接であり、ゆえに $K$ は
$(-n + 1)$-擬連接である。これはすべての $n$ に対して成り立つので、
$K$ は擬連接である。Cohomology on Sites, Definition
\ref{sites-cohomology-definition-pseudo-coherent} を参照されたい。

\medskip\noindent
(1) を仮定する。まず近似 $K_1 \to K$ を選ぶ。これは $(X, K, -2)$ の
完全複体 $K_1$ による近似である。Definitions \ref{spaces-perfect-definition-approximation-holds}、
\ref{spaces-perfect-definition-approximation}、および Theorem
\ref{spaces-perfect-theorem-approximation} を参照されたい。帰納的に
$$
K_1 \to K_2 \to \ldots \to K_n \to K
$$
が与えられ、$K_i$ は完全であり、
$\tau_{\geq -i}K_i \to \tau_{\geq -i}K$ はすべての $1 \leq i \leq n$ に
対して同型であると仮定する。
% P09-ERR-1011
次に $a \leq b$ を選ぶ。これは Lemma
\ref{spaces-perfect-lemma-ext-from-perfect-into-bounded-QCoh} における完全対象 $K_n$ に
対する整数である。
近似 $K_{n + 1} \to K$ を選ぶ。これは
$(X, K, \min(a - 1, -n - 1))$ の近似である。識別三角
$$
K_{n + 1} \to K \to C \to K_{n + 1}[1]
$$
を選ぶ。このとき $C \in D_\QCoh(\mathcal{O}_X)$ であり、
$H^i(C) = 0$ が $i \geq a$ に対して成り立つ。したがって $a, b$ の選び方により
$\Hom_{D(\mathcal{O}_X)}(K_n, C) = 0$ である。ゆえに合成
$K_n \to K \to C$ は零である。したがって Derived Categories, Lemma
\ref{derived-lemma-representable-homological} により、$K_n \to K$ を
$K_{n + 1}$ を通じて分解でき、帰納段階が証明される。

\medskip\noindent
なお $K = \text{hocolim} K_n$ を証明しなければならない。これは函手
$H^i( - ) : D(\mathcal{O}_X) \to \textit{Mod}(\mathcal{O}_X)$ と
$K_n$ の選び方に Derived Categories, Lemma
\ref{derived-lemma-cohomology-of-hocolim} を適用することで従う。
\end{proof}

\begin{lemma}
\label{spaces-perfect-lemma-pseudo-coherent-hocolim-with-support}
$X$ を準コンパクトかつ準分離的な概型とする。$T \subset X$ を閉部分集合で、
$X \setminus T$ が準コンパクトであるものとする。
$K \in D(\mathcal{O}_X)$ は $T$ 上に台をもつとする。次は同値である。
\begin{enumerate}
\item $K$ は擬連接である。
\item $K = \text{hocolim} K_n$ であり、ここで $K_n$ は完全で $T$ 上に台をもち、
$\tau_{\geq -n}K_n \to \tau_{\geq -n}K$ はすべての $n$ に対して同型である。
\end{enumerate}
\end{lemma}

\begin{proof}
この補題の証明は Lemma \ref{spaces-perfect-lemma-pseudo-coherent-hocolim} の証明と
まったく同じである。ただし近似を選ぶ際に三つ組 $(T, K, m)$ を用いる。
\end{proof}












\section{コヒーレータ再考}
\label{spaces-perfect-section-better-coherator}

\noindent
Section \ref{spaces-perfect-section-coherator} では、右随伴 $RQ_X$ を構成し、調べた。
これは標準函手
$$
D(\QCoh(\mathcal{O}_X)) \to D(\mathcal{O}_X)
$$
に対するものである。これはコヒーレータ
$Q_X : \textit{Mod}(\mathcal{O}_X) \to \QCoh(\mathcal{O}_X)$ の
右導来拡張として構成された。本節では包含函手
$$
D_\QCoh(\mathcal{O}_X) \longrightarrow D(\mathcal{O}_X)
$$
がいつ右随伴をもつかを調べる。この右随伴が存在するならば、
これを\footnote{これはおそらく標準的な記法ではない。しかし、すでに
コヒーレータには $Q_X$、その導来拡張には $RQ_X$ を用いた。}
$$
DQ_X :
D(\mathcal{O}_X) \longrightarrow D_\QCoh(\mathcal{O}_X)
$$
と記す。準コンパクトかつ準分離的な代数空間には、このような右随伴が
存在することが分かる。

\begin{lemma}
\label{spaces-perfect-lemma-better-coherator}
$S$ を概型とする。$X$ を $S$ 上の準コンパクトかつ準分離的な代数空間とする。
包含函手 $D_\QCoh(\mathcal{O}_X) \to D(\mathcal{O}_X)$ は右随伴をもつ。
\end{lemma}

\begin{proof}[第一の証明]
Lemma \ref{spaces-perfect-lemma-induction-principle} の帰納原理を用いて証明する。
$D(\QCoh(\mathcal{O}_X)) \to D_\QCoh(\mathcal{O}_X)$ が同値ならば、
Section \ref{spaces-perfect-section-coherator} の函手 $RQ_X$ が
$D(\QCoh(\mathcal{O}_X)) \to D(\mathcal{O}_X)$ の右随伴なので、
補題は成り立つ。特に、補題はアフィン代数空間に対して成り立つ。
Lemma \ref{spaces-perfect-lemma-affine-coherator} を参照されたい。したがって次を示せば十分である。
$(U \subset X, f : V \to X)$ が初等識別正方形で、$U$ は準コンパクト、
$V$ はアフィンであり、補題が $U$、$V$、および $U \times_X V$ に対して
成り立つならば、補題は $X$ に対して成り立つ。

\medskip\noindent
右随伴が存在することと、任意の対象 $K$（$D(\mathcal{O}_X)$ のもの）に対して
識別三角
% P09-ERR-1013
$$
E' \to E \to K \to E'[1]
$$
を $D(\mathcal{O}_X)$ において見いだせて、$E'$ が
$D_\QCoh(\mathcal{O}_X)$ に属し、かつ $\Hom(M, K) = 0$ がすべての
$M$（$D_\QCoh(\mathcal{O}_X)$ のもの）に対して成り立つこととは同値である。
Derived Categories, Lemma \ref{derived-lemma-right-adjoint} を参照されたい。
Lemma \ref{spaces-perfect-lemma-exact-sequence-j-star} の識別三角
$$
E \to Rj_{U, *}E|_U \oplus Rj_{V, *}E|_V \to
Rj_{U \times_X V, *}E|_{U \times_X V} \to E[1]
$$
を $D(\mathcal{O}_X)$ において考える。Derived Categories, Lemma
\ref{derived-lemma-prepare-adjoint} により、$Rj_{U, *}E|_U$、
$Rj_{V, *}E|_V$、および
$Rj_{U \times_X V, *}E|_{U \times_X V}$ に対して所望の識別三角を
構成すれば十分である。これで次の段落で論じる主張に帰着する。

\medskip\noindent
$j : U \to X$ を平展射とし、$U$ は準コンパクトかつ準分離的で、
補題が $U$ に対して成り立つとする。$L$ を $D(\mathcal{O}_U)$ の対象とする。
このとき識別三角
$$
E' \to Rj_*L \to K \to E'[1]
$$
が $D(\mathcal{O}_X)$ において存在し、$E'$ が $D_\QCoh(\mathcal{O}_X)$ に属し、
$\Hom(M, K) = 0$ がすべての $M$（$D_\QCoh(\mathcal{O}_X)$ のもの）に
対して成り立つ。これを見るため、識別三角
$$
L' \to L \to Q \to L'[1]
$$
を $D(\mathcal{O}_U)$ において選ぶ。ここで $L'$ は
$D_\QCoh(\mathcal{O}_U)$ に属し、$\Hom(N, Q) = 0$ がすべての
$N$（$D_\QCoh(\mathcal{O}_U)$ のもの）に対して成り立つ。
Derived Categories, Lemma \ref{derived-lemma-right-adjoint} の主張が
必要十分条件であるため、この選択が可能である。識別三角
$$
Rj_*L' \to Rj_*L \to Rj_*Q \to Rj_*L'[1]
$$
を $D(\mathcal{O}_X)$ において得る。Lemma
\ref{spaces-perfect-lemma-quasi-coherence-direct-image} により、$Rj_*L'$ は
$D_\QCoh(\mathcal{O}_X)$ に属する。一方、$M$ が
$D_\QCoh(\mathcal{O}_X)$ に属するならば
$$
\Hom(M, Rj_*Q) = \Hom(Lj^*M, Q) = 0
$$
である。実際、Lemma \ref{spaces-perfect-lemma-quasi-coherence-pullback} により
$Lj^*M$ は $D_\QCoh(\mathcal{O}_U)$ に属する。これで証明が終わる。
\end{proof}

\begin{proof}[第二の証明]
随伴は Derived Categories, Proposition \ref{derived-proposition-brown} により
存在する。仮定が満たされることを確認する。まず
$D_\QCoh(\mathcal{O}_X)$ は直和をもち、直和は包含函手と可換する
（Lemma \ref{spaces-perfect-lemma-quasi-coherence-direct-sums}）。一方、
$D_\QCoh(\mathcal{O}_X)$ はコンパクト生成される。実際、Theorem
\ref{spaces-perfect-theorem-bondal-van-den-Bergh} により完全生成対象をもち、
Proposition \ref{spaces-perfect-proposition-compact-is-perfect} により完全対象は
コンパクトである。% P09-ERR-1014
\end{proof}

\begin{lemma}
\label{spaces-perfect-lemma-pushforward-better-coherator}
$S$ を概型とする。$f : X \to Y$ を $S$ 上の代数空間の準コンパクトかつ
準分離的な射とする。右随伴 $DQ_X$ と $DQ_Y$（包含函手
$D_\QCoh \to D$ のもの）がそれぞれ $X$ と $Y$ に対して存在するならば
$$
Rf_* \circ DQ_X = DQ_Y \circ Rf_*
$$
である。
\end{lemma}

\begin{proof}
Lemma \ref{spaces-perfect-lemma-quasi-coherence-direct-image} により、$Rf_*$ は
$D_\QCoh(\mathcal{O}_X)$ を $D_\QCoh(\mathcal{O}_Y)$ の中へ写すので、
この主張には意味がある。また Lemma
\ref{spaces-perfect-lemma-quasi-coherence-pullback} により、$Lf^*$ も同様に
$D_\QCoh(\mathcal{O}_Y)$ を $D_\QCoh(\mathcal{O}_X)$ の中へ写す。
したがって $Rf_* \circ DQ_X$ と $DQ_Y \circ Rf_*$ はともに
$Lf^* : D_\QCoh(\mathcal{O}_Y) \to D(\mathcal{O}_X)$ の右随伴であり、
主張が従う。
\end{proof}

\begin{remark}
\label{spaces-perfect-remark-explain-consequence}
$S$ を概型とする。$(U \subset X, f : V \to X)$ を $S$ 上の代数空間の
初等識別正方形とする。$X$、$U$、$V$ は準コンパクトかつ準分離的であると仮定する。
Lemma \ref{spaces-perfect-lemma-better-coherator} により、函手 $DQ_X$、$DQ_U$、$DQ_V$、
$DQ_{U \times_X V}$ は存在する。さらに、標準的な識別三角
{\small
$$
DQ_X(K) \to Rj_{U, *}DQ_U(K|_U) \oplus Rj_{V, *}DQ_V(K|_V)
\to Rj_{U \times_X V, *}DQ_{U \times_X V}(K|_{U \times_X V}) \to
$$
}
が任意の $K \in D(\mathcal{O}_X)$ に対して存在する。% P09-ERR-1015
これは Lemma \ref{spaces-perfect-lemma-exact-sequence-j-star} の識別三角に完全函手 $DQ_X$ を
適用し、Lemma \ref{spaces-perfect-lemma-pushforward-better-coherator} を三回用いることで従う。
\end{remark}

\begin{lemma}
\label{spaces-perfect-lemma-boundedness-better-coherator}
$S$ を概型とする。$X$ を $S$ 上の準コンパクトかつ準分離的な代数空間とする。
Lemma \ref{spaces-perfect-lemma-better-coherator} の函手 $DQ_X$ は次の有界性をもつ。
ある整数 $N = N(X)$ が存在し、$K$ が $D(\mathcal{O}_X)$ に属して、
$H^i(U, K) = 0$ が、アフィン $U$ であって $X$ 上平展なものと
$i \not \in [a, b]$ に対して成り立つならば、コホモロジー層
$H^i(DQ_X(K))$ は $i \not \in [a, b + N]$ に対して零である。
\end{lemma}

\begin{proof}
Lemma \ref{spaces-perfect-lemma-induction-principle} の帰納原理を用いて証明する。

\medskip\noindent
$X$ がアフィンならば、$N = 0$ として補題は成り立つ。実際、このとき
$RQ_X = DQ_X$ は $R\Gamma(X, K)$ に付随する準連接層の複体をとることで
与えられる。Lemma \ref{spaces-perfect-lemma-affine-coherator} を参照されたい。

\medskip\noindent
$(U \subset W, f : V \to W)$ を初等識別正方形とし、$W$ は準コンパクトかつ
準分離的、$U \subset W$ は準コンパクト開、$V$ はアフィンであり、補題が
$U$、$V$、および $U \times_W V$ に対して成り立つとする。対応する整数を
$N(U)$、$N(V)$、および $N(U \times_W V)$ とする。ここで $K$ が
$D(\mathcal{O}_X)$ に属し、$H^i(W, K) = 0$ が、すべてのアフィン $W$ であって
$X$ 上平展なものと、すべての $i \not \in [a, b]$ に対して成り立つと仮定する。
% P09-ERR-1016
すると $K|_U$、$K|_V$、$K|_{U \times_W V}$ は同じ性質をもつ。したがって
$RQ_U(K|_U)$、$RQ_V(K|_V)$、および
$RQ_{U \cap V}(K|_{U \times_W V})$ のコホモロジー層は区間
% P09-ERR-1017
$[a, b + \max(N(U), N(V), N(U \times_W V))$
の外で消滅する。% P09-ERR-1018
Lemma \ref{spaces-perfect-lemma-quasi-coherence-direct-image} により、函手
$Rj_{U, *}$、$Rj_{V, *}$、$Rj_{U \times_W V, *}$ は $D_\QCoh$ 上有限な
コホモロジー次元をもつ。したがってある $N$ が存在し、
$Rj_{U, *}DQ_U(K|_U)$、$Rj_{V, *}DQ_V(K|_V)$、および
$Rj_{U \cap V, *}DQ_{U \times_W V}(K|_{U \times_W V})$ のコホモロジー層は
区間 $[a, b + N]$ の外で消滅する。% P09-ERR-1019
最後に Remark \ref{spaces-perfect-remark-explain-consequence} の識別三角から結論が従う。
\end{proof}

\begin{example}
\label{spaces-perfect-example-inverse-limit-quasi-coherent}
$S$ を概型とする。$X$ を $S$ 上の準コンパクトかつ準分離的な代数空間とする。
$(\mathcal{F}_n)$ を $X$ 上の準連接層の逆系とする。$DQ_X$ は右随伴なので
積と可換し、したがって導来極限とも可換する。ゆえに
$$
DQ_X(R\lim \mathcal{F}_n) =
(R\lim\text{ in }D_\QCoh(\mathcal{O}_X))(\mathcal{F}_n)
$$
である。ここで最初の $R\lim$ は $D(\mathcal{O}_X)$ においてとる。
実際、これについて $K = R\lim \mathcal{F}_n$ と書こう。任意のアフィン
$U$ であって $X$ 上平展なものに対して
$$
H^i(U, K) =
H^i(R\Gamma(U, R\lim \mathcal{F}_n)) =
H^i(R\lim R\Gamma(U, \mathcal{F}_n)) =
H^i(R\lim \Gamma(U, \mathcal{F}_n))
$$
である。これはコホモロジーが導来極限と可換し、準連接層
$\mathcal{F}_n$ はアフィン上で高次コホモロジーをもたないためである。
アーベル群の圏における $R\lim$ の計算により、
$H^i(U, K) = 0$ は $i \in [0, 1]$ でない限り成り立つ。
最後に、$R\lim$（$D_\QCoh(\mathcal{O}_X)$ におけるもの）は上により $DQ_X(K)$ であり、
$D^b_\QCoh(\mathcal{O}_X)$ に属し、Lemma
\ref{spaces-perfect-lemma-boundedness-better-coherator} により負次数のコホモロジー層は
消滅する。
\end{example}











\section{コホモロジーと基底変換、IV}
\label{spaces-perfect-section-cohomology-and-base-change-perfect}

\noindent
本節は Derived Categories of Schemes, Section
\ref{perfect-section-cohomology-and-base-change-perfect} の類似である。

\begin{lemma}
\label{spaces-perfect-lemma-cohomology-base-change}
$S$ を概型とする。$f : X \to Y$ を $S$ 上の代数空間の準コンパクトかつ
準分離的な射とする。$E$ が $D_\QCoh(\mathcal{O}_X)$ に属し、
$K$ が $D_\QCoh(\mathcal{O}_Y)$ に属するならば
$$
Rf_*(E) \otimes_{\mathcal{O}_Y}^\mathbf{L} K =
Rf_*(E \otimes_{\mathcal{O}_X}^\mathbf{L} Lf^*K)
$$
である。
\end{lemma}

\begin{proof}
何の仮定もなく、射
$Rf_*(E) \otimes_{\mathcal{O}_Y}^\mathbf{L} K \to
Rf_*(E \otimes_{\mathcal{O}_X}^\mathbf{L} Lf^*K)$
がある。すなわち、これは標準射
$$
Lf^*(Rf_*(E) \otimes_{\mathcal{O}_Y}^\mathbf{L} K) =
Lf^*(Rf_*(E)) \otimes_{\mathcal{O}_X}^\mathbf{L} Lf^*K
\longrightarrow
E \otimes_{\mathcal{O}_X}^\mathbf{L} Lf^*K
$$
の随伴であり、この標準射は射 $Lf^*Rf_*E \to E$ から生じる。
Cohomology on Sites, Lemmas
\ref{sites-cohomology-lemma-pullback-tensor-product} および
\ref{sites-cohomology-lemma-adjoint} を参照されたい。これが同型であることは
$Y$ 上平展局所的に確認できる。したがって $Y$ がアフィン概型である場合に帰着する。

\medskip\noindent
$K = \bigoplus K_i$ が複体 $K_i \in D_\QCoh(\mathcal{O}_Y)$ の直和であるとする。
主張が各 $K_i$ に対して成り立つならば、$K$ に対しても成り立つ。実際、函手
$Lf^*$ と $\otimes^\mathbf{L}$ は構成により直和を保ち、$Rf_*$ は
Lemma \ref{spaces-perfect-lemma-quasi-coherence-pushforward-direct-sums} により
（準連接コホモロジー層をもつ複体について）直和と可換する。さらに
$K \to L \to M \to K[1]$ が $D_\QCoh(Y)$ の識別三角であるとする。
補題の主張が $K, L, M$ のうち二つに対して成り立つならば、三つ目にも成り立つ
（関与する函手は三角圏の完全函手だからである）。

\medskip\noindent
$Y$ はアフィンで、$Y = \Spec(A)$ とする。函手
$\widetilde{\ } : D(A) \to D_\QCoh(\mathcal{O}_Y)$ は
Lemma \ref{spaces-perfect-lemma-derived-quasi-coherent-small-etale-site} および
Derived Categories of Schemes, Lemma
\ref{perfect-lemma-affine-compare-bounded} により同値である。
$T$ を、$K \in D(A)$ に対して補題の主張が $\widetilde{K}$ について
成り立つという性質とする。上の議論と More on Algebra, Remark
\ref{more-algebra-remark-P-resolution} により、$T$ が $A[k]$ に対して
成り立つことを証明すれば十分である。構造層のシフトについて補題の主張は
明らかなので、これで証明が終わる。
\end{proof}

\begin{definition}
\label{spaces-perfect-definition-tor-independent}
$S$ を概型とする。$B$ を $S$ 上の代数空間とする。$X$、$Y$ を $B$ 上の
代数空間とする。$X$ と $Y$ が {\it $B$ 上 Tor 独立} であるとは、すべての
可換図式
$$
\xymatrix{
\Spec(k) \ar[d]_{\overline{y}} \ar[dr]_{\overline{b}} \ar[r]_-{\overline{x}} &
X \ar[d] \\
Y \ar[r] & B
}
$$
で表される幾何学的点について、環
$\mathcal{O}_{X, \overline{x}}$ と $\mathcal{O}_{Y, \overline{y}}$ が
$\mathcal{O}_{B, \overline{b}}$ 上 Tor 独立であることをいう
（More on Algebra, Definition
\ref{more-algebra-definition-tor-independent} を参照）。
\end{definition}

\noindent
次の補題は、特にこの定義が表現可能な代数空間の場合の既存の定義と
一致することを示す。

\begin{lemma}
\label{spaces-perfect-lemma-tor-independent}
$S$ を概型とする。$B$ を $S$ 上の代数空間とする。$X$、$Y$ を $B$ 上の
代数空間とする。次は同値である。
\begin{enumerate}
\item $X$ と $Y$ は $B$ 上 Tor 独立である。
\item すべての可換図式
$$
\xymatrix{
U \ar[d] \ar[r] & W \ar[d] & V \ar[d] \ar[l] \\
X \ar[r] & B & Y \ar[l]
}
$$
で垂直射が平展であるものについて、$U$ と $V$ は $W$ 上 Tor 独立である。
\item (2) のような可換図式で、(a) $W \to B$ が平展かつ全射、(b)
$U \to X \times_B W$ が平展かつ全射、(c)
$V \to Y \times_B W$ が平展かつ全射であるものが存在し、その図式において
$U$ と $V$ は $W$ 上 Tor 独立である。
\item (3) のような可換図式で、$U$、$V$、$W$ が概型であるものが存在し、
概型 $U$ と $V$ は Derived Categories of Schemes, Definition
\ref{perfect-definition-tor-independent} の意味で $W$ 上 Tor 独立である。
\end{enumerate}
\end{lemma}

\begin{proof}
代数空間の平展射 $\varphi : U \to X$ と幾何学的点 $\overline{u}$ に対して、
局所環の射
$\mathcal{O}_{X, \varphi(\overline{u})} \to
\mathcal{O}_{U, \overline{u}}$
は同型である。したがって (1) と (2) の同値が従う。同じく
(1) $\Rightarrow$ (3) も従う。(3) を仮定し、Definition
\ref{spaces-perfect-definition-tor-independent} のような幾何学的点の図式を選ぶ。
仮定により、まず $\overline{b}$ を幾何学的点 $\overline{w}$（$W$ のもの）に
持ち上げ、次に幾何学的点 $(\overline{x}, \overline{b})$ を幾何学的点
$\overline{u}$（$U$ のもの）に持ち上げ、最後に幾何学的点
$(\overline{y}, \overline{b})$ を幾何学的点 $\overline{v}$（$V$ のもの）に
持ち上げることができる。持ち上げの存在には Properties of Spaces, Lemma
\ref{spaces-properties-lemma-geometric-lift-to-usual} を用いる。
上の局所環に関する注意を用いると、与えられた図式について定義の条件が
満たされることが分かる。

\medskip\noindent
これらの点を最初に選んでおけば、(4) は Definition
\ref{spaces-perfect-definition-tor-independent} と Derived Categories of Schemes, Definition
\ref{perfect-definition-tor-independent} が、$X$、$Y$、$B$ が概型であるとき
一致するという主張に帰着することは明らかである。

\medskip\noindent
$\overline{x}, \overline{b}, \overline{y}$ を Definition
\ref{spaces-perfect-definition-tor-independent} のような点で、それぞれ点 $x, y, b$ の
上にあるものとする。
$\mathcal{O}_{X, \overline{x}} = \mathcal{O}_{X, x}^{sh}$ であること
（Properties of Spaces, Lemma
\ref{spaces-properties-lemma-describe-etale-local-ring}）と、他の二つについての
同様の事実を思い出そう。Algebra, Lemma
\ref{algebra-lemma-strictly-henselian-functorial-improve} により、
$\mathcal{O}_{X, \overline{x}}$ は
$\mathcal{O}_{X, x} \otimes_{\mathcal{O}_{B, b}}
\mathcal{O}_{B, \overline{b}}$ の強 Hensel 化である。特に環準同型
$$
\mathcal{O}_{X, x} \otimes_{\mathcal{O}_{B, b}} \mathcal{O}_{B, \overline{b}}
\longrightarrow
\mathcal{O}_{X, \overline{x}}
$$
は平坦である（More on Algebra, Lemma
\ref{more-algebra-lemma-dumb-properties-henselization}）。More on Algebra,
Lemma \ref{more-algebra-lemma-tor-independent-flat} により
$$
\text{Tor}_i^{\mathcal{O}_{B, b}}(\mathcal{O}_{X, x}, \mathcal{O}_{Y, y})
\otimes_{\mathcal{O}_{X, x} \otimes_{\mathcal{O}_{B, b}} \mathcal{O}_{Y, y}}
(\mathcal{O}_{X, \overline{x}} \otimes_{\mathcal{O}_{B, \overline{b}}}
\mathcal{O}_{Y, \overline y})
=
\text{Tor}_i^{\mathcal{O}_{B, \overline{b}}}(
\mathcal{O}_{X, \overline{x}}, \mathcal{O}_{Y, \overline{y}})
$$
である。したがって $X$ と $Y$ が概型として $B$ 上 Tor 独立ならば、
$X$ と $Y$ は代数空間として $B$ 上 Tor 独立である。

\medskip\noindent
逆に、$X$、$Y$、$B$ はアフィンであると仮定してよい。上で述べたことにより
環準同型
$$
\mathcal{O}_{X, x} \otimes_{\mathcal{O}_{B, b}} \mathcal{O}_{Y, y}
\longrightarrow
\mathcal{O}_{X, \overline{x}} \otimes_{\mathcal{O}_{B, \overline{b}}}
\mathcal{O}_{Y, \overline y}
$$
は平坦である。さらに、スペクトル上の射の像は、すべての素イデアル
$$
\mathfrak s \subset
\mathcal{O}_{X, x} \otimes_{\mathcal{O}_{B, b}} \mathcal{O}_{Y, y}
$$
で $\mathfrak m_x$ と $\mathfrak m_y$ の上にあるものを含む。
したがってこのことと上の Tor の表示式から、$X$ と $Y$ が
代数空間として $B$ 上 Tor 独立ならば
$$
\text{Tor}_i^{\mathcal{O}_{B, b}}
(\mathcal{O}_{X, x}, \mathcal{O}_{Y, y})_\mathfrak s = 0
$$
がすべての $i > 0$ と上のようなすべての $\mathfrak s$ に対して成り立つ。
More on Algebra, Lemma \ref{more-algebra-lemma-tor-independent} を環準同型
$\Gamma(B, \mathcal{O}_B) \to \Gamma(X, \mathcal{O}_X)$ および
$\Gamma(B, \mathcal{O}_B) \to \Gamma(X, \mathcal{O}_X)$ に適用すると、
% P09-ERR-1020
$X$ と $Y$ は $B$ 上 Tor 独立であることが従う。
\end{proof}

\begin{lemma}
\label{spaces-perfect-lemma-compare-base-change}
$S$ を概型とする。$g : Y' \to Y$ を $S$ 上の代数空間の射とする。
$f : X \to Y$ を $S$ 上の代数空間の準コンパクトかつ準分離的な射とする。
基底変換図式
$$
\xymatrix{
X' \ar[r]_{g'} \ar[d]_{f'} &
X \ar[d]^f \\
Y' \ar[r]^g &
Y
}
$$
を考える。$X$ と $Y'$ が $Y$ 上 Tor 独立ならば、すべての
$E \in D_\QCoh(\mathcal{O}_X)$ に対して
$Rf'_*L(g')^*E = Lg^*Rf_*E$ である。
\end{lemma}

\begin{proof}
任意の対象 $E$（$D(\mathcal{O}_X)$ のもの）に対して、Cohomology on Sites, Remark
\ref{sites-cohomology-remark-base-change} から標準的な基底変換射
$Lg^*Rf_*E \to Rf'_*L(g')^*E$ を得る。これが同型であることは $Y'$ 上
平展局所的に確認できる。したがって $g : Y' \to Y$ はアフィン概型の射であると
仮定してよい。特に $g$ はアフィンであるから
$$
Rg_*Lg^*Rf_*E \to Rg_*Rf'_*L(g')^*E = Rf_*(Rg'_* L(g')^* E)
$$
が同型であることを示せば十分である。Lemma \ref{spaces-perfect-lemma-affine-morphism} を
参照されたい（また、対象 $Rf'_*L(g')^*E$ と $Lg^*Rf_*E$ が
準連接コホモロジー層をもつことには Lemmas
\ref{spaces-perfect-lemma-quasi-coherence-pullback}、
\ref{spaces-perfect-lemma-quasi-coherence-tensor-product}、および
\ref{spaces-perfect-lemma-quasi-coherence-direct-image} を用いる）。
$g'$ もアフィンであることに注意する（Morphisms of Spaces, Lemma
\ref{spaces-morphisms-lemma-base-change-affine}）。
Lemma \ref{spaces-perfect-lemma-affine-morphism-pull-push} により、この射は
$$
Rf_*E \otimes_{\mathcal{O}_Y}^\mathbf{L} g_*\mathcal{O}_{Y'}
\longrightarrow
Rf_*(E \otimes_{\mathcal{O}_X}^\mathbf{L} g'_*\mathcal{O}_{X'})
$$
となる。$g'_*\mathcal{O}_{X'} = f^*g_*\mathcal{O}_{Y'}$ である。
したがって Lemma \ref{spaces-perfect-lemma-cohomology-base-change} により、
$Lf^*g_*\mathcal{O}_{Y'} = f^*g_*\mathcal{O}_{Y'}$ を証明すれば十分である。
これは $X$ と $Y'$ が $Y$ 上 Tor 独立であるという仮定から従う。実際、
$X$ 上平展局所的に確認できるので、$X$ もアフィンであると仮定してよい。
$X = \Spec(A)$、$Y = \Spec(R)$、$Y' = \Spec(R')$ と書く。
仮定は $A$ と $R'$ が $R$ 上 Tor 独立であることを含意する
（Lemma \ref{spaces-perfect-lemma-tor-independent} および More on Algebra, Lemma
\ref{more-algebra-lemma-tor-independent}）。すなわち、
$\text{Tor}_i^R(A, R') = 0$ が $i > 0$ に対して成り立つ。言い換えると
$A \otimes_R^\mathbf{L} R' = A \otimes_R R'$ であり、これはまさに
$Lf^*g_*\mathcal{O}_{Y'} = f^*g_*\mathcal{O}_{Y'}$ を意味する。
\end{proof}

\noindent
次の補題は双対化複体の章で用いる。

\begin{lemma}
\label{spaces-perfect-lemma-affine-morphism-and-hom-out-of-perfect}
$g : S' \to S$ をアフィン概型の射とする。準コンパクトかつ準分離的な
代数空間の Cartesian 正方形
$$
\xymatrix{
X' \ar[r]_{g'} \ar[d]_{f'} & X \ar[d]^f \\
S' \ar[r]^g & S
}
$$
を考える。$g$ と $f$ は Tor 独立であると仮定する。
$S = \Spec(R)$ および $S' = \Spec(R')$ と書く。
$M, K \in D(\mathcal{O}_X)$ に対する標準射
$$
R\Hom_X(M, K) \otimes^\mathbf{L}_R R'
\longrightarrow
R\Hom_{X'}(L(g')^*M, L(g')^*K)
$$
は $D(R')$ において、次の二つの場合に同型である。
\begin{enumerate}
\item $M \in D(\mathcal{O}_X)$ は完全で、$K \in D_\QCoh(X)$ である。
\item $M \in D(\mathcal{O}_X)$ は擬連接で、$K \in D_\QCoh^+(X)$ であり、
$R'$ は $R$ 上有限 Tor 次元をもつ。
\end{enumerate}
\end{lemma}

\begin{proof}
大域 Hom 複体の標準射
$R\Hom_X(M, K) \to R\Hom_{X'}(L(g')^*M, L(g')^*K)$
が $D(\Gamma(X, \mathcal{O}_X))$ において存在する。Cohomology on Sites,
Section \ref{sites-cohomology-section-global-RHom} を参照されたい。
スカラーを制限すれば、これを $D(R)$ における射とみなせる。
次に、スカラー制限と $- \otimes_R^\mathbf{L} R'$ の随伴性を用いて、
補題の表示射を得る。射を定義したので、アーベル群の導来圏で同型であることを
証明すれば十分である。

\medskip\noindent
右辺は Lemma \ref{spaces-perfect-lemma-affine-morphism-pull-push} により
$$
R\Hom_X(M, R(g')_*L(g')^*K) =
R\Hom_X(M, K \otimes_{\mathcal{O}_X}^\mathbf{L} g'_*\mathcal{O}_{X'})
$$
に等しい。いずれの場合にも、複体 $R\SheafHom(M, K)$ は Lemma
\ref{spaces-perfect-lemma-quasi-coherence-internal-hom} により
$D_\QCoh(\mathcal{O}_X)$ の対象である。自然な射
$$
R\SheafHom(M, K) \otimes_{\mathcal{O}_X}^\mathbf{L} g'_*\mathcal{O}_{X'}
\longrightarrow
R\SheafHom(M, K \otimes_{\mathcal{O}_X}^\mathbf{L} g'_*\mathcal{O}_{X'})
$$
があり、いずれの場合にも Lemma
\ref{spaces-perfect-lemma-internal-hom-evaluate-tensor-isomorphism} により同型である。
% P09-ERR-1023
(2) の場合にこの補題を適用できることを見るため、
$g'_*\mathcal{O}_{X'} = Rg'_*\mathcal{O}_{X'} =
Lf^*g_*\mathcal{O}_X$ であることに注意する。第二の等式は Lemma
\ref{spaces-perfect-lemma-compare-base-change} による。% P09-ERR-1021
Derived Categories of Schemes, Lemma
\ref{perfect-lemma-tor-dimension-affine}、Lemma
\ref{spaces-perfect-lemma-descend-tor-amplitude}、および Cohomology on Sites, Lemma
\ref{sites-cohomology-lemma-tor-amplitude-pullback} を用いると、
$g'_*\mathcal{O}_{X'}$ は有限 Tor 次元をもつことが分かる。
したがって、いずれの場合にも $K$ を $R\SheafHom(M, K)$ で置き換えることで、
$$
R\Gamma(X, K) \otimes^\mathbf{L}_A A' \longrightarrow
R\Gamma(X, K \otimes^\mathbf{L}_{\mathcal{O}_X} g'_*\mathcal{O}_{X'})
$$
が同型であることに帰着する。% P09-ERR-1022
左辺は Lemma \ref{spaces-perfect-lemma-affine-morphism-pull-push} により
$R\Gamma(X', L(g')^*K)$ に等しいことに注意する。
したがって結果は Lemma \ref{spaces-perfect-lemma-compare-base-change} から従う。
\end{proof}

\begin{remark}
\label{spaces-perfect-remark-multiplication-map}
Lemma \ref{spaces-perfect-lemma-affine-morphism-and-hom-out-of-perfect} の記法を用いる。図式
$$
\xymatrix{
R\Hom_X(M, Rg'_*L) \otimes_R^\mathbf{L} R' \ar[r] \ar[d]_\mu &
R\Hom_{X'}(L(g')^*M, L(g')^*Rg'_*L) \ar[d]^a \\
R\Hom_X(M, R(g')_*L) \ar@{=}[r] &
R\Hom_{X'}(L(g')^*M, L)
}
$$
は可換である。ここで上の水平射は補題の射、$\mu$ は乗法射、$a$ は随伴射
$L(g')^*Rg'_*L \to L$ から生じる。乗法射とは、
随伴射 $K' \otimes_R^\mathbf{L} R' \to K'$ のことであり、これは任意の
$K' \in D(R')$ に対するものである。
\end{remark}

\begin{lemma}
\label{spaces-perfect-lemma-tor-independence-and-tor-amplitude}
$S$ を概型とする。代数空間の Cartesian 正方形
$$
\xymatrix{
X' \ar[r]_{g'} \ar[d]_{f'} & X \ar[d]^f \\
Y' \ar[r]^g & Y
}
$$
を $S$ 上で考える。$g$ と $f$ は Tor 独立であると仮定する。
\begin{enumerate}
\item $E \in D(\mathcal{O}_X)$ が $[a, b]$ に Tor 振幅をもち、これは
$f^{-1}\mathcal{O}_Y$-加群の複体としての振幅であるとする。このとき
$L(g')^*E$ は $[a, b]$ に Tor 振幅をもち、これは
$f^{-1}\mathcal{O}_{Y'}$-加群の複体としての振幅である。
% P09-ERR-1024
\item $\mathcal{G}$ が $\mathcal{O}_X$-加群で $Y$ 上平坦ならば、
$L(g')^*\mathcal{G} = (g')^*\mathcal{G}$ である。
\end{enumerate}
\end{lemma}

\begin{proof}
Tor 次元は茎で計算できる。Cohomology on Sites, Lemma
\ref{sites-cohomology-lemma-tor-amplitude-stalk} および Properties of Spaces,
Theorem \ref{spaces-properties-theorem-exactness-stalks} を参照されたい。
$\overline{x}'$ が $X'$ の幾何学的点で、その像が $\overline{x}$（$X$ の点）ならば
$$
(L(g')^*E)_{\overline{x}'} =
E_{\overline{x}}
\otimes_{\mathcal{O}_{X, \overline{x}}}^\mathbf{L}
\mathcal{O}_{X', \overline{x}'}
$$
である。$\overline{y}'$（$Y'$ の点）と $\overline{y}$（$Y$ の点）を、
それぞれ $\overline{x}'$ と $\overline{x}$ の像とする。
$X$ と $Y'$ は $Y$ 上 Tor 独立なので、More on Algebra, Lemma
\ref{more-algebra-lemma-base-change-comparison} を適用すると、表示式の右辺は
$$
E_{\overline{x}}
\otimes_{\mathcal{O}_{Y, \overline{y}}}^\mathbf{L}
\mathcal{O}_{Y', \overline{y}'}
$$
に $D(\mathcal{O}_{Y', \overline{y}'})$ において等しいことが分かる。
したがって (1) は More on Algebra, Lemma
\ref{more-algebra-lemma-pull-tor-amplitude} から従う。(2) を見るには、
$\mathcal{G}$ の平坦性が $\mathcal{G}[0]$ が $[0, 0]$ に Tor 振幅をもつという
条件と同値であることに注意する。(1) を適用すれば結論を得る。
\end{proof}










\section{コホモロジーと基底変換、V}
\label{spaces-perfect-section-cohomology-base-change}

\noindent
本節は Derived Categories of Schemes, Section
\ref{perfect-section-cohomology-base-change} の類似である。Section
\ref{spaces-perfect-section-cohomology-and-base-change-perfect} では、射が Tor 独立ならば
基底変換定理が成り立つことを見た。そのような条件がなければ、アフィンの場合でさえ
基底変換定理は存在しえない。More on Algebra, Section
\ref{more-algebra-section-tor-independence} を参照されたい。本節では、
「一度に一つの複体について」基底変換の結果がいつ得られるかを解析する。

\medskip\noindent
このため、$S$ を基礎概型とし、可換図式
$$
\xymatrix{
X' \ar[r]_{g'} \ar[d]_{f'} &
X \ar[d]^f \\
Y' \ar[r]^g &
Y
}
$$
が $S$ 上の代数空間からなると仮定する（通常は Cartesian であると仮定する）。
$K \in D_\QCoh(\mathcal{O}_X)$ とし、
$L(g')^*K \to K'$ を $D_\QCoh(\mathcal{O}_{X'})$ の射とする。
幾何学的点 $\overline{x}'$（$X'$ の点）に対して、幾何学的点
$\overline{x} = g'(\overline{x}')$、
$\overline{y}' = f'(\overline{x}')$、
$\overline{y} = f(\overline{x}) = g(\overline{y}')$
（それぞれ $X$、$Y'$、$Y$ の点）を考える。このとき射
$$
K_{\overline{x}}
\otimes_{\mathcal{O}_{Y, \overline{y}}}^\mathbf{L}
\mathcal{O}_{Y', \overline{y}'}
\to
K_{\overline{x}}
\otimes_{\mathcal{O}_{X, \overline{x}}}^\mathbf{L}
\mathcal{O}_{X', \overline{x}'} \to
K'_{\overline{x}'}
$$
を考えられる。ここで第一の矢印は More on Algebra, Equation
(\ref{more-algebra-equation-comparison-map}) であり、第二の矢印は
$$
(L(g')^*K)_{\overline{x}'} =
K_{\overline{x}} \otimes_{\mathcal{O}_{X, \overline{x}}}^\mathbf{L}
\mathcal{O}_{X', \overline{x}'}
$$
と与えられた射 $L(g')^*K \to K'$ から生じる。各
$i \in \mathbf{Z}$ に対して、
$\mathcal{O}_{X, \overline{x}} \otimes_{\mathcal{O}_{Y, \overline{y}}}
\mathcal{O}_{Y', \overline{y}'}$-加群構造が
$H^i(K_{\overline{x}}
\otimes_{\mathcal{O}_{Y, \overline{y}}}^\mathbf{L}
\mathcal{O}_{Y', \overline{y}'})$ 上に得られる。すべてを合わせると標準射
\begin{equation}
\label{spaces-perfect-equation-bc}
H^i(K_{\overline{x}}
\otimes_{\mathcal{O}_{Y, \overline{y}}}^\mathbf{L}
\mathcal{O}_{Y', \overline{y}'})
\otimes_{(\mathcal{O}_{X, \overline{x}}
\otimes_{\mathcal{O}_{Y, \overline{y}}}
\mathcal{O}_{Y', \overline{y}'})}
\mathcal{O}_{X', \overline{x}'}
\longrightarrow
H^i(K'_{\overline{x}'})
\end{equation}
が $\mathcal{O}_{X', \overline{x}'}$-加群の射として得られる。

\begin{lemma}
\label{spaces-perfect-lemma-single-complex-base-change-condition}
$S$ を概型とする。
$$
\xymatrix{
X' \ar[r]_{g'} \ar[d]_{f'} &
X \ar[d]^f \\
Y' \ar[r]^g &
Y
}
$$
を $S$ 上の代数空間の Cartesian 図式とする。
$K \in D_\QCoh(\mathcal{O}_X)$ とし、$L(g')^*K \to K'$ を
$D_\QCoh(\mathcal{O}_{X'})$ の射とする。次は同値である。
\begin{enumerate}
\item 任意の $x' \in X'$ と $i \in \mathbf{Z}$ に対して、射
(\ref{spaces-perfect-equation-bc}) は同型である。% P09-ERR-1026
\item 任意の可換図式
$$
\xymatrix{
& U \ar[d] \ar[rd]^a \\
V' \ar[r] \ar[rd]^c & V \ar[rd]^b & X \ar[d]^f \\
& Y' \ar[r]^g & Y
}
$$
で $a, b, c$ が平展、$U, V, V'$ が概型であり、
$U' = V' \times_V U$ であるものについて、Derived Categories of Schemes,
Lemma \ref{perfect-lemma-single-complex-base-change-condition} の同値な条件が
$(U \to X)^*K$ と $(U' \to X')^*K'$ に対して成り立つ。
\item (2) のような図式で $U' \to X'$ が全射であるものが一つ存在する。
\end{enumerate}
\end{lemma}

\begin{proof}
(1) は $X'$ 上平展局所的である。既知の形式的含意をたどると、
$X, Y, Y', X'$ がそれぞれ概型 $X_0, Y_0, Y'_0, X'_0$ によって
表現される場合に、この補題の条件 (1) と Derived Categories of Schemes,
Lemma \ref{perfect-lemma-single-complex-base-change-condition} の条件 (1) が
同値であることを証明すれば十分である。射 $f_0, g_0, g'_0, f'_0$ を、
これらの概型の間でそれぞれ $f, g, g', f'$ に対応するものとする。
Lemma \ref{spaces-perfect-lemma-derived-quasi-coherent-small-etale-site} により、
$K = \epsilon^*K_0$ および $K' = \epsilon^*K'_0$ と仮定してよい。ここで
$K_0 \in D_\QCoh(\mathcal{O}_{X_0})$ および
$K'_0 \in D_\QCoh(\mathcal{O}_{X'_0})$ である。
さらに、Remark \ref{spaces-perfect-remark-match-total-direct-images} の記法のもとで、射
$Lg^*K \to K'$ は射 $L(g_0)^*K_0 \to K'_0$ の引き戻しである。
% P09-ERR-1025
$\mathcal{O}_{X, \overline{x}}$ は $\mathcal{O}_{X, x}$ の強 Hensel 化である
（Properties of Spaces, Lemma
\ref{spaces-properties-lemma-describe-etale-local-ring}）こと、および
$$
K_{\overline{x}} =
K_{0, x} \otimes_{\mathcal{O}_{X, x}}^\mathbf{L} \mathcal{O}_{X, \overline{x}}
\quad\text{and}\quad
K'_{\overline{x}'} =
K'_{0, x'} \otimes_{\mathcal{O}_{X', x'}}^\mathbf{L}
\mathcal{O}_{X', \overline{x}'}
$$
であることを思い出そう（Properties of Spaces, Lemma
\ref{spaces-properties-lemma-stalk-quasi-coherent} と同様である）。
可換図式
$$
\xymatrix{
H^i(K_{\overline{x}}
\otimes_{\mathcal{O}_{Y, \overline{y}}}^\mathbf{L}
\mathcal{O}_{Y', \overline{y}'})
\otimes_{(\mathcal{O}_{X, \overline{x}}
\otimes_{\mathcal{O}_{Y, \overline{y}}}
\mathcal{O}_{Y', \overline{y}'})}
\mathcal{O}_{X', \overline{x}'}
\ar[r] &
H^i(K'_{\overline{x}'}) \\
H^i(K_{0, x} \otimes_{\mathcal{O}_{Y, y}}^\mathbf{L} \mathcal{O}_{Y', y'})
\otimes_{(\mathcal{O}_{X, x} \otimes_{\mathcal{O}_{Y, y}} \mathcal{O}_{Y', y'})}
\mathcal{O}_{X', x'}
\ar[u] \ar[r] &
H^i(K'_{0, x'}) \ar[u]
}
$$
を考える。下の水平矢印が同型であることと上の水平矢印が同型であることが
同値であることを示さなければならない。
$\mathcal{O}_{X', x'} \to \mathcal{O}_{X', \overline{x}'}$ は忠実平坦なので
（More on Algebra, Lemma
\ref{more-algebra-lemma-dumb-properties-henselization}）、上の矢印がこの射による
下の矢印の基底変換であることを示せば十分である。右側の対象については、
上で与えた茎の間の関係から直ちに従う。左側の対象については
\begin{align*}
&
H^i\left(
(K_{0, x}
\otimes_{\mathcal{O}_{X, x}}^\mathbf{L} \mathcal{O}_{X, \overline{x}})
\otimes_{\mathcal{O}_{Y, \overline{y}}}^\mathbf{L}
\mathcal{O}_{Y', \overline{y}'}\right) \\
& =
H^i(K_{0, x} \otimes_{\mathcal{O}_{Y, y}}^\mathbf{L} \mathcal{O}_{Y', y'})
\otimes_{(\mathcal{O}_{X, x} \otimes_{\mathcal{O}_{Y, y}} \mathcal{O}_{Y', y'})}
(\mathcal{O}_{X, \overline{x}}
\otimes_{\mathcal{O}_{Y, \overline{y}}}
\mathcal{O}_{Y', \overline{y}'})
\end{align*}
を示せば十分である。これは More on Algebra, Lemma
\ref{more-algebra-lemma-lemma-tor-independent-flat-compare} から従う。
この補題の平坦性の仮定は、上で述べたこと、および Algebra, Lemma
\ref{algebra-lemma-strictly-henselian-functorial-improve} が
$\mathcal{O}_{X, \overline{x}}$ は
$\mathcal{O}_{X, x} \otimes_{\mathcal{O}_{Y, y}}
\mathcal{O}_{Y, \overline{y}}$ の強 Hensel 化であり、
$\mathcal{O}_{Y', \overline{y}'}$ は
$\mathcal{O}_{Y', y'} \otimes_{\mathcal{O}_{Y, y}}
\mathcal{O}_{Y, \overline{y}}$ の強 Hensel 化であることを含意することから
満たされる。
\end{proof}

\begin{lemma}
\label{spaces-perfect-lemma-single-complex-base-change}
$S$ を概型とする。
$$
\xymatrix{
X' \ar[r]_{g'} \ar[d]_{f'} &
X \ar[d]^f \\
Y' \ar[r]^g &
Y
}
$$
を $S$ 上の代数空間の Cartesian 図式とする。
$K \in D_\QCoh(\mathcal{O}_X)$ とし、$L(g')^*K \to K'$ を
$D_\QCoh(\mathcal{O}_{X'})$ の射とする。次を仮定する。
\begin{enumerate}
\item Lemma \ref{spaces-perfect-lemma-single-complex-base-change-condition} の同値な条件が成り立つ。
\item $f$ は準コンパクトかつ準分離的である。
\end{enumerate}
このとき合成 $Lg^*Rf_*K \to Rf'_*L(g')^*K \to Rf'_*K'$ は同型である。
\end{lemma}

\begin{proof}
射が同型であることは $Y'$ 上平展局所的に確認できる。したがって
$g : Y' \to Y$ はアフィン概型の射であると仮定してよい。この場合、
Lemma \ref{spaces-perfect-lemma-induction-principle} の帰納原理を用いて、準コンパクトかつ
準分離的な代数空間 $U$ であって $X$ 上平展なものに対して、同様に構成した射
$Lg^*R(U \to Y)_*K|_U \to R(U' \to Y')_*K'|_{U'}$
が同型であることを証明する。ここで
$U' = X' \times_{g', X} U = Y' \times_{g, Y} U$ である。

\medskip\noindent
$U$ が概型ならば（例えばアフィンならば）結果は成り立つ。実際、その場合
$Y, Y', U, U'$ は概型であり、Lemma
\ref{spaces-perfect-lemma-derived-quasi-coherent-small-etale-site} により $K$ と $K'$ は
基礎にある概型の導来圏の対象から来る。また Derived Categories of Schemes,
Lemma \ref{perfect-lemma-single-complex-base-change-condition} の条件は、
Lemma \ref{spaces-perfect-lemma-single-complex-base-change-condition} によりこれらの複体に
対して成り立つ。したがって Remark
\ref{spaces-perfect-remark-match-total-direct-images} で説明した両立性により、概型の場合の
結果、すなわち Derived Categories of Schemes, Lemma
\ref{perfect-lemma-single-complex-base-change} を適用できる。

\medskip\noindent
帰納段階を行う。$(U \subset W, V \to W)$ を初等識別正方形とする。
$W$ は $X$ 上平展な準コンパクトかつ準分離的な代数空間、$U$ は
準コンパクト、$V$ はアフィンであり、結果が $U$、$V$、および
$U \times_W V$ に対して成り立つとする。記法を容易にするため $W$ を $X$ で
置き換える（この時点では許される）。% P09-ERR-1027
明らかな射を $a : U \to Y$、$b : V \to Y$、および
$c : U \times_X V \to Y$ と記す。$a' : U' \to Y'$、$b' : V' \to Y'$、および
$c' : U' \times_{X'} V' \to Y'$ を、それぞれ $a$、$b$、$c$ の基底変換とする。
相対 Mayer--Vietoris の
識別三角（Lemma \ref{spaces-perfect-lemma-unbounded-relative-mayer-vietoris}）を用いると、
可換図式
$$
\xymatrix{
Lg^*Rf_*K \ar[r] \ar[d] & Rf'_*K' \ar[d] \\
Lg^*Ra_*K|_U \oplus Lg^*Rb_*K|_V \ar[r] \ar[d] &
Ra'_* K'|_{U'} \oplus Rb'_* K'|_{V'} \ar[d] \\
Lg^*Rc_*K|_{U \times_X V} \ar[r] \ar[d] &
Rc'_*K'|_{U' \times_{X'} V'} \ar[d] \\
Lg^*Rf_* K[1] \ar[r] &
Rf'_* K'[1]
}
$$
を得る。第二および第三の水平矢印は同型なので、第一の水平矢印も同型である
（Derived Categories, Lemma
\ref{derived-lemma-third-isomorphism-triangle}）。これで補題の証明が終わる。
\end{proof}

\begin{lemma}
\label{spaces-perfect-lemma-single-complex-base-change-condition-inherited}
$S$ を概型とする。
$$
\xymatrix{
X' \ar[r]_{g'} \ar[d]_{f'} &
X \ar[d]^f \\
S' \ar[r]^g &
S
}
$$
を $S$ 上の代数空間の Cartesian 図式とする。
$K \in D_\QCoh(\mathcal{O}_X)$ とし、$L(g')^*K \to K'$ を
$D_\QCoh(\mathcal{O}_{X'})$ の射とする。Lemma
\ref{spaces-perfect-lemma-single-complex-base-change-condition} の同値な条件が成り立つならば、
\begin{enumerate}
\item $E \in D_\QCoh(\mathcal{O}_X)$ に対して、Lemma
\ref{spaces-perfect-lemma-single-complex-base-change-condition} の同値な条件が
$L(g')^*(E \otimes^\mathbf{L} K) \to
L(g')^*E \otimes^\mathbf{L} K'$ に対して成り立つ。
\item $E$ が $D(\mathcal{O}_X)$ に属して完全ならば、Lemma
\ref{spaces-perfect-lemma-single-complex-base-change-condition} の同値な条件が
$L(g')^*R\SheafHom(E, K) \to R\SheafHom(L(g')^*E, K')$ に対して成り立つ。
\item $K$ が下に有界で、$E$ が $D(\mathcal{O}_X)$ に属して擬連接ならば、
Lemma \ref{spaces-perfect-lemma-single-complex-base-change-condition} の同値な条件が
$L(g')^*R\SheafHom(E, K) \to R\SheafHom(L(g')^*E, K')$ に対して成り立つ。
\end{enumerate}
\end{lemma}

\begin{proof}
関与する複体が準連接コホモロジー層をもつので、主張には意味がある。これは Lemmas
\ref{spaces-perfect-lemma-quasi-coherence-pullback}、
\ref{spaces-perfect-lemma-quasi-coherence-tensor-product}、
\ref{spaces-perfect-lemma-quasi-coherence-internal-hom}、および Cohomology on Sites, Lemmas
\ref{sites-cohomology-lemma-pseudo-coherent-pullback}、
\ref{sites-cohomology-lemma-perfect-pullback} による。
このことを踏まえると、(1) の場合に射 (\ref{spaces-perfect-equation-bc}) が同型であることは、
テンソル積の推移性を用いて (\ref{spaces-perfect-equation-bc}) の始域と終域を計算すれば
確認できる。More on Algebra, Lemma
\ref{more-algebra-lemma-triple-tensor-product} を参照されたい。
(2) と (3) の射は合成
$$
L(g')^*R\SheafHom(E, K) \to R\SheafHom(L(g')^*E, L(g')^*K)
\to R\SheafHom(L(g')^*E, K')
$$
である。第一の矢印は Cohomology on Sites, Remark
\ref{sites-cohomology-remark-prepare-fancy-base-change} であり、
第二の矢印は与えられた射 $L(g')^*K \to K'$ から生じる。
射 (\ref{spaces-perfect-equation-bc}) が同型であることを証明するには、(2) の場合
$E_x$ を有限生成射影 $\mathcal{O}_{X. x}$-加群の有界複体で表し、
% P09-ERR-1028
(3) の場合は有限自由加群の上に有界な複体で表して、矢印の始域と終域を計算する。
いくつかの詳細は省略する。
\end{proof}
\begin{lemma}
\label{spaces-perfect-lemma-base-change-tensor}
$S$ を概型とする。$f : X \to Y$ を $S$ 上の代数空間の準コンパクトかつ
準分離的な射とする。$E \in D_\QCoh(\mathcal{O}_X)$ とする。
$\mathcal{G}^\bullet$ を、準連接 $\mathcal{O}_X$-加群で $Y$ 上平坦なものからなる
上に有界な複体とする。このとき
$$
Rf_*(E \otimes^\mathbf{L}_{\mathcal{O}_X} \mathcal{G}^\bullet)
$$
の形成は任意の基底変換と可換する（正確な主張は証明を参照）。
\end{lemma}

\begin{proof}
主張の意味は次のとおりである。$g : Y' \to Y$ を代数空間の射とし、基底変換図式
$$
\xymatrix{
X' \ar[r]_{g'} \ar[d]_{f'} &
X \ar[d]^f \\
Y' \ar[r]^g &
Y
}
$$
を考える。言い換えると $X' = Y' \times_Y X$ である。補題は
$$
Lg^*Rf_*(E \otimes^\mathbf{L}_{\mathcal{O}_X} \mathcal{G}^\bullet)
\longrightarrow
Rf'_*(L(g')^*E \otimes^\mathbf{L}_{\mathcal{O}_{X'}} (g')^*\mathcal{G}^\bullet)
$$
が同型であると主張する。右辺では $\mathcal{G}^\bullet$ に導来引き戻しを
{\bf 用いていない}ことに注意する。これを証明するため Lemmas
\ref{spaces-perfect-lemma-single-complex-base-change} および
\ref{spaces-perfect-lemma-single-complex-base-change-condition-inherited} を適用すると、
標準射
$$
L(g')^*\mathcal{G}^\bullet \to (g')^*\mathcal{G}^\bullet
$$
が Lemma \ref{spaces-perfect-lemma-single-complex-base-change-condition} の同値な条件を
満たすことを証明すれば十分である。これは茎で条件を確認すれば従う。
実際、複体に関する仮定により
$$
\mathcal{G}^\bullet_{\overline{x}}
\otimes_{\mathcal{O}_{Y, \overline{y}}}
\mathcal{O}_{Y', \overline{y}'}
$$
が導来テンソル積を計算する。これは複体 $\mathcal{G}^\bullet$ に関する仮定による。
\end{proof}

\begin{lemma}
\label{spaces-perfect-lemma-base-change-RHom}
$S$ を概型とする。$f : X \to Y$ を $S$ 上の代数空間の準コンパクトかつ
準分離的な射とする。$E$ を $D(\mathcal{O}_X)$ の対象とする。
$\mathcal{G}^\bullet$ を準連接 $\mathcal{O}_X$-加群の複体とする。次のいずれかを
仮定する。
\begin{enumerate}
\item $E$ は完全で、$\mathcal{G}^\bullet$ は上に有界であり、
$\mathcal{G}^n$ は $Y$ 上平坦である。% P09-ERR-1030
\item $E$ は擬連接で、$\mathcal{G}^\bullet$ は有界であり、
$\mathcal{G}^n$ は $Y$ 上平坦である。
\end{enumerate}
このとき
$$
Rf_*R\SheafHom(E, \mathcal{G}^\bullet)
$$
の形成は任意の基底変換と可換する（正確な主張は証明を参照）。
\end{lemma}

\begin{proof}
主張の意味は次のとおりである。$g : Y' \to Y$ を代数空間の射とし、
基底変換図式
$$
\xymatrix{
X' \ar[r]_h \ar[d]_{f'} &
X \ar[d]^f \\
Y' \ar[r]^g &
Y
}
$$
を考える。% P09-ERR-1029
言い換えると $X' = Y' \times_Y X$ である。補題は
$$
Lg^*Rf_*R\SheafHom(E, \mathcal{G}^\bullet)
\longrightarrow
R(f')_*R\SheafHom(L(g')^*E, (g')^*\mathcal{G}^\bullet)
$$
が同型であると主張する。右辺では $\mathcal{G}^\bullet$ に導来引き戻しを
{\bf 用いていない}ことに注意する。これを証明するため Lemmas
\ref{spaces-perfect-lemma-single-complex-base-change} および
\ref{spaces-perfect-lemma-single-complex-base-change-condition-inherited} を適用すると、
標準射
$$
L(g')^*\mathcal{G}^\bullet \to (g')^*\mathcal{G}^\bullet
$$
が Lemma \ref{spaces-perfect-lemma-single-complex-base-change-condition} の同値な条件を
満たすことを証明すれば十分である。これは Lemma
\ref{spaces-perfect-lemma-base-change-tensor} の証明で示した。
\end{proof}

\section{完全複体の構成}
\label{spaces-perfect-section-producing-perfect}

\noindent
次の補題は、完全複体を構成するための主要な技術的道具である。
この結果の後の諸版は、ネーター近似によってこれに帰着される。

\begin{lemma}
\label{spaces-perfect-lemma-perfect-direct-image}
$S$ を概型とする。$Y$ を $S$ 上のネーター代数空間とする。
$f : X \to Y$ を局所有限型かつ準分離的な代数空間の射とする。
$E \in D(\mathcal{O}_X)$ が次を満たすとする。
\begin{enumerate}
\item $E \in D^b_{\textit{Coh}}(\mathcal{O}_X)$ である。
\item $H^i(E)$ の台は $Y$ 上固有である（すべての $i$ について）。
\item $E$ は $D(f^{-1}\mathcal{O}_Y)$ の対象として有限 tor 次元をもつ。
\end{enumerate}
このとき $Rf_*E$ は $D(\mathcal{O}_Y)$ の完全対象である。
\end{lemma}

\begin{proof}
Lemma \ref{spaces-perfect-lemma-direct-image-coherent} により、$Rf_*E$ は
$D^b_{\textit{Coh}}(\mathcal{O}_Y)$ の対象である。したがって $Rf_*E$ は
擬連接である（Lemma \ref{spaces-perfect-lemma-identify-pseudo-coherent-noetherian}）。
ゆえに $Rf_*E$ が有限 tor 次元をもつことを示せば十分である。Cohomology
on Sites, Lemma \ref{sites-cohomology-lemma-perfect} を参照されたい。
Lemma \ref{spaces-perfect-lemma-tor-qc-qs} により、
$Rf_*(E) \otimes_{\mathcal{O}_Y}^\mathbf{L} \mathcal{F}$ が一様に有界な
コホモロジーをもつことを、任意の準連接層 $\mathcal{F}$（$Y$ 上）について
確かめれば十分である。$Rf_*(E)$ は上に有界
なので、上からの有界性は明らかである。$T \subset |X|$ を
$H^i(E)$ の台（すべての $i$ について）の合併とする。このとき仮定 (1), (2)
および Lemma \ref{spaces-perfect-lemma-union-closed-proper-over-base} により、$T$ は
$Y$ 上固有である。特に、準コンパクト開部分空間 $X' \subset X$ が存在し、
これは $T$ を含む。$f' = f|_{X'}$ とおくと
$Rf_*(E) = Rf'_*(E|_{X'})$ である。実際、$E$ は
$X \setminus T$ 上で零に制限される。したがって $X$ を $X'$ で置き換え、
$f$ が準コンパクトであると仮定してよい。$f$ は準分離的であるとも
仮定している。よって
$$
Rf_*(E) \otimes_{\mathcal{O}_Y}^\mathbf{L} \mathcal{F} =
Rf_*\left(E \otimes_{\mathcal{O}_X}^\mathbf{L} Lf^*\mathcal{F}\right) =
Rf_*\left(E \otimes_{f^{-1}\mathcal{O}_Y}^\mathbf{L} f^{-1}\mathcal{F}\right)
$$
である。これは Lemma \ref{spaces-perfect-lemma-cohomology-base-change} および
Cohomology on Sites, Lemma
\ref{sites-cohomology-lemma-variant-derived-pullback} による。
仮定 (3) により、複体
$E \otimes_{f^{-1}\mathcal{O}_Y}^\mathbf{L} f^{-1}\mathcal{F}$ の
コホモロジー層は、ある一定の有限区間、たとえば $[a, b]$ に入る。
これに $Rf_*$ を施したものは区間 $[a, \infty)$ にコホモロジーをもち、
これでよい。
\end{proof}

\begin{lemma}
\label{spaces-perfect-lemma-tensor-perfect}
$S$ を概型とする。$B$ を $S$ 上のネーター代数空間とする。
$f : X \to B$ を局所有限型かつ準分離的な代数空間の射とする。
$E \in D(\mathcal{O}_X)$ を完全対象とする。$\mathcal{G}^\bullet$ を、
連接 $\mathcal{O}_X$-加群からなる有界複体であって、$B$ 上平坦かつ
台が $B$ 上固有なものとする。このとき
$K = Rf_*(E \otimes^\mathbf{L}_{\mathcal{O}_X} \mathcal{G}^\bullet)$ は
$D(\mathcal{O}_B)$ の完全対象である。
\end{lemma}

\begin{proof}
Lemma \ref{spaces-perfect-lemma-perfect-direct-image} により対象 $K$ は完全である。
同補題を適用できることを確かめよう。局所的に $E$ は有限階数自由
$\mathcal{O}_X$-加群からなる有限複体と同型である。したがって局所的に
$E \otimes^\mathbf{L}_{\mathcal{O}_X} \mathcal{G}^\bullet$ は、各項が
$$
\bigoplus\nolimits_{i = a, \ldots, b} (\mathcal{G}^i)^{\oplus r_i}
$$
という形の有限複体と同型である。ここで $a, b, r_a, \ldots, r_b$ は
いくつかの整数である。このことからコホモロジー層
$H^i(E \otimes^\mathbf{L}_{\mathcal{O}_X} \mathcal{G})$ が連接であることが
直ちに従う。tor 次元に関する仮定も、$\mathcal{G}^i$ が
$f^{-1}\mathcal{O}_Y$ 上平坦であることから従う。% P09-ERR-1031
\end{proof}

\begin{lemma}
\label{spaces-perfect-lemma-ext-perfect}
$S$ を概型とする。$B$ を $S$ 上のネーター代数空間とする。
$f : X \to B$ を局所有限型かつ準分離的な代数空間の射とする。
$E \in D(\mathcal{O}_X)$ を完全対象とする。$\mathcal{G}^\bullet$ を、
連接 $\mathcal{O}_X$-加群からなる有界複体であって、$B$ 上平坦かつ
台が $B$ 上固有なものとする。このとき
$K = Rf_*R\SheafHom(E, \mathcal{G})$ は $D(\mathcal{O}_B)$ の完全対象である。
% P09-ERR-1032
\end{lemma}

\begin{proof}
$E$ は完全複体なので双対完全複体 $E^\vee$ が存在する。Cohomology on
Sites, Lemma \ref{sites-cohomology-lemma-dual-perfect-complex} を参照されたい。
$$
R\SheafHom(E, \mathcal{G}^\bullet) =
E^\vee \otimes^\mathbf{L}_{\mathcal{O}_X} \mathcal{G}^\bullet
$$
であることに注意する。したがって $K$ の完全性は
Lemma \ref{spaces-perfect-lemma-tensor-perfect} から従う。
\end{proof}






\section{Ext に対する射影公式}
\label{spaces-perfect-section-ext}

\noindent
Lemma \ref{spaces-perfect-lemma-compute-ext}（または文献中の同様の結果）は、Quot 関手、
Hilbert 概型、およびその他のモジュライ問題に対する Artin の判定条件の
一つを確かめるために必要な障害理論の性質を検証する際に有用なことがある。
$f : X \to Y$ を固有、平坦かつ有限表示な代数空間の射とし、
$E \in D(\mathcal{O}_X)$ を完全とする。ここで補題は
$$
\Ext^i_X(E, f^*\mathcal{F}) =
\Ext^i_Y((Rf_*E^\vee)^\vee, \mathcal{F})
$$
が $\mathcal{F}$（$Y$ 上準連接）に対して成り立つと述べている。
このように書くと Ext に対する射影公式のように見え、実際この結果は
Lemma \ref{spaces-perfect-lemma-cohomology-base-change} からかなり容易に従う。

\begin{lemma}
\label{spaces-perfect-lemma-compute-tensor-perfect}
仮定と記法を Lemma \ref{spaces-perfect-lemma-tensor-perfect} と同じとする。このとき、
関手的同型
$$
H^i(B, K \otimes^\mathbf{L}_{\mathcal{O}_B} \mathcal{F})
\longrightarrow
H^i(X, E \otimes^\mathbf{L}_{\mathcal{O}_X}
(\mathcal{G}^\bullet \otimes_{\mathcal{O}_X} f^*\mathcal{F}))
$$
が存在する。これは $\mathcal{F}$（$B$ 上準連接）について定まり、
境界写像と両立する（証明参照）。
\end{lemma}

\begin{proof}
まず
$$
\mathcal{G}^\bullet \otimes_{\mathcal{O}_X}^\mathbf{L} Lf^*\mathcal{F} =
\mathcal{G}^\bullet \otimes_{f^{-1}\mathcal{O}_B}^\mathbf{L} f^{-1}\mathcal{F} =
\mathcal{G}^\bullet \otimes_{f^{-1}\mathcal{O}_B} f^{-1}\mathcal{F} =
\mathcal{G}^\bullet \otimes_{\mathcal{O}_X} f^*\mathcal{F}
$$
である。第一の等号は Cohomology on Sites, Lemma
\ref{sites-cohomology-lemma-variant-derived-pullback} による。第二の等号は
$\mathcal{G}^n$ が平坦 $f^{-1}\mathcal{O}_B$-加群であることによる。
第三の等号は引き戻しの定義による。したがって
\begin{align*}
H^i(X, E \otimes^\mathbf{L}_{\mathcal{O}_X}
(\mathcal{G}^\bullet \otimes_{\mathcal{O}_X} f^*\mathcal{F}))
& =
H^i(X, E \otimes^\mathbf{L}_{\mathcal{O}_X} \mathcal{G}^\bullet
\otimes_{\mathcal{O}_X}^\mathbf{L} Lf^*\mathcal{F}) \\
& =
H^i(B,
Rf_*(E \otimes^\mathbf{L}_{\mathcal{O}_X} \mathcal{G}^\bullet
\otimes^\mathbf{L}_{\mathcal{O}_X} Lf^*\mathcal{F})) \\
& =
H^i(B,
Rf_*(E \otimes^\mathbf{L}_{\mathcal{O}_X} \mathcal{G}^\bullet)
\otimes^\mathbf{L}_{\mathcal{O}_B} \mathcal{F}) \\
& =
H^i(B, K \otimes^\mathbf{L}_{\mathcal{O}_B} \mathcal{F})
\end{align*}
を得る。第一の等号は上の議論による。第二の等号は Leray
（Cohomology on Sites, Remark \ref{sites-cohomology-remark-before-Leray}）により、
第三の等号は Lemma \ref{spaces-perfect-lemma-cohomology-base-change} による。
境界写像に関する主張の意味は次のとおりである。短完全列
$0 \to \mathcal{F}_1 \to \mathcal{F}_2 \to \mathcal{F}_3 \to 0$ が与えられると、
これらの同型は可換図式
$$
\xymatrix{
H^i(B, K \otimes^\mathbf{L}_{\mathcal{O}_B} \mathcal{F}_3)
\ar[r] \ar[d]_\delta &
H^i(X, E \otimes^\mathbf{L}_{\mathcal{O}_X}
(\mathcal{G}^\bullet \otimes_{\mathcal{O}_X} f^*\mathcal{F}_3)) \ar[d]^\delta \\
H^{i + 1}(B, K \otimes^\mathbf{L}_{\mathcal{O}_B} \mathcal{F}_1)
\ar[r] &
H^{i + 1}(X, E \otimes^\mathbf{L}_{\mathcal{O}_X}
(\mathcal{G}^\bullet \otimes_{\mathcal{O}_X} f^*\mathcal{F}_1))
}
$$
に組み込まれる。ここで境界写像は、識別三角形
$$
K \otimes^\mathbf{L}_{\mathcal{O}_B} \mathcal{F}_1 \to
K \otimes^\mathbf{L}_{\mathcal{O}_B} \mathcal{F}_2 \to
K \otimes^\mathbf{L}_{\mathcal{O}_B} \mathcal{F}_3 \to
K \otimes^\mathbf{L}_{\mathcal{O}_B} \mathcal{F}_1[1]
$$
および $D(\mathcal{O}_X)$ において短完全列
$$
0 \to
\mathcal{G}^\bullet \otimes_{\mathcal{O}_X} f^*\mathcal{F}_1 \to
\mathcal{G}^\bullet \otimes_{\mathcal{O}_X} f^*\mathcal{F}_2 \to
\mathcal{G}^\bullet \otimes_{\mathcal{O}_X} f^*\mathcal{F}_3 \to 0
$$
に付随する識別三角形から来る。この列が完全なのは
$\mathcal{G}^n$ が $B$ 上平坦だからである。表示した図式の可換性の検証は
省略する。
\end{proof}

\begin{lemma}
\label{spaces-perfect-lemma-compute-ext-perfect}
仮定と記法を Lemma \ref{spaces-perfect-lemma-ext-perfect} と同じとする。このとき、
関手的同型
$$
H^i(B, K \otimes^\mathbf{L}_{\mathcal{O}_B} \mathcal{F})
\longrightarrow
\Ext^i_{\mathcal{O}_X}(E,
\mathcal{G}^\bullet \otimes_{\mathcal{O}_X} f^*\mathcal{F})
$$
が存在する。これは $\mathcal{F}$（$B$ 上準連接）について定まり、
境界写像と両立する（証明参照）。
\end{lemma}

\begin{proof}
Lemma \ref{spaces-perfect-lemma-ext-perfect} の証明と同様に、$E^\vee$ を双対完全複体とし、
$$
K = Rf_*(E^\vee \otimes_{\mathcal{O}_X}^\mathbf{L} \mathcal{G}^\bullet)
$$
であったことを思い出す。また
$$
\Ext^i_{\mathcal{O}_X}(E,
\mathcal{G}^\bullet \otimes_{\mathcal{O}_X} f^*\mathcal{F})
=
H^i(X, E^\vee \otimes^\mathbf{L}_{\mathcal{O}_X}
(\mathcal{G}^\bullet \otimes_{\mathcal{O}_X} f^*\mathcal{F}))
$$
が $E^\vee$ の構成により成り立つので、同型の存在は
Lemma \ref{spaces-perfect-lemma-compute-tensor-perfect} を $E^\vee$ と
$\mathcal{G}^\bullet$ に適用することから従う。
境界写像に関する主張の意味は次のとおりである。短完全列
$0 \to \mathcal{F}_1 \to \mathcal{F}_2 \to \mathcal{F}_3 \to 0$ が与えられると、
これらの同型は可換図式
$$
\xymatrix{
H^i(B, K \otimes^\mathbf{L}_{\mathcal{O}_B} \mathcal{F}_3)
\ar[r] \ar[d]_\delta &
\Ext^i_{\mathcal{O}_X}(E,
\mathcal{G}^\bullet \otimes_{\mathcal{O}_X} f^*\mathcal{F}_3) \ar[d]^\delta \\
H^{i + 1}(B, K \otimes^\mathbf{L}_{\mathcal{O}_B} \mathcal{F}_1)
\ar[r] &
\Ext^{i + 1}_{\mathcal{O}_X}(E,
\mathcal{G}^\bullet \otimes_{\mathcal{O}_X} f^*\mathcal{F}_1)
}
$$
に組み込まれる。ここで境界写像は、識別三角形
$$
K \otimes^\mathbf{L}_{\mathcal{O}_B} \mathcal{F}_1 \to
K \otimes^\mathbf{L}_{\mathcal{O}_B} \mathcal{F}_2 \to
K \otimes^\mathbf{L}_{\mathcal{O}_B} \mathcal{F}_3 \to
K \otimes^\mathbf{L}_{\mathcal{O}_B} \mathcal{F}_1[1]
$$
および $D(\mathcal{O}_X)$ において短完全列
$$
0 \to
\mathcal{G}^\bullet \otimes_{\mathcal{O}_X} f^*\mathcal{F}_1 \to
\mathcal{G}^\bullet \otimes_{\mathcal{O}_X} f^*\mathcal{F}_2 \to
\mathcal{G}^\bullet \otimes_{\mathcal{O}_X} f^*\mathcal{F}_3 \to 0
$$
に付随する識別三角形から来る。この列が完全なのは
$\mathcal{G}^n$ が $B$ 上平坦だからである。表示した図式の可換性の検証は
省略する。
\end{proof}

\begin{lemma}
\label{spaces-perfect-lemma-compute-ext}
$S$ を概型とする。$f : X \to B$ を $S$ 上の代数空間の射とし、
$E \in D(\mathcal{O}_X)$ とする。また $\mathcal{F}^\bullet$ を
$\mathcal{O}_X$-加群の複体とする。次を仮定する。% P09-ERR-1033
\begin{enumerate}
\item $B$ はネーター的である。
\item $f$ は局所有限型かつ準分離的である。
\item $E \in D^-_{\textit{Coh}}(\mathcal{O}_X)$ である。
\item $\mathcal{G}^\bullet$ は、連接 $\mathcal{O}_X$-加群からなる
有界複体であって、$B$ 上平坦かつ台が $B$ 上固有なものである。
\end{enumerate}
このとき次の二つの主張が成り立つ。
\begin{enumerate}
\item[(A)] すべての $m \in \mathbf{Z}$ に対して、完全対象
$K$（$D(\mathcal{O}_B)$ の）と関手的写像
$$
\alpha^i_\mathcal{F} :
\Ext^i_{\mathcal{O}_X}(E,
\mathcal{G}^\bullet \otimes_{\mathcal{O}_X} f^*\mathcal{F})
\longrightarrow
H^i(B, K \otimes^\mathbf{L}_{\mathcal{O}_B} \mathcal{F})
$$
が存在する。これらは $\mathcal{F}$（$B$ 上準連接）について定まり、
境界写像と両立し（証明参照）、$\alpha^i_\mathcal{F}$ は $i \leq m$ に
対して同型である。
\item[(B)] 擬連接な $L \in D(\mathcal{O}_B)$ と関手的同型
$$
\Ext^i_{\mathcal{O}_B}(L, \mathcal{F}) \longrightarrow
\Ext^i_{\mathcal{O}_X}(E,
\mathcal{G}^\bullet \otimes_{\mathcal{O}_X} f^*\mathcal{F})
$$
が存在する。これらは $\mathcal{F}$（$B$ 上準連接）について定まり、
境界写像と両立する。
\end{enumerate}
\end{lemma}

\begin{proof}
(A) の証明。$\mathcal{G}^i$ が零でないのは $i \in [a, b]$ のときだけと
する。$X$ を $\mathcal{G}^i$ の台の合併の準コンパクト開近傍で置き換えて
よい。したがって $X$ はネーター的であると仮定してよい。この場合、
$X$ と $f$ は準コンパクトかつ準分離的である。
近似 $P \to E$ を選ぶ。これは完全複体 $P$ による
$(X, E, -m - 1 + a)$ の近似である
（Theorem \ref{spaces-perfect-theorem-approximation} により可能である）。すると誘導写像
$$
\Ext^i_{\mathcal{O}_X}(E,
\mathcal{G}^\bullet \otimes_{\mathcal{O}_X} f^*\mathcal{F})
\longrightarrow
\Ext^i_{\mathcal{O}_X}(P,
\mathcal{G}^\bullet \otimes_{\mathcal{O}_X} f^*\mathcal{F})
$$
は $i \leq m$ に対して同型である。実際、この写像の核と余核はそれぞれ
$$
\Ext^i_{\mathcal{O}_X}(C,
\mathcal{G}^\bullet \otimes_{\mathcal{O}_X} f^*\mathcal{F})
\quad\text{および}\quad
\Ext^{i + 1}_{\mathcal{O}_X}(C,
\mathcal{G}^\bullet \otimes_{\mathcal{O}_X} f^*\mathcal{F})
$$
の商と部分加群である。ここで $C$ は $P \to E$ の錐である。$C$ の
コホモロジー層は次数 $\geq -m - 1 + a$ で消えるので、Derived Categories,
Lemma \ref{derived-lemma-negative-exts} によりこれらの $\Ext$-群は
$i \leq m + 1$ に対して零である。これにより $E$ が完全複体の場合へ帰着し、
これは Lemma \ref{spaces-perfect-lemma-compute-ext-perfect} である。境界写像に関する主張は
Lemma \ref{spaces-perfect-lemma-compute-ext-perfect} の証明で説明されている。

\medskip\noindent
(B) の証明。(A) の証明と同様に、$X$ はネーター的であると仮定してよい。
Lemma \ref{spaces-perfect-lemma-identify-pseudo-coherent-noetherian} により $E$ は擬連接で
あることに注意する。Lemma \ref{spaces-perfect-lemma-pseudo-coherent-hocolim} により
$E = \text{hocolim} E_n$ と書ける。ここで $E_n$ は完全で、$E_n \to E$ は
切断 $\tau_{\geq -n}$ 上の同型を誘導する。$E_n^\vee$ を双対完全複体とする
（Cohomology on Sites, Lemma
\ref{sites-cohomology-lemma-dual-perfect-complex}）。完全対象の逆系
$\ldots \to E_3^\vee \to E_2^\vee \to E_1^\vee$ を得る。これはさらに逆系
$$
\ldots \to K_3 \to K_2 \to K_1\quad\text{with}\quad
K_n = Rf_*(E_n^\vee \otimes_{\mathcal{O}_X}^\mathbf{L} \mathcal{G}^\bullet)
$$
を与える。これは $Y$ 上完全である。Lemma \ref{spaces-perfect-lemma-tensor-perfect} を
参照されたい。% P09-ERR-1034
Lemma \ref{spaces-perfect-lemma-compute-ext-perfect} とその証明、および前段落の議論
（$P = E_n$ とする）により、任意の準連接 $\mathcal{F}$（$Y$ 上）に対して
関手的な標準写像
$$
\xymatrix{
& \Ext^i_{\mathcal{O}_X}(E,
\mathcal{G}^\bullet \otimes_{\mathcal{O}_X} f^*\mathcal{F})
\ar[ld] \ar[rd] \\
H^i(Y, K_{n + 1} \otimes_{\mathcal{O}_Y}^\mathbf{L} \mathcal{F})
\ar[rr] & &
H^i(Y, K_n \otimes_{\mathcal{O}_Y}^\mathbf{L} \mathcal{F})
}
$$
を得る。これらは $i \leq n + a$ に対して同型である。
$L_n = K_n^\vee$ を双対完全複体とする。このとき
$L_1 \to L_2 \to L_3 \to \ldots$ は $D(\mathcal{O}_Y)$ における完全対象の
系であり、任意の準連接 $\mathcal{F}$（$Y$ 上）に対して写像
$$
\Ext^i_{\mathcal{O}_Y}(L_{n + 1}, \mathcal{F})
\longrightarrow
\Ext^i_{\mathcal{O}_Y}(L_n, \mathcal{F})
$$
は $i \leq n + a - 1$ に対して同型である。これは
$L_n \to L_{n + 1}$ が切断 $\tau_{\geq -n - a + 2}$ 上の同型を誘導することを
意味する（ヒント：$L_n \to L_{n + 1}$ の錐を取り、その最後の非消滅
コホモロジー層を見る）。したがって $L = \text{hocolim} L_n$ は擬連接である。
Lemma \ref{spaces-perfect-lemma-pseudo-coherent-hocolim} を参照されたい。ホモトピー余極限の
写像性質により
$\Ext^i_{\mathcal{O}_Y}(L, \mathcal{F}) =
\Ext^i_{\mathcal{O}_Y}(L_n, \mathcal{F})$ が $i \leq n + a - 3$ に対して
成り立ち、証明が完了する。
\end{proof}

\begin{remark}
\label{spaces-perfect-remark-base-change-of-L}
Lemma \ref{spaces-perfect-lemma-compute-ext} の部分 (B) の擬連接複体 $L$ は、この状況に
標準的に付随する。たとえば、(B) における $L$ の形成は基底変換と両立する。
言い換えると、カルテシアン図式
$$
\xymatrix{
X' \ar[r]_{g'} \ar[d]_{f'} &
X \ar[d]^f \\
Y' \ar[r]^g &
Y
}
$$
が概型について与えられると、関手的な標準同型
$$
\Ext^i_{\mathcal{O}_{Y'}}(Lg^*L, \mathcal{F}') \longrightarrow
\Ext^i_{\mathcal{O}_X}(L(g')^*E,
(g')^*\mathcal{G}^\bullet \otimes_{\mathcal{O}_{X'}} (f')^*\mathcal{F}')
$$
を得る。ここで $\mathcal{F}'$ は $Y'$ 上準連接である。右辺の
$\mathcal{G}^\bullet$ には導来引き戻しを {\bf 用いていない}ことに注意する。
% P09-ERR-1035 P09-ERR-1036
この結果が必要になれば、ここで精密な結果を定式化して詳しい証明を与える。
\end{remark}








\section{極限と導来圏}
\label{spaces-perfect-section-limits}

\noindent
この節では、代数空間の逆系の極限である代数空間の導来圏に関する
いくつかの結果を集める。より正確には、以下の設定で議論する。

\begin{situation}
\label{spaces-perfect-situation-descent}
$S$ を概型とする。$X = \lim_{i \in I} X_i$ を、$S$ 上の代数空間の
有向系の極限とし、その遷移射
$f_{i'i} : X_{i'} \to X_i$ はアフィンであるとする。射影を
$f_i : X \to X_i$ と書く。$X_i$ は準コンパクトかつ準分離的であると
仮定する（すべての $i \in I$ について）。また元 $0 \in I$ を一つ選ぶ。
\end{situation}

\begin{lemma}
\label{spaces-perfect-lemma-descend-homomorphisms}
Situation \ref{spaces-perfect-situation-descent} のもとで、$E_0$ と $K_0$ を
$D(\mathcal{O}_{X_0})$ の対象とする。$E_i = Lf_{i0}^*E_0$ および
$K_i = Lf_{i0}^*K_0$ とおき（$i \geq 0$ について）、
$E = Lf_0^*E_0$ および $K = Lf_0^*K_0$ とおく。このとき写像
$$
\colim_{i \geq 0} \Hom_{D(\mathcal{O}_{X_i})}(E_i, K_i)
\longrightarrow
\Hom_{D(\mathcal{O}_X)}(E, K)
$$
は、次のいずれかが成り立てば同型である。
\begin{enumerate}
\item $E_0$ は完全で、$K_0 \in D_\QCoh(\mathcal{O}_{X_0})$ である。
\item $E_0$ は擬連接で、$K_0 \in D_\QCoh(\mathcal{O}_{X_0})$ は
有限 tor 次元をもつ。
\end{enumerate}
\end{lemma}

\begin{proof}
準コンパクトかつ準分離的な各対象 $U_0$（$(X_0)_{spaces, \etale}$ の）に対し、
条件 $P$ を考える。これは標準写像
$$
\colim_{i \geq 0} \Hom_{D(\mathcal{O}_{U_i})}(E_i|_{U_i}, K_i|_{U_i})
\longrightarrow
\Hom_{D(\mathcal{O}_U)}(E|_U, K|_U)
$$
が同型であるという条件である。ここで
$U = X \times_{X_0} U_0$ および $U_i = X_i \times_{X_0} U_0$ である。
Lemma \ref{spaces-perfect-lemma-induction-principle} の帰納原理により、$P$ が成り立つことを
各 $U_0$ に対して示す。同補題の条件 (2) は、導来圏における hom の
Mayer--Vietoris から直ちに従う。Lemma
\ref{spaces-perfect-lemma-mayer-vietoris-hom} を参照されたい。したがって $X_0$ が
アフィンの場合に補題を証明すれば十分である。

\medskip\noindent
$X_0$ がアフィンならば、結果は概型の場合から従う。Derived Categories of
Schemes, Lemma \ref{perfect-lemma-descend-homomorphisms} を参照されたい。
これを見るには Lemma \ref{spaces-perfect-lemma-derived-quasi-coherent-small-etale-site} の
同値、および Lemmas \ref{spaces-perfect-lemma-descend-pseudo-coherent},
\ref{spaces-perfect-lemma-descend-tor-amplitude}, and
\ref{spaces-perfect-lemma-descend-perfect} で説明された性質の翻訳を用いる。
\end{proof}

\begin{lemma}
\label{spaces-perfect-lemma-perfect-on-limit}
Situation \ref{spaces-perfect-situation-descent} のもとで、$D(\mathcal{O}_X)$ の完全対象の圏は
$D(\mathcal{O}_{X_i})$ の完全対象の諸圏の余極限である。
\end{lemma}

\begin{proof}
準コンパクトかつ準分離的な各対象 $U_0$（$(X_0)_{spaces, \etale}$ の）に対し、
条件 $P$ を考える。これは関手
$$
\colim_{i \geq 0} D_{perf}(\mathcal{O}_{U_i})
\longrightarrow
D_{perf}(\mathcal{O}_U)
$$
が同値であるという条件である。ここで ${}_{perf}$ は完全対象からなる
充満部分圏を表し、$U = X \times_{X_0} U_0$ および
$U_i = X_i \times_{X_0} U_0$ である。Lemma
\ref{spaces-perfect-lemma-induction-principle} の帰納原理により、$P$ が成り立つことを
すべての $U_0$ に対して示す。まず Lemma
\ref{spaces-perfect-lemma-descend-homomorphisms} によって関手が忠実充満であることはすでに
分かっている。したがって本質的全射性を示せば十分である。

\medskip\noindent
まず帰納原理の条件 (2) を確かめる。したがって基本識別正方形
$(U_0 \subset X_0, V_0 \to X_0)$ があり、$P$ が $U_0$, $V_0$, および
$U_0 \times_{X_0} V_0$ に対して成り立つと仮定する。
$E$ を $D(\mathcal{O}_X)$ の完全対象とする。ある $i \geq 0$ と、完全対象
$E_{U, i}$（$U_i$ 上）および完全対象 $E_{V, i}$（$V_i$ 上）を取れて、これらの
$U$ と $V$ への引き戻しはそれぞれ $E|_U$ と $E|_V$ に同型である。
写像
$$
a : E_{U, i} \to (R(X \to X_i)_*E)|_{U_i}
\quad\text{and}\quad
b : E_{V, i} \to (R(X \to X_i)_*E)|_{V_i}
$$
を考える。これらは同型 $L(U \to U_i)^*E_{U, i} \to E|_U$ および
$L(V \to V_i)^*E_{V, i} \to E|_V$ に随伴する写像である。
忠実充満性により、$i$ を大きくした後、同型
$c : E_{U, i}|_{U_i \times_{X_i} V_i} \to E_{V, i}|_{U_i \times_{X_i} V_i}$
であって、引き戻すと同一視
$$
L(U \to U_i)^*E_{U, i}|_{U \times_X V} \to E|_{U \times_X V} \to
L(V \to V_i)^*E_{V, i}|_{U \times_X V}.
$$
になるものを取れる。Lemma \ref{spaces-perfect-lemma-glue} を適用すると、対象
$E_i$（$X_i$ 上）と写像 $d : E_i \to R(X \to X_i)_*E$ を得て、これは写像
$a$ および $b$ に制限される（それぞれ $U_i$ および $V_i$ 上で）。このとき $E_i$ が完全で、
$d$ が同型 $L(X \to X_i)^*E_i \to E$ に随伴することは明らかである。

\medskip\noindent
最後に帰納原理の条件 (1)、すなわち $X_0$ がアフィンの場合に補題が
成り立つことを確かめる。これは概型の場合から従う。Derived Categories of
Schemes, Lemma \ref{perfect-lemma-descend-perfect} を参照されたい。
これを見るには Lemma \ref{spaces-perfect-lemma-derived-quasi-coherent-small-etale-site} の
同値、および Lemma \ref{spaces-perfect-lemma-descend-perfect} の翻訳を用いる。
\end{proof}






\section{コホモロジーと基底変換 VI}
\label{spaces-perfect-section-cohomology-and-base-change-final}

\noindent
コホモロジーと基底変換に関する最後の節であり、Sections
\ref{spaces-perfect-section-cohomology-and-base-change-perfect},
\ref{spaces-perfect-section-cohomology-base-change}, and
\ref{spaces-perfect-section-producing-perfect} の議論を続ける。理解しやすい特別な場合は
Remark \ref{spaces-perfect-remark-explain-perfect-direct-image} に与えられている。

\begin{lemma}
\label{spaces-perfect-lemma-base-change-tensor-perfect}
$S$ を概型とする。$f : X \to Y$ を、$S$ 上の代数空間の間の有限表示射とする。
$E \in D(\mathcal{O}_X)$ を完全対象とする。$\mathcal{G}^\bullet$ を
有限表示 $\mathcal{O}_X$-加群からなる有界複体であって、$Y$ 上平坦かつ
台が $Y$ 上固有なものとする。このとき
$$
K = Rf_*(E \otimes_{\mathcal{O}_X}^\mathbf{L} \mathcal{G}^\bullet)
$$
は $D(\mathcal{O}_Y)$ の完全対象であり、その形成は任意の基底変換と可換する。
\end{lemma}

\begin{proof}
基底変換に関する主張は Lemma \ref{spaces-perfect-lemma-base-change-tensor} である。
したがって $K$ が完全対象であることを示せば十分である。$Y$ がネーター的
ならば、これは Lemma \ref{spaces-perfect-lemma-tensor-perfect} から従う。ネーター近似により
この場合へ帰着する。読者には以下の証明を読み飛ばすことを勧める。

\medskip\noindent
問題は $Y$ 上局所的なので、$Y$ はアフィンであると仮定してよい。
$Y = \Spec(R)$ と書く。$R = \colim R_i$ をネーター環 $R_i$ のフィルター
余極限として書く。Limits of Spaces, Lemma
\ref{spaces-limits-lemma-descend-finite-presentation} により、ある $i$ と
代数空間 $X_i$（$R_i$ 上有限表示）が存在し、その $R$ への基底変換は
$X$ である。Limits of Spaces, Lemma
\ref{spaces-limits-lemma-descend-modules-finite-presentation} により、$i$ を
大きくした後、有限表示 $\mathcal{O}_{X_i}$-加群からなる有界複体
$\mathcal{G}_i^\bullet$ で、その $X$ への引き戻しが
$\mathcal{G}^\bullet$ であるものが存在すると仮定してよい。さらに $i$ を
大きくした後、$\mathcal{G}_i^n$ は $R_i$ 上平坦であると仮定してよい。
Limits of Spaces, Lemma \ref{spaces-limits-lemma-descend-flat} を参照されたい。
さらに $i$ を大きくした後、$\mathcal{G}_i^n$ の台は $R_i$ 上固有であると
仮定してよい。Limits of Spaces, Lemma
\ref{spaces-limits-lemma-eventually-proper-support} を参照されたい。
最後に Lemma \ref{spaces-perfect-lemma-perfect-on-limit} により、$i$ を大きくした後、
完全対象 $E_i$（$D(\mathcal{O}_{X_i})$ の）で、その $X$ への引き戻しが
$E$ であるものが存在すると仮定してよい。Lemma
\ref{spaces-perfect-lemma-tensor-perfect} により
$$
K_i =
Rf_{i, *}(E_i \otimes_{\mathcal{O}_{X_i}}^\mathbf{L} \mathcal{G}_i^\bullet)
$$
は $\Spec(R_i)$ 上完全である。ここで $f_i : X_i \to \Spec(R_i)$ は
構造射である。基底変換の結果（Lemma \ref{spaces-perfect-lemma-base-change-tensor}）により、
$K_i$ の $Y = \Spec(R)$ への引き戻しは $K$ であり、結論を得る。
\end{proof}

\begin{remark}
\label{spaces-perfect-remark-explain-perfect-direct-image}
$R$ を環とする。$X$ を $R$ 上有限表示な代数空間とする。
$\mathcal{G}$ を有限表示 $\mathcal{O}_X$-加群であって、$R$ 上平坦かつ
台が $R$ 上固有なものとする。Lemma
\ref{spaces-perfect-lemma-base-change-tensor-perfect} により、有限生成射影 $R$-加群からなる
有限複体 $M^\bullet$ で
$$
R\Gamma(X_{R'}, \mathcal{G}_{R'}) = M^\bullet \otimes_R R'
$$
となるものが存在し、これは $R$-代数 $R'$ について関手的である。
\end{remark}

\begin{lemma}
\label{spaces-perfect-lemma-base-change-tensor-pseudo-coherent}
$S$ を概型とする。$f : X \to Y$ を、$S$ 上の代数空間の間の有限表示射とする。
$E \in D(\mathcal{O}_X)$ を擬連接対象とする。$\mathcal{G}^\bullet$ を
有限表示 $\mathcal{O}_X$-加群からなる上に有界な複体であって、$Y$ 上平坦
かつ台が $Y$ 上固有なものとする。このとき
$$
K = Rf_*(E \otimes_{\mathcal{O}_X}^\mathbf{L} \mathcal{G}^\bullet)
$$
は $D(\mathcal{O}_Y)$ の擬連接対象であり、その形成は任意の基底変換と
可換する。
\end{lemma}

\begin{proof}
基底変換に関する主張は Lemma \ref{spaces-perfect-lemma-base-change-tensor} である。
したがって $K$ が擬連接対象であることを示せば十分である。これは完全複体に
よる近似を用いて Lemma \ref{spaces-perfect-lemma-base-change-tensor-perfect} から従う。
読者には以下の証明を読み飛ばすことを勧める。

\medskip\noindent
問題は $Y$ 上エタール局所的なので、$Y$ はアフィンであると仮定してよい。
すると $X$ は準コンパクトかつ準分離的である。さらに、ある整数 $N$ が存在し、
全導来直像 $Rf_* : D_\QCoh(\mathcal{O}_X) \to D_\QCoh(\mathcal{O}_Y)$ は
コホモロジー次元 $N$ をもつ。Lemma
\ref{spaces-perfect-lemma-quasi-coherence-direct-image} の説明を参照されたい。
整数 $b$ を選び、$\mathcal{G}^i = 0$ が $i > b$ に対して成り立つようにする。
$K$ が $m$-擬連接であることをすべての $m$ に対して示せば十分である。
近似 $P \to E$ を選ぶ。これは完全複体 $P$ による
$(X, E, m - N - 1 - b)$ の近似である。これは Theorem
\ref{spaces-perfect-theorem-approximation} により可能である。識別三角形
$$
P \to E \to C \to P[1]
$$
を $D_\QCoh(\mathcal{O}_X)$ において選ぶ。$C$ のコホモロジー層は次数
$\geq m - N - 1 - b$ で零である。
したがって $C \otimes^\mathbf{L} \mathcal{G}^\bullet$ のコホモロジー層は
次数 $\geq m - N - 1$ で零である。ゆえに
$Rf_*(C \otimes^\mathbf{L} \mathcal{G})$ のコホモロジー層は
次数 $\geq m - 1$ で零である。% P09-ERR-1037
したがって
$$
Rf_*(P \otimes^\mathbf{L} \mathcal{G}) \to
Rf_*(E \otimes^\mathbf{L} \mathcal{G})
$$
は次数 $\geq m$ のコホモロジー層上で同型である。
次に $H^i(P) = 0$ が $i > a$ に対して成り立つと仮定する。このとき
$
P \otimes^\mathbf{L} \sigma_{\geq m - N - 1 - a}\mathcal{G}^\bullet
\longrightarrow
P \otimes^\mathbf{L} \mathcal{G}^\bullet
$
は次数 $\geq m - N - 1$ のコホモロジー層上で同型である。したがって再び
$$
Rf_*(P \otimes^\mathbf{L} \sigma_{\geq m - N - 1 - a}\mathcal{G}^\bullet) \to
Rf_*(P \otimes^\mathbf{L} \mathcal{G}^\bullet)
$$
が次数 $\geq m$ のコホモロジー層上で同型であると分かる。Lemma
\ref{spaces-perfect-lemma-base-change-tensor-perfect} により左辺は完全複体である。
したがって望みどおり $K$ は $m$-擬連接である。
\end{proof}

\begin{lemma}
\label{spaces-perfect-lemma-flat-proper-perfect-direct-image-general}
$S$ を概型とする。$f : X \to Y$ を、$S$ 上の代数空間の固有な有限表示射とする。
\begin{enumerate}
\item $E \in D(\mathcal{O}_X)$ を完全とし、$f$ を平坦とする。このとき
$Rf_*E$ は $D(\mathcal{O}_Y)$ の完全対象であり、その形成は任意の基底変換と
可換する。
\item $\mathcal{G}$ を有限表示 $\mathcal{O}_X$-加群であって、$S$ 上平坦な
ものとする。このとき $Rf_*\mathcal{G}$ は $D(\mathcal{O}_Y)$ の完全対象で
あり、その形成は任意の基底変換と可換する。% P09-ERR-1038
\end{enumerate}
\end{lemma}

\begin{proof}
Lemma \ref{spaces-perfect-lemma-base-change-tensor-perfect} の特別な場合である。(1) では
$\mathcal{G}^\bullet$ を $\mathcal{O}_X$（次数 $0$ に置いたもの）とし、(2) では
$E = \mathcal{O}_X$ として $\mathcal{G}^\bullet$ を
$\mathcal{G}$（次数 $0$ に置いたもの）からなるものとする。
\end{proof}

\begin{lemma}
\label{spaces-perfect-lemma-flat-proper-pseudo-coherent-direct-image-general}
$S$ を概型とする。$f : X \to Y$ を、$S$ 上の代数空間の平坦かつ固有な
有限表示射とする。$E \in D(\mathcal{O}_X)$ を擬連接とする。このとき
$Rf_*E$ は $D(\mathcal{O}_Y)$ の擬連接対象であり、その形成は任意の
基底変換と可換する。
\end{lemma}

\noindent
より一般に、$f : X \to Y$ が固有で、$E$（$X$ 上）が $Y$ に関して擬連接
（More on Morphisms of Spaces, Definition
\ref{spaces-more-morphisms-definition-relative-pseudo-coherence}）ならば、
$Rf_*E$ は擬連接である（ただし、この一般性では形成は基底変換と可換しない）。
この概型の場合は \cite{Kiehl} で証明されている。

\begin{proof}
Lemma \ref{spaces-perfect-lemma-base-change-tensor-pseudo-coherent} を
$\mathcal{G} = \mathcal{O}_X$ として適用した特別な場合である。
\end{proof}

\begin{lemma}
\label{spaces-perfect-lemma-pullback-and-limits}
$R$ を環とする。$X$ を代数空間とし、$f : X \to \Spec(R)$ を固有、平坦かつ
有限表示とする。$(M_n)$ を、全射な遷移写像をもつ $R$-加群の逆系とする。
このとき標準写像
$$
\mathcal{O}_X \otimes_R (\lim M_n)
\longrightarrow
\lim \mathcal{O}_X \otimes_R M_n
$$
は、左辺から右辺に $DQ_X$ を適用したものへの同型を誘導する。
\end{lemma}

\begin{proof}
この主張の意味は、任意の対象 $E$（$D_\QCoh(\mathcal{O}_X)$ の）に対して、
誘導写像
$$
\Hom(E, \mathcal{O}_X \otimes_R (\lim M_n))
\longrightarrow
\Hom(E, \lim \mathcal{O}_X \otimes_R M_n)
$$
が同型であるということである。$D_\QCoh(\mathcal{O}_X)$ は完全生成子をもつ
（Theorem \ref{spaces-perfect-theorem-bondal-van-den-Bergh}）ので、完全な $E$ について
確かめれば十分である。Lemma \ref{spaces-perfect-lemma-Rlim-quasi-coherent} により
$\lim \mathcal{O}_X \otimes_R M_n =
R\lim \mathcal{O}_X \otimes_R M_n$ である。Cohomology on Sites, Section
\ref{sites-cohomology-section-global-RHom} の完全関手
$R\Hom_X(E, -) : D_\QCoh(\mathcal{O}_X) \to D(R)$ は積と可換し、したがって
導来極限とも可換する。ゆえに
$$
R\Hom_X(E, \lim \mathcal{O}_X \otimes_R M_n) =
R\lim R\Hom_X(E, \mathcal{O}_X \otimes_R M_n)
$$
$E^\vee$ を双対完全複体とする。Cohomology on Sites, Lemma
\ref{sites-cohomology-lemma-dual-perfect-complex} を参照されたい。Lemma
\ref{spaces-perfect-lemma-cohomology-base-change} により
$$
R\Hom_X(E, \mathcal{O}_X \otimes_R M_n) =
R\Gamma(X, E^\vee \otimes_{\mathcal{O}_X}^\mathbf{L} Lf^*M_n) =
R\Gamma(X, E^\vee) \otimes_R^\mathbf{L} M_n
$$
である。Lemma \ref{spaces-perfect-lemma-flat-proper-perfect-direct-image-general} により
$R\Gamma(X, E^\vee)$ は完全な $R$-加群複体であると分かる。特にこれは
擬連接複体なので、More on Algebra, Lemma
\ref{more-algebra-lemma-pseudo-coherent-tensor-limit} により
$$
R\lim R\Gamma(X, E^\vee) \otimes_R^\mathbf{L} M_n =
R\Gamma(X, E^\vee) \otimes_R^\mathbf{L} \lim M_n
$$
を得て、望む結果が従う。
\end{proof}

\begin{lemma}
\label{spaces-perfect-lemma-perfect-enough}
$A$ を環とする。$X$ を $A$ 上の準コンパクトかつ準分離的な代数空間とする。
$K \in D^-_\QCoh(\mathcal{O}_X)$ とする。
$R\Gamma(X, E \otimes^\mathbf{L} K)$ が $D(A)$ において擬連接であることを、
すべての完全な $E$（$D(\mathcal{O}_X)$ の）について仮定する。このとき
$R\Gamma(X, E \otimes^\mathbf{L} K)$ は $D(A)$ において擬連接である
（すべての擬連接な $E$（$D(\mathcal{O}_X)$ の）について）。
\end{lemma}

\begin{proof}
ある整数 $N$ が存在し、$R\Gamma(X, -) : D_\QCoh(\mathcal{O}_X) \to D(A)$ は
コホモロジー次元 $N$ をもつ。Lemma
\ref{spaces-perfect-lemma-quasi-coherence-direct-image} の説明を参照されたい。
$b \in \mathbf{Z}$ を、$H^i(K) = 0$ が $i > b$ に対して成り立つように選ぶ。
$E$ を $X$ 上擬連接とする。$R\Gamma(X, E \otimes^\mathbf{L} K)$ が
$m$-擬連接であることをすべての $m$ に対して示せば十分である。
近似 $P \to E$ を選ぶ。これは完全複体 $P$ による
$(X, E, m - N - 1 - b)$ の近似である。これは Theorem
\ref{spaces-perfect-theorem-approximation} により可能である。識別三角形
$$
P \to E \to C \to P[1]
$$
を $D_\QCoh(\mathcal{O}_X)$ において選ぶ。$C$ のコホモロジー層は次数
$\geq m - N - 1 - b$ で零である。したがって
$C \otimes^\mathbf{L} K$ のコホモロジー層は次数 $\geq m - N - 1$ で零で
ある。ゆえに $R\Gamma(X, C \otimes^\mathbf{L} K)$ のコホモロジーは次数
$\geq m - 1$ で零である。したがって
$$
R\Gamma(X, P \otimes^\mathbf{L} K) \to R\Gamma(X, E \otimes^\mathbf{L} K)
$$
は次数 $\geq m$ のコホモロジー上で同型である。仮定により左辺は
擬連接である。したがって望みどおり
$R\Gamma(X, E \otimes^\mathbf{L} K)$ は $m$-擬連接である。
\end{proof}

\begin{lemma}
\label{spaces-perfect-lemma-base-change-RHom-perfect}
$S$ を概型とする。$f : X \to Y$ を、$S$ 上の代数空間の間の有限表示射とする。
$E \in D(\mathcal{O}_X)$ を完全対象とする。$\mathcal{G}^\bullet$ を
有限表示 $\mathcal{O}_X$-加群からなる有界複体であって、$Y$ 上平坦かつ
台が $Y$ 上固有なものとする。このとき
$$
K = Rf_*R\SheafHom(E, \mathcal{G}^\bullet)
$$
は $D(\mathcal{O}_Y)$ の完全対象であり、その形成は任意の基底変換と可換する。
\end{lemma}

\begin{proof}
基底変換に関する主張は Lemma \ref{spaces-perfect-lemma-base-change-RHom} である。
したがって $K$ が完全対象であることを示せば十分である。$Y$ がネーター的
ならば、これは Lemma \ref{spaces-perfect-lemma-ext-perfect} から従う。ネーター近似により
この場合へ帰着する。読者には以下の証明を読み飛ばすことを勧める。

\medskip\noindent
問題は $Y$ 上局所的なので、$Y$ はアフィンであると仮定してよい。
$Y = \Spec(R)$ と書く。$R = \colim R_i$ をネーター環 $R_i$ のフィルター
余極限として書く。Limits of Spaces, Lemma
\ref{spaces-limits-lemma-descend-finite-presentation} により、ある $i$ と
代数空間 $X_i$（$R_i$ 上有限表示）が存在し、その $R$ への基底変換は
$X$ である。Limits of Spaces, Lemma
\ref{spaces-limits-lemma-descend-modules-finite-presentation} により、$i$ を
大きくした後、有限表示 $\mathcal{O}_{X_i}$-加群からなる有界複体
$\mathcal{G}_i^\bullet$ で、その $X$ への引き戻しが $\mathcal{G}$ であるものが
存在すると仮定してよい。% P09-ERR-1039
さらに $i$ を大きくした後、$\mathcal{G}_i^n$ は $R_i$ 上平坦であると
仮定してよい。Limits of Spaces, Lemma
\ref{spaces-limits-lemma-descend-flat} を参照されたい。さらに $i$ を
大きくした後、$\mathcal{G}_i^n$ の台は $R_i$ 上固有であると仮定してよい。
Limits of Spaces, Lemma
\ref{spaces-limits-lemma-eventually-proper-support} を参照されたい。
最後に Lemma \ref{spaces-perfect-lemma-descend-perfect} により、$i$ を大きくした後、
完全対象 $E_i$（$D(\mathcal{O}_{X_i})$ の）で、その $X$ への引き戻しが
$E$ であるものが存在すると仮定してよい。Lemma
\ref{spaces-perfect-lemma-compute-ext-perfect} を $X_i \to \Spec(R_i)$, $E_i$,
$\mathcal{G}_i^\bullet$ に適用し、すでに示した基底変換性を用いると結果を得る。
\end{proof}








\section{完全複体}
\label{spaces-perfect-section-perfect-complexes}

\noindent
まず完全複体の Betti 数の飛躍軌跡について論じる。そのため、最初に
Betti 数を定義しなければならない。

\medskip\noindent
$S$ を概型とする。$X$ を $S$ 上の代数空間とする。$E$ を
$D(\mathcal{O}_X)$ の対象とする。$x \in |X|$ とする。ここで
$\beta_i(x) \in \{0, 1, 2, \ldots \} \cup \{\infty\}$ を定義したい。
そのため、射 $f : \Spec(k) \to X$ を選ぶ。これは $x$ の同値類に属する。
このとき
$Lf^*E$ は $D(\Spec(k)_\etale, \mathcal{O})$ の対象である。\'Etale
Cohomology, Lemma
\ref{etale-cohomology-lemma-all-modules-quasi-coherent} および Theorem
\ref{etale-cohomology-theorem-quasi-coherent} により、
$D(\Spec(k)_\etale, \mathcal{O}) = D(k)$ は $k$-ベクトル空間の導来圏である。
したがって $Lf^*E$ は $k$-ベクトル空間の複体であり、
$\beta_i(x) = \dim_k H^i(Lf^*E)$ とおける。これが $x$ に属する代表の選択に
依存しないことは容易に分かる。さらに、$X$ が概型ならば、これは Derived
Categories of Schemes, Section \ref{perfect-section-perfect-complexes} で
用いられた概念と同じである。

\begin{lemma}
\label{spaces-perfect-lemma-jump-loci}
$S$ を概型とする。$X$ を $S$ 上の代数空間とする。
$E \in D(\mathcal{O}_X)$ を擬連接（たとえば完全）とする。任意の
$i \in \mathbf{Z}$ に対して、上で定義した関数
$$
\beta_i : |X| \longrightarrow \{0, 1, 2, \ldots\}
$$
を考える。このとき次が成り立つ。
\begin{enumerate}
\item $\beta_i$ の形成は任意の基底変換と可換する。
\item 関数 $\beta_i$ は上半連続である。
\item $\beta_i$ の水準集合はエタール局所構成可能である。
\end{enumerate}
\end{lemma}

\begin{proof}
概型 $U$ と全射エタール射 $\varphi : U \to X$ を選ぶ。このとき
$L\varphi^*E$ は概型 $U$ 上の擬連接複体である（Lemma
\ref{spaces-perfect-lemma-descend-pseudo-coherent} を用いる）。したがって概型に対する結果を
適用できる。Derived Categories of Schemes, Lemma
\ref{perfect-lemma-jump-loci} を参照されたい。部分 (3) の意味は、水準集合の
$U$ への逆像が局所構成可能であるということである。Properties of Spaces,
Definition \ref{spaces-properties-definition-locally-constructible} を
参照されたい。
\end{proof}

\begin{lemma}
\label{spaces-perfect-lemma-jump-loci-geometric}
$Y$ を概型とし、$X$ を $Y$ 上の代数空間であって、構造射
$f : X \to Y$ が平坦、固有かつ有限表示であるものとする。
$\mathcal{F}$ を有限表示 $\mathcal{O}_X$-加群であって、$Y$ 上平坦なものと
する。固定した $i \in \mathbf{Z}$ に対して関数
$$
\beta_i : |Y| \to \{0, 1, 2, \ldots\},\quad
y \longmapsto \dim_{\kappa(y)} H^i(X_y, \mathcal{F}_y)
$$
を考える。このとき次が成り立つ。
\begin{enumerate}
\item $\beta_i$ の形成は任意の基底変換と可換する。
\item 関数 $\beta_i$ は上半連続である。
\item $\beta_i$ の水準集合は $Y$ において局所構成可能である。
\end{enumerate}
\end{lemma}

\begin{proof}
コホモロジーと基底変換（より正確には Lemma
\ref{spaces-perfect-lemma-flat-proper-perfect-direct-image-general}）により、対象
$K = Rf_*\mathcal{F}$ は $Y$ の導来圏の完全対象であり、その形成は任意の
基底変換と可換する。特に
$$
H^i(X_y, \mathcal{F}_y) = H^i(K \otimes_{\mathcal{O}_Y}^\mathbf{L} \kappa(y))
$$
である。したがって補題は Lemma \ref{spaces-perfect-lemma-jump-loci} から従う。
\end{proof}

\begin{lemma}
\label{spaces-perfect-lemma-chi-locally-constant}
$S$ を概型とする。$X$ を $S$ 上の代数空間とする。
$E \in D(\mathcal{O}_X)$ を完全とする。関数
$$
\chi_E : |X| \longrightarrow \mathbf{Z},\quad
x \longmapsto \sum (-1)^i \beta_i(x)
$$
は $X$ 上局所定数である。
\end{lemma}

\begin{proof}
省略する。ヒント：エタール局所化により概型の場合から従う。Derived
Categories of Schemes, Lemma
\ref{perfect-lemma-chi-locally-constant} を参照されたい。
\end{proof}

\begin{lemma}
\label{spaces-perfect-lemma-open-where-cohomology-in-degree-i-rank-r}
$S$ を概型とする。$X$ を $S$ 上の代数空間とする。
$E \in D(\mathcal{O}_X)$ を完全とする。$i, r \in \mathbf{Z}$ が与えられると、
次で特徴づけられる開部分空間 $U \subset X$ が存在する。% P09-ERR-1040
\begin{enumerate}
\item $E|_U \cong H^i(E|_U)[-i]$ であり、$H^i(E|_U)$ は局所自由
$\mathcal{O}_U$-加群で階数 $r$ である。
\item 射 $f : Y \to X$ が $U$ を経由することと、$Lf^*E$ が階数 $r$ の
局所自由加群を次数 $i$ に置いたものと同型であることは同値である。
\end{enumerate}
\end{lemma}

\begin{proof}
省略する。ヒント：エタール局所化により概型の場合から従う。Derived
Categories of Schemes, Lemma
\ref{perfect-lemma-open-where-cohomology-in-degree-i-rank-r} を参照されたい。
\end{proof}

\begin{lemma}
\label{spaces-perfect-lemma-open-where-cohomology-in-degree-i-rank-r-geometric}
$S$ を概型とする。$f : X \to Y$ を、$S$ 上の代数空間の固有、平坦かつ
有限表示な射とする。$\mathcal{F}$ を有限表示 $\mathcal{O}_X$-加群であって、
$Y$ 上平坦なものとする。$i, r \in \mathbf{Z}$ を固定する。このとき次の
性質をもつ開部分空間 $V \subset Y$ が存在する。射 $T \to Y$ が $V$ を
経由することと、$Rf_{T, *}\mathcal{F}_T$ が階数 $r$ の有限局所自由加群を
次数 $i$ に置いたものと同型であることは同値である。
\end{lemma}

\begin{proof}
コホモロジーと基底変換（Lemma
\ref{spaces-perfect-lemma-flat-proper-perfect-direct-image-general}）により、対象
$K = Rf_*\mathcal{F}$ は $Y$ の導来圏の完全対象であり、その形成は任意の
基底変換と可換する。したがってこの補題は Lemma
\ref{spaces-perfect-lemma-open-where-cohomology-in-degree-i-rank-r} から直ちに従う。
\end{proof}

\begin{lemma}
\label{spaces-perfect-lemma-locally-closed-where-H0-locally-free}
$S$ を概型とする。$X$ を $S$ 上の代数空間とする。
$E \in D(\mathcal{O}_X)$ を、tor 振幅が $[a, b]$ に入る完全対象とする。
ここで $a, b \in \mathbf{Z}$ とする。$r \geq 0$ とする。このとき、次で
特徴づけられる局所閉部分空間 $j : Z \to X$ が存在する。
\begin{enumerate}
\item $H^a(Lj^*E)$ は局所自由 $\mathcal{O}_Z$-加群で階数 $r$ である。
\item 射 $f : Y \to X$ が $Z$ を経由することと、すべての射
$g : Y' \to Y$ に対して $\mathcal{O}_{Y'}$-加群
$H^a(L(f \circ g)^*E)$ が階数 $r$ の局所自由加群であることは同値である。
\end{enumerate}
さらに $j : Z \to X$ は有限表示であり、次が成り立つ。
\begin{enumerate}
\item[(3)] $f : Y \to X$ が $Y \xrightarrow{g} Z \to X$ と分解するならば、
$H^a(Lf^*E) = g^*H^a(Lj^*E)$ である。
\item[(4)] $\beta_a(x) \leq r$ がすべての $x \in |X|$ に対して成り立つならば、
$j$ は閉埋め込みであり、$f : Y \to X$ が与えられたとき次は同値である。
\begin{enumerate}
\item $f : Y \to X$ は $Z$ を経由する。
\item $H^0(Lf^*E)$ は局所自由 $\mathcal{O}_Y$-加群で階数 $r$ である。
\end{enumerate}
さらに $r = 1$ ならば、これらは次とも同値である。
\begin{enumerate}
\item[(c)] $\mathcal{O}_Y \to
\SheafHom_{\mathcal{O}_Y}(H^0(Lf^*E), H^0(Lf^*E))$ は単射である。
\end{enumerate}
\end{enumerate}
\end{lemma}

\begin{proof}
省略する。ヒント：エタール局所化により概型の場合から従う。Derived
Categories of Schemes, Lemma
\ref{perfect-lemma-locally-closed-where-H0-locally-free} を参照されたい。
\end{proof}

\begin{lemma}
\label{spaces-perfect-lemma-proper-flat-h0}
$S$ を概型とする。$f : X \to Y$ を $S$ 上の代数空間の射とする。次を仮定する。
\begin{enumerate}
\item $f$ は固有、平坦かつ有限表示である。
\item 射 $\Spec(k) \to Y$ を取り、$k$ が体であるとする。このとき
$k = H^0(X_k, \mathcal{O}_{X_k})$ である。% P09-ERR-1041
\end{enumerate}
このとき次が成り立つ。
\begin{enumerate}
\item[(a)] $f_*\mathcal{O}_X = \mathcal{O}_Y$ であり、これは任意の基底変換後も
成り立つ。
\item[(b)] $Y$ 上エタール局所的に
$$
Rf_*\mathcal{O}_X = \mathcal{O}_Y \oplus P
$$
が $D(\mathcal{O}_Y)$ において成り立つ。ここで $P$ は tor 振幅が
$[1, \infty)$ に入る完全対象である。
\end{enumerate}
\end{lemma}

\begin{proof}
(a) と (b) は $Y$ 上エタール局所的に証明すれば十分なので、$Y$ は
アフィン概型であると仮定する。コホモロジーと基底変換（Lemma
\ref{spaces-perfect-lemma-flat-proper-perfect-direct-image-general}）により、複体
$E = Rf_*\mathcal{O}_X$ は完全であり、その形成は任意の基底変換と可換する。
特に $y \in Y$ に対して
$H^0(E \otimes^\mathbf{L} \kappa(y)) =
H^0(X_y, \mathcal{O}_{X_y}) = \kappa(y)$ である。したがって Lemma
\ref{spaces-perfect-lemma-jump-loci} の記法で、$\beta_0(y) \leq 1$ がすべての $y \in Y$ に
対して成り立つ。Lemma
\ref{spaces-perfect-lemma-locally-closed-where-H0-locally-free} を $a = 0$ および $r = 1$ として
適用する。普遍閉部分概型 $j : Z \to Y$ を得る。これは
$H^0(Lj^*E)$ が可逆であること、および同補題の (4)(a), (b), (c) の同値性に
よって特徴づけられる。$E$ の形成は基底変換と可換するので
$$
Lf^*E = R\text{pr}_{1, *}\mathcal{O}_{X \times_Y X}
$$
である。射 $\text{pr}_1 : X \times_Y X$ は切断、すなわち対角射
$\Delta$（$X$ の $Y$ 上の）をもつ。したがって合成が恒等となる写像
$$
\mathcal{O}_X \longrightarrow R\text{pr}_{1, *}\mathcal{O}_{X \times_Y X}
\longrightarrow \mathcal{O}_X
$$
を $D(\mathcal{O}_X)$ において得る。ゆえに
$R\text{pr}_{1, *}\mathcal{O}_{X \times_Y X} = \mathcal{O}_X \oplus E'$
が $D(\mathcal{O}_X)$ において成り立つ。したがって $\mathcal{O}_X$ は
$H^0(Lf^*E)$ の直和因子であり、上で引用した補題の (4)(c) と (4)(a) の
同値性により $X \to Y$ は $Z$ を経由すると結論する。
$\{X \to Y\}$ は fppf 被覆なので $Z = Y$ である。したがって
$f_*\mathcal{O}_X$ は可逆 $\mathcal{O}_Y$-加群である。したがって
$\mathcal{O}_Y \to f_*\mathcal{O}_X$ は同型であると結論する。実際、環準同型
$A \to B$ において $B$ が $A$-加群として可逆ならばその準同型は同型である。
仮定は基底変換で保たれるので、(a) が成り立つ。

\medskip\noindent
(b) の証明。上で、すべての $y \in Y$ に対して写像
$\mathcal{O}_Y \to H^0(E \otimes^\mathbf{L} \kappa(y))$ が全射であることを
示した。したがって More on Algebra, Lemma
\ref{more-algebra-lemma-better-cut-complex-in-two} を適用すると、$y$ の開近傍で
分解 $Rf_*\mathcal{O}_X = \mathcal{O}_Y \oplus P$ を得る。
\end{proof}

\begin{lemma}
\label{spaces-perfect-lemma-proper-flat-geom-red-connected}
$S$ を概型とする。$f : X \to Y$ を $S$ 上の代数空間の射とする。次を仮定する。
\begin{enumerate}
\item $f$ は固有、平坦かつ有限表示である。
\item $f$ の幾何学的ファイバーは被約かつ連結である。
\end{enumerate}
このとき $f_*\mathcal{O}_X = \mathcal{O}_Y$ であり、これは任意の基底変換後も
成り立つ。
\end{lemma}

\begin{proof}
Lemma \ref{spaces-perfect-lemma-proper-flat-h0} により、等式
$k = H^0(X_k, \mathcal{O}_{X_k})$ を、すべての射 $\Spec(k) \to Y$
（$k$ が体）に対して示せば十分で
ある。これは Spaces over Fields, Lemma
\ref{spaces-over-fields-lemma-proper-geometrically-reduced-global-sections} と、
$X_k$ が幾何学的連結かつ幾何学的被約であることから従う。
\end{proof}







\section{その他の応用}
\label{spaces-perfect-section-other-applications}

\noindent
この節では、上で展開した理論から導けるいくつかの結果を述べ、証明する。

\begin{lemma}
\label{spaces-perfect-lemma-countable-cohomology}
$S$ を概型とする。$X$ を $S$ 上の準コンパクトかつ準分離的な代数空間とする。
$K$ を $D_\QCoh(\mathcal{O}_X)$ の対象であって、コホモロジー層 $H^i(K)$ の
$X$ 上エタールなアフィン概型上の切断集合が可算であるものとする。このとき、
任意の準コンパクトかつ準分離的なエタール射 $U \to X$ と、
任意の完全対象 $E$（$D(\mathcal{O}_X)$ の）に対して、集合
$$
H^i(U, K \otimes^\mathbf{L} E),\quad \Ext^i(E|_U, K|_U)
$$
は可算である。
\end{lemma}

\begin{proof}
Cohomology on Sites, Lemma
\ref{sites-cohomology-lemma-dual-perfect-complex} を用いると、群
$H^i(U, K \otimes^\mathbf{L} E)$ について結果を証明すれば十分であると
分かる。Lemma \ref{spaces-perfect-lemma-induction-principle} の帰納原理を用いて補題を
証明する。

\medskip\noindent
$U = \Spec(A)$ がアフィンならば、結果は概型の場合から従う。Derived
Categories of Schemes, Lemma
\ref{perfect-lemma-countable-cohomology} を参照されたい。

\medskip\noindent
証明を完了するには次を示せば十分である。$(U \subset W, V \to W)$ が
基本識別三角形であり、結果が $U$, $V$, および $U \times_W V$ に対して
成り立つならば、結果は $W$ に対しても成り立つ。% P09-ERR-1042
これは Mayer--Vietoris 列の直ちに分かる帰結である。Lemma
\ref{spaces-perfect-lemma-unbounded-mayer-vietoris} を参照されたい。
\end{proof}

\begin{lemma}
\label{spaces-perfect-lemma-countable}
$S$ を概型とする。$X$ を $S$ 上の準コンパクトかつ準分離的な代数空間とする。
$\mathcal{O}_X$ の $X$ 上エタールなアフィン概型上の切断集合が可算であると
仮定する。$K$ を $D_\QCoh(\mathcal{O}_X)$ の対象とする。次は同値である。
\begin{enumerate}
\item $K = \text{hocolim} E_n$ であり、$E_n$ は
$D(\mathcal{O}_X)$ の完全対象である。
\item コホモロジー層 $H^i(K)$ の $X$ 上エタールなアフィン概型上の
切断集合は可算である。
\end{enumerate}
\end{lemma}

\begin{proof}
(1) が成り立つならば (2) が成り立つ。実際、ホモトピー余極限は
コホモロジー層を取ることと可換する（Derived Categories, Lemma
\ref{derived-lemma-cohomology-of-hocolim} による）。また完全複体は局所的に
有限階数自由 $\mathcal{O}_X$-加群からなる有限複体と同型なので、$X$ に
関する仮定により (2) を満たす。% P09-ERR-1043

\medskip\noindent
(2) を仮定する。K-入射複体 $\mathcal{K}^\bullet$ を選び、これで $K$ を表す。
完全生成子 $E$（$D_\QCoh(\mathcal{O}_X)$ の）を選び、これを K-入射複体
$\mathcal{I}^\bullet$ で表す。Theorem \ref{spaces-perfect-theorem-DQCoh-is-Ddga} とその
証明によれば、三角圏の同値
$F : D_\QCoh(\mathcal{O}_X) \to D(A, \text{d})$ が存在する。ここで
$(A, \text{d})$ は微分次数付き代数
$$
(A, \text{d}) =
\Hom_{\text{Comp}^{dg}(\mathcal{O}_X)}
(\mathcal{I}^\bullet, \mathcal{I}^\bullet)
$$
である。この同値は $K$ を微分次数付き加群
$$
M = \Hom_{\text{Comp}^{dg}(\mathcal{O}_X)}
(\mathcal{I}^\bullet, \mathcal{K}^\bullet)
$$
に写す。$H^i(A) = \Ext^i(E, E)$ および
$H^i(M) = \Ext^i(E, K)$ であることに注意する。さらに $F$ は同値なので、
この同値とその擬逆はホモトピー余極限と可換する。したがって $M$ を
$D(A, \text{d})$ のコンパクト対象のホモトピー余極限として書けば十分である。
Differential Graded Algebra, Lemma \ref{dga-lemma-countable} により、
$\Ext^i(E, E)$ と $\Ext^i(E, K)$ が可算であることを各 $i$ について示せば
十分である。
これは Lemma \ref{spaces-perfect-lemma-countable-cohomology} から従う。
\end{proof}

\begin{lemma}
\label{spaces-perfect-lemma-computing-sections-as-colim}
$A$ を環とする。$f : U \to X$ を、$A$ 上の代数空間の平坦な有限表示射とする。
このとき次が成り立つ。
\begin{enumerate}
\item 完全対象 $L_n$（$D(\mathcal{O}_X)$ の）の逆系で
$$
R\Gamma(U, Lf^*K) = \text{hocolim}\ R\Hom_X(L_n, K)
$$
が $D(A)$ において成り立つものが存在し、これは
$K$（$D_\QCoh(\mathcal{O}_X)$ の）について関手的である。
\item 完全対象 $E_n$（$D(\mathcal{O}_X)$ の）の系で
$$
R\Gamma(U, Lf^*K) = \text{hocolim}\ R\Gamma(X, E_n \otimes^\mathbf{L} K)
$$
が $D(A)$ において成り立つものが存在し、これは
$K$（$D_\QCoh(\mathcal{O}_X)$ の）について関手的である。
\end{enumerate}
\end{lemma}

\begin{proof}
Lemma \ref{spaces-perfect-lemma-cohomology-base-change} により
$$
R\Gamma(U, Lf^*K) = R\Gamma(X, Rf_*\mathcal{O}_U \otimes^\mathbf{L} K)
$$
が $K$ について関手的に成り立つ。$R\Gamma(X, -)$ は直和と可換するので
ホモトピー余極限と可換することに注意する。Lemma
\ref{spaces-perfect-lemma-quasi-coherence-pushforward-direct-sums} を参照されたい。同様に、
$- \otimes^\mathbf{L} K$ は導来余極限と可換する。実際、
$- \otimes^\mathbf{L} K$ は直和と可換する
（$D(\mathcal{O}_X)$ における直和は、代表複体の直和によって与えられる）。
したがって (2) を証明するには
$Rf_*\mathcal{O}_U = \text{hocolim} E_n$ と書けば十分である。ここで $E_n$ は
$D(\mathcal{O}_X)$ の完全対象の系である。これができれば
$L_n = E_n^\vee$ とおくことで (1) を得る。Cohomology on Sites, Lemma
\ref{sites-cohomology-lemma-dual-perfect-complex} を参照されたい。

\medskip\noindent
$A = \colim A_i$ と書く。ここで $A_i$ は $\mathbf{Z}$ 上有限型である。
Limits of Spaces, Lemma
\ref{spaces-limits-lemma-descend-finite-presentation} により、ある $i$ と有限表示射
$U_i \to X_i \to \Spec(A_i)$ を取れて、その $\Spec(A)$ への基底変換は
$U \to X \to \Spec(A)$ を復元する。$i$ を大きくした後、
$f_i : U_i \to X_i$ は平坦であると仮定してよい。Limits of Spaces, Lemma
\ref{spaces-limits-lemma-descend-flat} を参照されたい。Lemma
\ref{spaces-perfect-lemma-compare-base-change} により、$Rf_{i, *}\mathcal{O}_{U_i}$ の
$g : X \to X_i$ による導来引き戻しは $Rf_*\mathcal{O}_U$ に等しい。
$Lg^*$ は導来余極限と可換するので、$f_i$ について望む主張を証明すれば
十分である。したがって $U$ と $X$ は $\mathbf{Z}$ 上有限型であると仮定して
よい。

\medskip\noindent
$f : U \to X$ を $\mathbf{Z}$ 上有限型な代数空間の射と仮定する。
証明を完了するため、$Rf_*\mathcal{O}_U$ が完全複体のホモトピー余極限で
あることを示す。そのため Lemma \ref{spaces-perfect-lemma-countable} を適用する。したがって
$R^if_*\mathcal{O}_U$ が $X$ 上エタールなアフィン概型上で可算な切断集合を
もつことを示せば十分である。これは構造層に Lemma
\ref{spaces-perfect-lemma-countable-cohomology} を適用することから従う。
\end{proof}








\section{分解性質}
\label{spaces-perfect-section-resolution-property}

\noindent
この節は代数空間に対する Derived Categories of Schemes, Section
\ref{perfect-section-resolution-property} の類似である。まず同節を
読まれたい。体上の滑らかで固有な代数空間が常に分解性質をもつか、あるいは
そうでないかは現在知られていない。この問いの答えをご存じなら、
\href{mailto:stacks.project@gmail.com}{stacks.project@gmail.com} まで
電子メールを送られたい。

\medskip\noindent
一般の代数空間について考えることにはほとんど意味がないものの、次の定義を
設けることができる。

\begin{definition}
\label{spaces-perfect-definition-resolution-property}
$S$ を概型とする。$X$ を $S$ 上の代数空間とする。$X$ が
{\it 分解性質} をもつとは、有限型のすべての準連接 $\mathcal{O}_X$-加群が
有限局所自由 $\mathcal{O}_X$-加群の商であることをいう。
\end{definition}

\noindent
$X$ が準コンパクトかつ準分離的な代数空間ならば、有限表示のすべての
$\mathcal{O}_X$-加群加群（自動的に準連接）が有限局所自由
$\mathcal{O}_X$-加群の商であることを確かめれば十分である。% P09-ERR-1044
Limits of Spaces, Lemma
\ref{spaces-limits-lemma-finite-directed-colimit-surjective-maps} を参照されたい。
$X$ がネーター代数空間ならば、有限型準連接加群はちょうど連接
$\mathcal{O}_X$-加群である。Cohomology of Spaces, Lemma
\ref{spaces-cohomology-lemma-coherent-Noetherian} を参照されたい。

\begin{lemma}
\label{spaces-perfect-lemma-resolution-property-ample-relative}
$S$ を概型とする。$f : X \to Y$ を $S$ 上の代数空間の射とする。次を仮定する。
\begin{enumerate}
\item $Y$ は準コンパクトかつ準分離的で、分解性質をもつ。
\item $f$-豊富な可逆加群が $X$ 上に存在する（Divisors on Spaces,
Definition \ref{spaces-divisors-definition-relatively-ample}）。
\end{enumerate}
このとき $X$ は分解性質をもつ。
\end{lemma}

\begin{proof}
$\mathcal{F}$ を有限型準連接 $\mathcal{O}_X$-加群とする。
$\mathcal{L}$ を $f$-豊富な可逆加群とする。アフィン概型 $V$ と全射エタール射
$V \to Y$ を選ぶ。$U = V \times_Y X$ とおく。このとき
$\mathcal{L}|_U$ は $U$ 上豊富である。Properties, Proposition
\ref{properties-proposition-characterize-ample} により、有限個の写像
$s_i : \mathcal{L}^{\otimes n_i}|_U \to \mathcal{F}|_U$ であって、一緒にすると
全射であるものが存在する。準連接 $\mathcal{O}_Y$-加群
$$
\mathcal{H}_n =
f_*(\mathcal{F} \otimes_{\mathcal{O}_X} \mathcal{L}^{\otimes n})
$$
を考える。$s_i$ を $V$ 上の層 $\mathcal{H}_{-n_i}$ の切断とみなしてよい。
有限局所自由 $\mathcal{O}_Y$-加群 $\mathcal{E}_i$ と写像
$\mathcal{E}_i \to \mathcal{H}_{-n_i}$ であって、$s_i$ が像に入るものを
取れると仮定しよう。このとき対応する写像
$$
f^*\mathcal{E}_i
\otimes_{\mathcal{O}_X}
\mathcal{L}^{\otimes n_i} \longrightarrow \mathcal{F}
$$
は一緒にすると全射となり、補題が証明される。Limits of Spaces, Lemma
\ref{spaces-limits-lemma-directed-colimit-finite-type} により、各 $i$ に対して
有限型準連接部分加群 $\mathcal{H}'_i \subset  \mathcal{H}_{-n_i}$ で、
切断 $s_i$（$V$ 上の）を含むものを取れる。したがって $Y$ の分解性質により
全射 $\mathcal{E}_i \to \mathcal{H}'_i$ を得て、結論が従う。
\end{proof}

\begin{lemma}
\label{spaces-perfect-lemma-resolution-property-goes-up-affine}
$S$ を概型とする。$f : X \to Y$ を、$S$ 上の代数空間のアフィンまたは
準アフィン射とし、$Y$ は準コンパクトかつ準分離的であるとする。
$Y$ が分解性質をもつならば $X$ も分解性質をもつ。
\end{lemma}

\begin{proof}
Divisors on Spaces, Lemma
\ref{spaces-divisors-lemma-quasi-affine-O-ample} により、これは Lemma
\ref{spaces-perfect-lemma-resolution-property-ample-relative} の特別な場合である。
\end{proof}

\noindent
分解性質が下へ移ることを証明できる場合を挙げよう。

\begin{lemma}
\label{spaces-perfect-lemma-resolution-property-goes-down-finite-flat}
$S$ を概型とする。$f : X \to Y$ を、$S$ 上の代数空間の全射な有限局所自由射と
する。$X$ が分解性質をもつならば $Y$ も分解性質をもつ。
\end{lemma}

\begin{proof}
この条件の意味は、$f$ がアフィンであり、$f_*\mathcal{O}_X$ が正の階数の
有限局所自由 $\mathcal{O}_Y$-加群であるということである。
$\mathcal{G}$ を有限型準連接 $\mathcal{O}_Y$-加群とする。仮定により、ある
全射 $\mathcal{E} \to f^*\mathcal{G}$ が存在する。ここで
$\mathcal{O}_X$-加群 $\mathcal{E}$ は有限局所自由である。$f_*$ は完全なので
（Cohomology of Spaces, Section
\ref{spaces-cohomology-section-finite-morphisms}）、全射
$$
f_*\mathcal{E}
\longrightarrow
f_*f^*\mathcal{G} = \mathcal{G} \otimes_{\mathcal{O}_Y} f_*\mathcal{O}_X
$$
を得る。双対を取ると全射
$$
f_*\mathcal{E}
\otimes_{\mathcal{O}_Y}
\SheafHom_{\mathcal{O}_Y}(f_*\mathcal{O}_X, \mathcal{O}_Y)
\longrightarrow
\mathcal{G}
$$
を得る。$f_*\mathcal{E}$ は有限局所自由なので、結論が従う。
\end{proof}

\noindent
代数空間の分解性質についてさらに詳しくは、More on Morphisms of Spaces,
Section \ref{spaces-more-morphisms-section-resolution-property} を参照されたい。













\section{有界性の検出}
\label{spaces-perfect-section-detecting-boundedness}

\noindent
この節では、Section \ref{spaces-perfect-section-generating} で構成された、準コンパクトかつ
準分離的な概型の $D_\QCoh$ のコンパクト生成子が特別な性質をもつことを
示す。この節と非常によく似ているので、まず同節を読むことを勧める。

\begin{lemma}
\label{spaces-perfect-lemma-bounded-truncation}
$S$ を概型とする。$X$ を $S$ 上の準コンパクトかつ準分離的な代数空間とする。
$P \in D_{perf}(\mathcal{O}_X)$ および $E \in D_{\QCoh}(\mathcal{O}_X)$ とする。
$a \in \mathbf{Z}$ とする。次は同値である。
\begin{enumerate}
\item $\Hom_{D(\mathcal{O}_X)}(P[-i], E) = 0$ が $i \gg 0$ に対して成り立つ。
\item $\Hom_{D(\mathcal{O}_X)}(P[-i], \tau_{\geq a} E) = 0$ が
$i \gg 0$ に対して成り立つ。
\end{enumerate}
\end{lemma}

\begin{proof}
三角形 $ \tau_{< a} E \to E \to \tau_{\geq a} E \to$ を用いると、
$$
\Hom_{D(\mathcal{O}_X)}(P[-i], \tau_{< a} E) =
\Hom_{D(\mathcal{O}_X)}(P, (\tau_{< a} E)[i]) = 0
$$
が $i \gg 0$ に対して成り立つことを示せば同値性が従うと分かる。
$P$ は完全なので、これは Lemma
\ref{spaces-perfect-lemma-ext-from-perfect-into-bounded-QCoh} により成り立つ。
\end{proof}

\begin{lemma}
\label{spaces-perfect-lemma-bounded-below-truncation}
$S$ を概型とする。$X$ を $S$ 上の準コンパクトかつ準分離的な代数空間とする。
$P \in D_{perf}(\mathcal{O}_X)$ および $E \in D_{\QCoh}(\mathcal{O}_X)$ とする。
$a \in \mathbf{Z}$ とする。次は同値である。
\begin{enumerate}
\item $\Hom_{D(\mathcal{O}_X)}(P[-i], E) = 0$ が $i \ll 0$ に対して成り立つ。
\item $\Hom_{D(\mathcal{O}_X)}(P[-i], \tau_{\leq a} E) = 0$ が
$i \ll 0$ に対して成り立つ。
\end{enumerate}
\end{lemma}

\begin{proof}
三角形 $ \tau_{\leq a} E \to E \to \tau_{> a} E \to$ を用いると、
$$
\Hom_{D(\mathcal{O}_X)}(P[-i], \tau_{> a} E) =
\Hom_{D(\mathcal{O}_X)}(P, (\tau_{> a} E)[i]) = 0
$$
が $i \ll 0$ に対して成り立つことを示せば同値性が従うと分かる。
$P$ は完全なので、これは Lemma
\ref{spaces-perfect-lemma-ext-from-perfect-into-bounded-QCoh} により成り立つ。
\end{proof}

\begin{proposition}
\label{spaces-perfect-proposition-detecting-bounded-above}
$S$ を概型とする。$X$ を $S$ 上の準コンパクトかつ準分離的な代数空間とする。
$G \in D_{perf}(\mathcal{O}_X)$ を $D_\QCoh (\mathcal{O}_X)$ を生成する
完全複体とする。$E \in D_\QCoh (\mathcal{O}_X)$ とする。次は同値である。
\begin{enumerate}
\item $E \in D^-_\QCoh (\mathcal{O}_X)$ である。
\item $\Hom_{D(\mathcal{O}_X)}(G[-i], E) = 0$ が $i \gg 0$ に対して成り立つ。
\item $\Ext^i_X(G, E) = 0$ が $i \gg 0$ に対して成り立つ。
\item $R\Hom_X(G, E)$ は $D^-(\mathbf{Z})$ に属する。
\item $H^i(X, G^\vee \otimes_{\mathcal{O}_X}^\mathbf{L} E) = 0$ が
$i \gg 0$ に対して成り立つ。
\item $R\Gamma(X, G^\vee \otimes_{\mathcal{O}_X}^\mathbf{L} E)$ は
$D^-(\mathbf{Z})$ に属する。
\item すべての完全対象 $P$（$D(\mathcal{O}_X)$ の）に対して次が成り立つ。
\begin{enumerate}
\item 主張 (2), (3), (4) は $G$ を $P$ で置き換えても成り立つ。
\item $H^i(X, P \otimes_{\mathcal{O}_X}^\mathbf{L} E) = 0$ が
$i \gg 0$ に対して成り立つ。
\item $R\Gamma(X, P \otimes_{\mathcal{O}_X}^\mathbf{L} E)$ は
$D^-(\mathbf{Z})$ に属する。
\end{enumerate}
\end{enumerate}
\end{proposition}

\begin{proof}
(1) を仮定する。
$\Hom_{D(\mathcal{O}_X)}(G[-i], E) =
\Hom_{D(\mathcal{O}_X)}(G, E[i])$ なので、Lemma
\ref{spaces-perfect-lemma-ext-from-perfect-into-bounded-QCoh} によりこれは $i \gg 0$ に
対して零である。これで (1) が (2) を含意することが示された。

\medskip\noindent
部分 (2), (3), (4) は Cohomology on Sites, Section
\ref{sites-cohomology-section-global-RHom} の議論により同値である。
部分 (5) と (6) は、定義により
$H^i(X, -) = H^i(R\Gamma(X, -))$ なので同値である。同値な条件 (2), (3), (4)
と同値な条件 (5), (6) は、Cohomology on Sites, Lemma
\ref{sites-cohomology-lemma-dual-perfect-complex} および
$(G[-i])^\vee = G^\vee[i]$ という事実により互いに同値である。

\medskip\noindent
(7) が (2) を含意することは明らかである。逆に、同値な条件 (2)--(6) が
(7) を含意することを証明しよう。Remark
\ref{spaces-perfect-remark-classical-generator} により、$G$ は
$D_{perf}(\mathcal{O}_X)$ の古典的生成子であることを思い出す。
$P \in D_{perf}(\mathcal{O}_X)$ に対して、$T(P)$ を
$R\Hom_X(P, E)$ が $D^-(\mathbf{Z})$ に属するという主張とする。
明らかに $T$ は直和へ遺伝し、識別三角形に対する三者中二者性を満たし、
直和因子へ遺伝し、シフトで保たれる。したがって Derived Categories, Remark
\ref{derived-remark-check-on-generator} により、(4) は $T$ が
$D_{perf}(\mathcal{O}_X)$ のすべてで成り立つことを含意する。他のすべての
性質にも同じ議論が使える。ただし性質 (7)(b), (7)(c) については、さらに
$P \mapsto P^\vee$ が $D_{perf}(\mathcal{O}_X)$ の自己同値であることを用いる。
細部は省略する。

\medskip\noindent
同値な条件 (2)--(7) が (1) を含意することを、Lemma
\ref{spaces-perfect-lemma-induction-principle} の帰納原理を用いて証明する。

\medskip\noindent
まず (2)--(7) $\Rightarrow$ (1) を、$X$ がアフィンの場合に証明する。
これは概型の場合から従う。Derived Categories of Schemes, Proposition
\ref{perfect-proposition-detecting-bounded-above} を参照されたい。

\medskip\noindent
次に $(U \subset X, j : V \to X)$ を $S$ 上の準コンパクトかつ準分離的な
代数空間の基本識別正方形とし、(2)--(7) $\Rightarrow$ (1) が
$U$, $V$, および $U \times_X V$ に対して既知であると仮定する。証明を
完了するには (2)--(7) $\Rightarrow$ (1) を $X$ に対して示さなければならない。
$E \in D_\QCoh(\mathcal{O}_X)$ が (2)--(7) を満たすとする。Lemma
\ref{spaces-perfect-lemma-direct-summand-of-a-restriction} および Theorem
\ref{spaces-perfect-theorem-bondal-van-den-Bergh} により、完全複体 $Q$（$X$ 上）であって、
$Q|_U$ が $D_\QCoh (\mathcal{O}_U)$ を生成するものが存在する。

\medskip\noindent
$V = \Spec(A)$ と書く。$Z \subset V$ を $X \setminus U$ の逆像であり、
それへ同型に写る被約閉部分概型とする（Definition
\ref{spaces-perfect-definition-elementary-distinguished-square} を参照）。これは $V$ の
レトロコンパクト閉部分集合である。$f_1, \ldots, f_r \in A$ を
$Z = V(f_1, \ldots, f_r)$ となるように選ぶ。
$K \in D(\mathcal{O}_V)$ を、$f_1, \ldots, f_r$ に関する Koszul
複体（$A$ 上の）に対応する完全対象とする。$K$ は $Z$ 上に台をもつので、直像
$K' = Rj_*K$ は $D(\mathcal{O}_X)$ の完全対象であり、その $V$ への制限は
$K$ であることに注意する（Lemmas \ref{spaces-perfect-lemma-pushforward-perfect} および
\ref{spaces-perfect-lemma-pushforward-with-support-in-open} を参照）。仮定により、
$R\Hom_{\mathcal{O}_X}(Q, E)$ と $R\Hom_{\mathcal{O}_X}(K', E)$ は上に有界で
ある。

\medskip\noindent
Lemma \ref{spaces-perfect-lemma-pushforward-with-support-in-open} により
$K' =  j_!K$ であり、したがって
$$
\Hom_{D(\mathcal{O}_X)}(K'[-i], E) =
\Hom_{D(\mathcal{O}_V)}(K[-i], E|_V) = 0
$$
が $i \gg 0$ に対して成り立つ。ゆえに Derived Categories of Schemes,
Lemma \ref{perfect-lemma-orthogonal-koszul-first-variant} を $E|_V$ に適用して、
整数 $a$ であって
$$
\tau_{\geq a}(E|_V) =
\tau_{\geq a} R (U \times_X V \to V)_* (E|_{U \times_X V})
$$
となるものを得る。このとき
$\tau_{\geq a} E = \tau_{\geq a} R (U \to X)_* (E |_U)$ である
（標準写像を $U$ および $V$ に制限すると同型になることを確かめればよい）。
したがって Lemma \ref{spaces-perfect-lemma-bounded-truncation} を二度用いると
$$
\Hom_{D(\mathcal{O}_U)}(Q|_U [-i], E|_U) =
\Hom_{D(\mathcal{O}_X)}(Q[-i], R (U \to X)_* (E|_U)) = 0
$$
が $i \gg 0$ に対して成り立つと分かる。命題は $U$ と生成子 $Q|_U$ に
対して成り立つので、$E|_U \in D^-_\QCoh(\mathcal{O}_U)$ である。
さらに関手 $R (U \to X)_*$ は $D^-_\QCoh$ を保つので（Lemma
\ref{spaces-perfect-lemma-quasi-coherence-direct-image} による）、
$\tau_{\geq a}E \in D^-_\QCoh(\mathcal{O}_X)$ を得る。したがって
$E \in D^-_\QCoh (\mathcal{O}_X)$ である。
\end{proof}

\begin{proposition}
\label{spaces-perfect-proposition-detecting-bounded-below}
$S$ を概型とする。$X$ を $S$ 上の準コンパクトかつ準分離的な代数空間とする。
$G \in D_{perf}(\mathcal{O}_X)$ を $D_\QCoh (\mathcal{O}_X)$ を生成する
完全複体とする。$E \in D_\QCoh (\mathcal{O}_X)$ とする。次は同値である。
\begin{enumerate}
\item $E \in D^+_\QCoh (\mathcal{O}_X)$ である。
\item $\Hom_{D(\mathcal{O}_X)}(G[-i], E) = 0$ が $i \ll 0$ に対して成り立つ。
\item $\Ext^i_X(G, E) = 0$ が $i \ll 0$ に対して成り立つ。
\item $R\Hom_X(G, E)$ は $D^+(\mathbf{Z})$ に属する。
\item $H^i(X, G^\vee \otimes_{\mathcal{O}_X}^\mathbf{L} E) = 0$ が
$i \ll 0$ に対して成り立つ。
\item $R\Gamma(X, G^\vee \otimes_{\mathcal{O}_X}^\mathbf{L} E)$ は
$D^+(\mathbf{Z})$ に属する。
\item すべての完全対象 $P$（$D(\mathcal{O}_X)$ の）に対して次が成り立つ。
\begin{enumerate}
\item 主張 (2), (3), (4) は $G$ を $P$ で置き換えても成り立つ。
\item $H^i(X, P \otimes_{\mathcal{O}_X}^\mathbf{L} E) = 0$ が
$i \ll 0$ に対して成り立つ。
\item $R\Gamma(X, P \otimes_{\mathcal{O}_X}^\mathbf{L} E)$ は
$D^+(\mathbf{Z})$ に属する。
\end{enumerate}
\end{enumerate}
\end{proposition}

\begin{proof}
(1) を仮定する。
$\Hom_{D(\mathcal{O}_X)}(G[-i], E) =
\Hom_{D(\mathcal{O}_X)}(G, E[i])$ なので、Lemma
\ref{spaces-perfect-lemma-ext-from-perfect-into-bounded-QCoh} によりこれは $i \ll 0$ に
対して零である。これで (1) が (2) を含意することが示された。

\medskip\noindent
部分 (2), (3), (4) は Cohomology on Sites, Section
\ref{sites-cohomology-section-global-RHom} の議論により同値である。
部分 (5) と (6) は、定義により
$H^i(X, -) = H^i(R\Gamma(X, -))$ なので同値である。同値な条件 (2), (3), (4)
と同値な条件 (5), (6) は、Cohomology on Sites, Lemma
\ref{sites-cohomology-lemma-dual-perfect-complex} および
$(G[-i])^\vee = G^\vee[i]$ という事実により互いに同値である。

\medskip\noindent
(7) が (2) を含意することは明らかである。逆に、同値な条件 (2)--(6) が
(7) を含意することを証明しよう。Remark
\ref{spaces-perfect-remark-classical-generator} により、$G$ は
$D_{perf}(\mathcal{O}_X)$ の古典的生成子であることを思い出す。
$P \in D_{perf}(\mathcal{O}_X)$ に対して、$T(P)$ を
$R\Hom_X(P, E)$ が $D^+(\mathbf{Z})$ に属するという主張とする。
明らかに $T$ は直和へ遺伝し、識別三角形に対する三者中二者性を満たし、
直和因子へ遺伝し、シフトで保たれる。したがって Derived Categories, Remark
\ref{derived-remark-check-on-generator} により、(4) は $T$ が
$D_{perf}(\mathcal{O}_X)$ のすべてで成り立つことを含意する。他のすべての
性質にも同じ議論が使える。ただし性質 (7)(b), (7)(c) については、さらに
$P \mapsto P^\vee$ が $D_{perf}(\mathcal{O}_X)$ の自己同値であることを用いる。
細部は省略する。

\medskip\noindent
同値な条件 (2)--(7) が (1) を含意することを、Lemma
\ref{spaces-perfect-lemma-induction-principle} の帰納原理を用いて証明する。

\medskip\noindent
まず (2)--(7) $\Rightarrow$ (1) を、$X$ がアフィンの場合に証明する。
これは概型の場合から従う。Derived Categories of Schemes, Proposition
\ref{perfect-proposition-detecting-bounded-below} を参照されたい。

\medskip\noindent
次に $(U \subset X, j : V \to X)$ を $S$ 上の準コンパクトかつ準分離的な
代数空間の基本識別正方形とし、(2)--(7) $\Rightarrow$ (1) が
$U$, $V$, および $U \times_X V$ に対して既知であると仮定する。証明を
完了するには (2)--(7) $\Rightarrow$ (1) を $X$ に対して示さなければならない。
$E \in D_\QCoh(\mathcal{O}_X)$ が (2)--(7) を満たすとする。Lemma
\ref{spaces-perfect-lemma-direct-summand-of-a-restriction} および Theorem
\ref{spaces-perfect-theorem-bondal-van-den-Bergh} により、完全複体 $Q$（$X$ 上）であって、
$Q|_U$ が $D_\QCoh (\mathcal{O}_U)$ を生成するものが存在する。

\medskip\noindent
$V = \Spec(A)$ と書く。$Z \subset V$ を $X \setminus U$ の逆像であり、
それへ同型に写る被約閉部分概型とする（Definition
\ref{spaces-perfect-definition-elementary-distinguished-square} を参照）。これは $V$ の
レトロコンパクト閉部分集合である。$f_1, \ldots, f_r \in A$ を
$Z = V(f_1, \ldots, f_r)$ となるように選ぶ。
$K \in D(\mathcal{O}_V)$ を、$f_1, \ldots, f_r$ に関する Koszul
複体（$A$ 上の）に対応する完全対象とする。$K$ は $Z$ 上に台をもつので、直像
$K' = Rj_*K$ は $D(\mathcal{O}_X)$ の完全対象であり、その $V$ への制限は
$K$ であることに注意する（Lemmas \ref{spaces-perfect-lemma-pushforward-perfect} および
\ref{spaces-perfect-lemma-pushforward-with-support-in-open} を参照）。仮定により、
$R\Hom_{\mathcal{O}_X}(Q, E)$ と $R\Hom_{\mathcal{O}_X}(K', E)$ は下に有界で
ある。

\medskip\noindent
Lemma \ref{spaces-perfect-lemma-pushforward-with-support-in-open} により
$K' =  j_!K$ であり、したがって
$$
\Hom_{D(\mathcal{O}_X)}(K'[-i], E) =
\Hom_{D(\mathcal{O}_V)}(K[-i], E|_V) = 0
$$
が $i \ll 0$ に対して成り立つ。ゆえに Derived Categories of Schemes,
Lemma \ref{perfect-lemma-orthogonal-koszul-second-variant} を $E|_V$ に適用して、
整数 $a$ であって
$$
\tau_{\leq a}(E|_V) =
\tau_{\leq a} R (U \times_X V \to V)_* (E|_{U \times_X V})
$$
となるものを得る。このとき
$\tau_{\leq a} E = \tau_{\leq a} R (U \to X)_* (E |_U)$ である
（標準写像を $U$ および $V$ に制限すると同型になることを確かめればよい）。
したがって Lemma \ref{spaces-perfect-lemma-bounded-below-truncation} を二度用いると
$$
\Hom_{D(\mathcal{O}_U)}(Q|_U [-i], E|_U) =
\Hom_{D(\mathcal{O}_X)}(Q[-i], R (U \to X)_* (E|_U)) = 0
$$
が $i \ll 0$ に対して成り立つと分かる。命題は $U$ と生成子 $Q|_U$ に
対して成り立つので、$E|_U \in D^+_\QCoh(\mathcal{O}_U)$ である。
さらに関手 $R (U \to X)_*$ は $D^+_\QCoh$ を保つので（Lemma
\ref{spaces-perfect-lemma-quasi-coherence-direct-image} による）、
$\tau_{\leq a}E \in D^+_\QCoh(\mathcal{O}_X)$ を得る。したがって
$E \in D^+_\QCoh (\mathcal{O}_X)$ である。
\end{proof}











\section{導来圏の準連接対象}
\label{spaces-perfect-section-QC}

\noindent
$S$ を概型とする。$X$ を $S$ 上の代数空間とする。
$X_{affine, \etale}$ は、標準エタール被覆が定める位相を備えた
$X_\etale$ のアフィン対象の圏を表すことを思い出そう。Properties of Spaces,
Definition \ref{spaces-properties-definition-affine-etale-site} を参照されたい。
$X_{affine, \etale}$ のトポスは $X$ の小エタール・トポスであることも
思い出そう。Properties of Spaces, Lemma
\ref{spaces-properties-lemma-alternative} を参照されたい。サイト $X_\etale$ は
構造層 $\mathcal{O}_X$ をもち、その $X_{affine, \etale}$ への制限も
$\mathcal{O}_X$ と書く。このとき環付きトポスの同値
$$
(\Sh(X_{affine, \etale}), \mathcal{O}_X)
\longrightarrow
(\Sh(X_\etale), \mathcal{O}_X)
$$
が存在する。Descent on Spaces, Equation
(\ref{spaces-descent-equation-alternative-small-ringed}) および Descent on
Spaces, Section
\ref{spaces-descent-section-alternative-quasi-coherent} の議論を参照されたい。

\medskip\noindent
この節では、$X_{affine}$ を $X_{affine, \etale}$ の基礎圏に無秩序位相、
すなわち層と前層が一致するような位相を入れたものとする。特に構造層
$\mathcal{O}_X$ は $X_{affine}$ 上の層にもなる。Cohomology on Sites, Section
\ref{sites-cohomology-section-compare} と同様に、環付きサイトの射
$$
\epsilon :
(X_{affine, \etale}, \mathcal{O}_X)
\longrightarrow
(X_{affine}, \mathcal{O}_X)
$$
を得る。% P09-ERR-1045
この節では $D_\QCoh(\mathcal{O}_X)$ を、Cohomology on Sites, Section
\ref{sites-cohomology-section-modules-cohomology} で導入された圏
$\mathit{QC}(X_{affine}, \mathcal{O}_X)$ と同一視する。

\begin{lemma}
\label{spaces-perfect-lemma-DQCoh-alternative-small}
上の状況において、次の三角圏の間に標準的な完全同値が存在する。
\begin{enumerate}
\item $D_\QCoh(\mathcal{O}_X)$,
\item $D_\QCoh(X_{affine, \etale}, \mathcal{O}_X)$,
\item $D_\QCoh(X_{affine}, \mathcal{O}_X)$, and
\item $\mathit{QC}(X_{affine}, \mathcal{O}_X)$.
\end{enumerate}
\end{lemma}

\begin{proof}
$U \to V \to X$ をエタール射とし、$U$ と $V$ はアフィンであるとする。
このとき環準同型 $\mathcal{O}_X(V) \to \mathcal{O}_X(U)$ は平坦である。
したがって (3) と (4) の同値は Cohomology on Sites, Lemma
\ref{sites-cohomology-lemma-cartesian-flat-quasi-coherent} の特別な場合である
（その証明は主張の意味も明らかにする）。

\medskip\noindent
補題の前の議論により、環付きトポスの同値
$(\Sh(X_{affine, \etale}), \mathcal{O}_X) \to
(\Sh(X_\etale), \mathcal{O}_X)$ があり、したがって加群のアーベル圏の間に
同値がある。準連接加群の概念は内在的なので（Modules on Sites, Lemma
\ref{sites-modules-lemma-special-locally-free}）、この同値は準連接加群の
部分圏を保つと分かる。したがって (1) と (2) の三角圏の間に標準的な
完全同値を得る。

\medskip\noindent
(2) と (3) の三角圏の間の完全同値を得るため、Cohomology on Sites, Lemma
\ref{sites-cohomology-lemma-compare-topologies-derived-adequate-modules} を、上の射
$\epsilon : (X_{affine, \etale}, \mathcal{O}_X) \to
(X_{affine}, \mathcal{O}_X)$ に適用する。
$\mathcal{B} = \Ob(X_{affine})$ とし、
$\mathcal{A} \subset \textit{PMod}(X_{affine}, \mathcal{O})$ を、前層
$\mathcal{F}$ であって
$\mathcal{F}(V) \otimes_{\mathcal{O}_X(V)} \mathcal{O}_X(U)
\to \mathcal{F}(U)$ が同型となるものの充満部分圏とする。Descent on Spaces,
Lemma \ref{spaces-descent-lemma-quasi-coherent-alternative-small} により、
$\mathcal{A}$ の対象はちょうど、エタール位相で層であり
$(X_{affine, \etale}, \mathcal{O}_X)$ 上の準連接加群であるものに一致する。
他方、Modules on Sites, Lemma
\ref{sites-modules-lemma-chaotic-quasi-coherent} により、$\mathcal{A}$ の対象は
ちょうど $(X_{affine}, \mathcal{O}_X)$ 上、すなわち無秩序位相における
準連接加群である。したがって Cohomology on Sites, Lemma
\ref{sites-cohomology-lemma-compare-topologies-derived-adequate-modules} が適用
できることを示せば、$\epsilon^*$ と $R\epsilon_*$ を用いて (2) と (3) の
圏の間の標準同値を実際に得る。

\medskip\noindent
四つの条件を検証しなければならない。
\begin{enumerate}
\item $\mathcal{A}$ のすべての対象はエタール位相に対する層である。これは
上で見た。
\item $\mathcal{A}$ は
$\textit{Mod}(X_{affine, \etale}, \mathcal{O}_X)$ の弱 Serre 部分圏である。
上で $\mathcal{A} = \QCoh(X_{affine, \etale}, \mathcal{O}_X)$ と分かった。
また、同値
$\textit{Mod}(X_{affine, \etale}, \mathcal{O}) =
\textit{Mod}(X_\etale, \mathcal{O}_X)$ を介して、これらは $X$ 上の準連接加群に
対応することも上で見た。したがって結果は Properties of Spaces, Lemma
\ref{spaces-properties-lemma-properties-quasi-coherent} および Homology, Lemma
\ref{homology-lemma-characterize-weak-serre-subcategory} による。
% P09-ERR-1046
\item $X_{affine}$ のすべての対象は、各要素が $\mathcal{B}$ に属するような
無秩序位相の被覆をもつ。$\mathcal{B}$ はすべての対象を含むので、これは
成り立つ。
\item すべての対象 $U$（$X_{affine}$ の）と、$\mathcal{F}$（$\mathcal{A}$ に
属する）に対して、$H^p_{Zar}(U, \mathcal{F}) = 0$ が $p > 0$ に対して
成り立つ。これはアフィン上の準連接加群のコホモロジーの消滅による。
Cohomology of Spaces, Section
\ref{spaces-cohomology-section-higher-direct-image} の議論、および Cohomology of
Schemes, Lemma
\ref{coherent-lemma-quasi-coherent-affine-cohomology-zero} を参照されたい。
\end{enumerate}
これで証明が完了する。
\end{proof}

\begin{remark}
\label{spaces-perfect-remark-QC-compare}
$S$ を概型とする。$X$ を $S$ 上の代数空間とする。後に
$\mathit{QC}((\textit{Aff}/X), \mathcal{O})$ も
$D_\QCoh(\mathcal{O}_X)$ に標準的に同値であることを示す。Sheaves on Stacks,
Proposition \ref{stacks-sheaves-proposition-QC-compare} を参照されたい。
\end{remark}
